% =====================================================================
%  Harald Høffding, KORT OVERSIGT OVER DEN NYERE FILOSOFIS HISTORIE.
%  Fjerde, paany omarbejdede Udgave.  Kjøbenhavn og Kristiania: Gyldendalske Boghandel, Nordisk Forlag, 1910.
%  Forlagets Bogtrykkeri.  95 pp.
%
%  GENERATED; DO NOT EDIT BY HAND.  Built by ebook-method/kort-oversigt/ (rebuild.sh) from the scan, the
%  collation decisions and the patches there.
%
%  SOURCE: the National Library of Norway's scan (nb.no item 17246e2b574fbd978af6a80c866e4e4b,
%  URN:NBN:no-nb_digibok_2010102906103), ~/bibliotek/Høffding, Harald/1910-kort-oversigt-4udg.pdf (SCANS.tsv).
%  PDF page = printed page + 10.
%
%  EDITION (2026-10-17): the 1910 fourth edition, chosen by Hans; the only edition with an open scan (NB also holds
%  1906, 1919 and 1922, not open; Runeberg's kortfil is the same 1910 NB copy, unproofread OCR).  The SAGA /
%  Lindhardt og Ringhof e-book (ISBN 9788726363074, "Copyright 1898, 2020") is NOT the 1898 text but a LATER edition
%  (it gives Wundt's death, 1920; a second Tillæg „Vurderingsprincipets Historie“; its note cites Den menneskelige
%  Tanke 1910 where 1910's cites Mindre Arbejder; verb plurals modernised): a voting witness only where it agrees.
%
%  HOW ACCURACY WAS ESTABLISHED (2026-10-17; ebook-method/kort-oversigt/, README entry there):
%  (a) Draft: NB's own OCR (ALTO) laid over the page images (nb_scan.py), turned into LaTeX by rule (draft.py), with
%  running heads and sheet signatures dropped; headings from the Indhold; the one note placed at its mark.
%  (b) Word collation (collate/check.py, punctuation included) against an INDEPENDENT tesseract reading of the same
%  images: ~490 disagreements; 216 by the e-book's vote (rules.py), but where it voted for CAPITALS the word is in small
%  capitals in the print (SC, rules2.py: 87), 15 NB å for a, 5 umlauts; 19 by hand where three witnesses agree; 125 at
%  the image.  OPEN 0; italics OPEN 0.
%  (c) Blind crops (b1: two readers): the open checks, the e-book's non-word differences, italics, and an AUDIT of 30
%  rule decisions (30/30 right); controls 24/24.
%  (d) Structure read at the image (s1: 80 crops, two readers; s2: 14 crops, main session): half-title, title page, the
%  blank versos, Indhold, Fortale, the end leaves; every chapter and section head (and the short rule that closes each
%  part); the run-in names in small capitals (29 openings; 13 names mended); the α. β. γ. labels; doubtful paragraph
%  starts (8); the one footnote whole.
%  (e) Sweeps, read at line crops in the main session: quotation marks (the print uses „…“ only; NB dropped „ or read it
%  « or »: 45 places mended, every paragraph now balances), unknown words against the e-book and the other Høffding
%  transcriptions (13 mended: ü/ä/à/è/œ, l' read V), names half in small capitals (6 mended).  The whole text was also
%  compared with the (later-edition) e-book (ebookdiff.py: 1238 differences, mostly edition wording; its 32 non-word
%  suspects went to b1).
%  (f) A blind sampled page check (pagecheck.py): 6 pages from four strata (pp. 5, 12, 43, 44, 58, 82), 1,764 words:
%  0 word errors, 0 punctuation errors, 1 markup error (an italic title, mended); 95 % upper bound 1.7 word errors per
%  1000 words.
%  (g) The translators' 54 odd spots (tr/odd-*.tsv), read at line crops in the main session (o6-o8): 33 faults of the
%  transcription mended: the first line of p. 2, dropped as a running head (draft.py now keeps a line ending in a
%  hyphen); twelve doubled dashes; four specks read as stops or a comma; „done“ for donc; a doubled og; Göttingen,
%  Système (2); a lost paragraph break (p. 32); eight titles only partly italic.  Nine misprints of the print kept and
%  marked PRINTED AS IS.  The rest were questions of grammar, fact or period spelling, settled in translating.
%  NOT done: a whole-book reading at the image.
%
%  EDITORIAL: diplomatic.  Printed as is, and marked: „Er Sjælen er særegen“ (p. 1), „og og“ (p. 3), „Bevægelsee“ (p. 10),
%  „hviken“ (p. 29), „Skrifter“ without its stop (p. 32), „almægtigt Endeligt“ (p. 43), „Væsen Hvad“ (p. 44), „raade
%  Den“ (p. 48), „og og“ (p. 66), „Tænkeré“ (p. 68), „Aarhundres“ (p. 77), „Hypetese“ and „staar, En“ (p. 79),
%  „kvantitive“ and „concience“ (p. 83), „Livsvurdening“ (p. 94).  Kept as printed, unmarked:
%  Burckardt, Melanchton, Réné, Oldenbarnewelts, Darvin (Indhold), Darwin's death 1881, reduit, Holbach in italics.
%  The house dash --- for the print's ---; the publisher's device on the title page is described, not drawn.
% =====================================================================
\documentclass[12pt,a4paper,openany]{book}
\usepackage[utf8]{inputenc}
\usepackage[T1]{fontenc}
\usepackage{libertinus}
\usepackage{textalpha}
\usepackage[danish]{babel}
\usepackage{setspace}
\usepackage[margin=1.2in]{geometry}
\usepackage{fancyhdr}
\setlength{\headheight}{28pt}
\addtolength{\topmargin}{-13.5pt}
\usepackage{hyperref}
\hypersetup{pdftitle={Kort Oversigt over den nyere Filosofis Historie (1910)}, pdfauthor={Harald Høffding}, hidelinks}
\pagestyle{fancy}
\fancyhf{}
\fancyhead[LE,RO]{\thepage}
\fancyhead[RE]{\textit{Harald Høffding: Kort Oversigt}}
\fancyhead[LO]{\textit{\leftmark}}
\fancypagestyle{plain}{\fancyhf{}\fancyhead[LE,RO]{\thepage}}
\setstretch{1.3}
\usepackage{marginnote}
\usepackage{graphicx}
\DeclareUnicodeCharacter{0254}{\reflectbox{c}}  % ɔ: the old Danish id-est mark (a reversed c)
\newcommand{\starmark}[1]{\ifcase#1\or*\or**\or***\fi)}   % the print marks its one note *)
\renewcommand{\thefootnote}{\starmark{\value{footnote}}}   % the print numbers its notes 1) 2) per page, a few *)
\newcommand{\opage}[1]{\marginnote{\footnotesize\textit{[#1]}}}
\newcommand{\apage}[1]{\marginnote{\footnotesize\textit{[#1]}}}
\newcommand{\whiteline}{\par\vspace{\baselineskip}\par}
% the printed Indhold: an essay or chapter line (title, page), a section line
\newcommand{\ientry}[2]{\par\smallskip\noindent #1\dotfill\,#2\par}
\newcommand{\isub}[2]{\par\noindent\hspace*{2em}#1\dotfill\,#2\par}
% the headings (no \addcontentsline: the print has its own Indhold, and a footnote in a heading breaks a toc entry)
\pdfstringdefDisableCommands{\renewcommand{\footnote}[2][]{}}
\newcommand{\prefacehead}[1]{\chapter*{#1}\markboth{Fortale}{Fortale}}
\newcommand{\sectionhead}[1]{\par\medskip\begin{center}\textbf{#1}\end{center}\smallskip}   % the print: centred, bold
\newcommand{\essayhead}[2][]{\clearpage\begin{center}{\Large #2}\end{center}\if\relax\detokenize{#1}\relax\markboth{#2}{#2}\else\markboth{#1}{#1}\fi\medskip}   % [running head] when the title carries a note
\newcommand{\essaysource}[1]{\begin{center}{\small #1}\end{center}\medskip}
\newcommand{\subhead}[1]{\par\medskip\begin{center}#1\end{center}\smallskip}
\newcommand{\introhead}[1]{\begin{center}{\large #1}\\[4pt]\rule{2em}{0.4pt}\end{center}\markboth{Indledning}{Indledning}\medskip}
\newcommand{\closerule}{\par\bigskip\begin{center}\rule{4em}{0.4pt}\end{center}\par}   % the short rule that closes each part of the print
\newcommand{\chapterheadone}[2]{\bigskip\begin{center}{\large #1}\\[4pt]\rule{2em}{0.4pt}\end{center}\markboth{#2}{#2}\medskip}   % VIII: numeral and title on one line
\newcommand{\chapterhead}[2]{\bigskip\begin{center}#1\\[3pt]{\large #2}\\[4pt]\rule{2em}{0.4pt}\end{center}\markboth{#2}{#2}\medskip}   % numeral, title in capitals, short rule
\newcommand{\romanhead}[1]{\par\medskip\begin{center}#1.\end{center}\smallskip}

\begin{document}

% ---------------------------------------------------------------------
%  FRONT MATTER (scan PDF 3-9), read at the image (structure pass s1, 2026-10-17).  PDF 4, 6, 8, 10 are blank;
%  nothing is printed after p. 95 (PDF 106-109: blank leaf, endpapers, board).
% ---------------------------------------------------------------------
\frontmatter
\thispagestyle{empty}
\vspace*{5cm}
\begin{center}
{\large KORT OVERSIGT OVER DEN\\[4pt] NYERE FILOSOFIS HISTORIE}
\end{center}
\clearpage
\begin{titlepage}
\centering
\vspace*{2cm}
{HARALD HØFFDING}\par
\vspace{6pt}\rule{2em}{0.4pt}\par
\vspace{1.2cm}
{\Huge KORT OVERSIGT OVER\\[6pt] DEN NYERE FILOSOFIS\\[6pt] HISTORIE}\par
\vspace{1cm}
{\small FJERDE, PAANY OMARBEJDEDE UDGAVE}\par
\vfill
{\small [the publisher's device: an oval wreath with the monogram GB]}\par
\vfill
{KJØBENHAVN OG KRISTIANIA}\par
{GYLDENDALSKE BOGHANDEL}\par
{NORDISK FORLAG}\par
\vspace{4pt}\rule{2em}{0.4pt}\par\vspace{4pt}
{\footnotesize FORLAGETS BOGTRYKKERI}\par
{\footnotesize 1910}\par
\end{titlepage}
\clearpage

% the Indhold (PDF 7): from NB's ALTO, read at the image (s1-03-01: every entry and page number as printed, Réné and Darvin included)
\chapter*{INDHOLD}
\markboth{Indhold}{Indhold}
\noindent\hfill{\small SIDE}\par
\ientry{Indledning. De filosofiske Problemer}{1}
\ientry{I. Den nyere Filosofis Tilbliven}{5}
\isub{1. Opdagelsen af det naturlige Menneske}{5}
\isub{2. Det nye Verdensbillede}{11}
\isub{3. Den nye Videnskab}{14}
\ientry{II. Systemdannelser}{17}
\isub{a. Réné Descartes}{18}
\isub{b. Thomas Hobbes}{21}
\isub{c. Baruch Spinoza}{26}
\isub{d. Gottfried Wilhelm Leibniz}{32}
\ientry{III. Erfaringsfilosofien}{35}
\isub{a. John Locke}{36}
\isub{b. George Berkeley}{38}
\isub{c. David Hume}{39}
\ientry{IV. Oplysningsfilosofien}{43}
\isub{1. Voltaire, Holbach og Rousseau}{44}
\isub{2. Wolff og Lessing}{48}
\ientry{V. Immanuel Kant og den kritiske Filosofi}{50}
\isub{Biografi og Karakteristik}{50}
\isub{1. Erkendelsesteori}{53}
\isub{2. Etik}{59}
\isub{3. Enhedslære paa æstetisk og biologisk Grundlag}{61}
\ientry{VI. Romantikens Filosofi}{63}
\isub{a. Fichte, Schelling og Hegel}{63}
\isub{b. Schopenhauer, Kierkegaard og Feuerbach}{67}
\ientry{VII. Positivismen}{74}
\isub{1. Saint Simon og Comte}{75}
\isub{2. Stuart Mill}{77}
\isub{3. Darvin og Spencer}{78}
\ientry{VIII. Nyere Problembehandlinger}{81}
\isub{a. Det psykologiske Problem}{82}
\isub{b. Tilværelsesproblemet}{85}
\isub{c. Erkendelsesproblemet}{88}
\isub{d. Vurderingsproblemet}{92}
\clearpage

% the Fortale til fjerde Udgave (PDF 9, unpaged): from NB's ALTO, read at the image (s1-04-01: one paragraph; the date italic, the signature bold)
\prefacehead{FORTALE TIL FJERDE UDGAVE}

I denne nye Udgave har jeg helt omarbejdet den lille Bog efter
de Erfaringer, jeg har gjort ved i en Række af Aar at benytte den
ved det propædeutiske Kursus. Jeg har udeladt Meget, som der
sjelden blev Brug for ved dette Kursus, men har til Gengæld optaget
biografiske Oplysninger, og jeg har desuden søgt at gøre
Fremstillingen mere afrundet og forstaaelig. Dette gælder ikke mindst
det sidste Afsnit, som skulde give en Forestilling om Filosofiens
nuværende Arbejde. --- Jeg bringer min bedste Tak til Hr. Dr.
Anton Thomsen for en Række værdifulde Bemærkninger, der ere
komne denne nye Udgave tilgode.

\textit{Den 28. Marts 1910.}

\hfill\textbf{Harald Høffding.}
\clearpage

\mainmatter
% --- p. 1 ---
\setcounter{footnote}{0}%
\opage{1}\introhead{INDLEDNING.}

\sectionhead{Filosofiske Problemer.}

Filosofiens Historie fremstiller de Forsøg, Tænkere have
gjort paa at drøfte Erkendelsens og Tilværelsens sidste Problemer.
I den nyere Filosofi, d. v. s. de sidste halvfjerde Aarhundreders
Filosofi, ere især fire Hovedproblemer blevne behandlede,
saaledes som den følgende Fremstilling vil vise.

1. \textit{Det psykologiske Problem} opstaar ved Spørgsmaalet om
% PRINTED AS IS (read at the image, odd spots o6)
Sjælelivets Hovedegenskaber. Er Sjælen er særegen Substans,
ellers bestaar dens Væsen i en ejendommelig Virken? Er den
sammensat af mangfoldige selvstændige Elementer, eller har
den Karakteren af Enhed og Helhed? Hvilken af Sjælelivets
forskellige Sider eller Beskaffenheder er den fundamentale?
Hvilke Love forbinde de forskellige sjælelige Elementer med
hverandre indbyrdes?

Medens den psykologiske Betragtning har det faktisk givne
Sjæleliv til Genstand, gives der to Problemer, som ere betingede
ved Modsætningen mellem det faktiske og det gyldige
eller værdifulde Sjæleliv, --- mellem Sjælelivets virkelige Beskaffenhed
og de ideale Fordringer, der kunne stilles til det:
Erkendelsesproblemet og Vurderingsproblemet. Ved disse to
Problemer drejer det sig om Paavisning og Anvendelse af
bestemte Maalestokke for Gyldighed eller Værdi.

2. \textit{Erkendelsesproblemet} (det logiske Problem) opstaar,
naar der spørges om Betingelserne og Grænserne for vor
% --- p. 2 ---
\setcounter{footnote}{0}%
\opage{2}Erkendelses Gyldighed. Erkendelse er menneskelig Aandsvirksomhed,
et Arbejde, vi foretage med vore aandelige Kræfter
og efter Lovene for disse Kræfters Virksomhed; hvorvidt ere
vi da berettigede til at tillægge dens Resultater Gyldighed ud
over vor egen Aand? Hvori bestaar Erkendelsens Sandhed? ---
\textit{Den formelle Logik} viser, hvorledes vi maa tænke for at være
i Overensstemmelse med os selv, være konsekvente. Men det
egentlige Erkendelsesproblem opstaar først, naar vi spørge,
hvilken Betydning den Orden og Sammenhæng, vi tilvejebringe
mellem vore Tanker, har for de Emner, vi ikke selv
frembringe, --- for Tilværelsen, der er mere omfattende end
vor Tænken. \textit{Erkendelsesteorien} (Logik i videre Betydning)
bliver derfor ikke staaende ved den formelle Konsekvens eller
Systematik, men opkaster Spørgsmaalet om selve denne Konsekvens'
reale Betydning.

3. Et analogt Spørgsmaal opkaster \textit{Vurderingsproblemet} (det
etisk-religiøse Problem) med Hensyn til vore Følelses- og
Villiesytringer. Vi bedømme eller vurdere menneskelige Handlinger
(vore egne og andres) og menneskelige Samfundsforhold
og kalde dem fuldkomne eller ufuldkomne, gode eller
onde, ædle eller lave. Hvilken Maalestok have vi for saadan
Vurdering? Er det muligt at begrunde en fast og almengyldig
Maalestok og at gennemføre den med Konsekvens? Det er
\textit{Etikens} Opgave at drøfte disse Spørgsmaal. Men vi vurdere
ogsaa Tilværelsen idethele efter det Forhold, i hvilket den
viser sig at staa til vor etiske Maalestok, til vore Idealer, vore
højeste Værdier. Viser Tingenes Gang i Naturen og i Historien
sig gunstig eller ugunstig, fremmende eller hæmmende
for vore etiske Idealer? Paa Besvarelsen af dette Spørgsmaal
beror den religiøse Opfattelse (i Ordets videste Betydning).
Det er \textit{Religionsfilosofiens} Opgave at drøfte de etiske Idealers
Forhold til Erfaringen, ligesom det er Erkendelsesteoriens
Opgave at undersøge vore logiske Idealers Forhold til Erfaringen.

4. Endeligt kan der spørges om Beskaffenheden af den Tilværelse,
af hvilken det tænkende, følende og villende Væsen
kun er et enkelt Led. Derved stilles \textit{Tilværelsesproblemet} (det
% --- p. 3 ---
\setcounter{footnote}{0}%
\opage{3}kosmologiske eller metafysiske Problem). Det gælder her Muligheden
af en videnskabelig Verdensanskuelse. Formaa vi
at sammenfatte alle Erfaringer til en Helhed, og hvilken Gyldighed
kan et Forsøg i denne Retning have? Ethvert saadant
Forsøg vil faa sit Særpræg ved det Erfaringsomraade eller den
Side af Erfaringen, det særligt betoner, og i hvis Lys det ser
de andre Omraader eller Sider. Man kan saaledes betragte
det Aandelige som det Væsentlige, som Urfænomenet, og
kommer da til en idealistisk Verdensanskuelse, eller gøre det
i Rummet Bevægelige til Urfænomen, og kommer da til en
% PRINTED AS IS (read at the image, odd spots o6)
materialistisk Verdensanskuelse. Man kan betone Enheden og
og Sammenhængen i Verden, og kommer da til en monistisk
Verdensanskuelse, eller Mangfoldigheden og Forskelligheden,
og kommer da til en pluralistisk Verdensanskuelse. Indenfor
enhver af disse Anskuelser kan der igen spørges, om Tilværelsen
er i Udvikling, eller om den i sin inderste Grund er
uforanderlig, og om der virkeliggøres Værdier i Tilværelsen,
eller om dennes Væsen er udtømt i et System af Aarsager og
Virkninger. Dette sidste Spørgsmaal fører til at forbinde Tilværelsesproblemet
og Vurderingsproblemet.

Medens Historien viser, at der har været en Tendens til,
formedelst sangvinsk Tillid til Tankens Magt, at gaa umiddelbart
til Tilværelsesproblemet, fremtræder der senere en afgjort
Bestræbelse for at behandle Psykologi, Erkendelsesteori og
Etik uafhængigt af „Metafysiken“ eller den systematiske Behandling
af Tilværelsesproblemet. Men Interessen for Filosofi
kan udspringe af fire forskellige Motiver, svarende til de fire
Hovedproblemer, og det ene Problem fører naturligt over til
det andet, forsaavidt der er en Analogi imellem dem. Et og
samme Grundproblem drager sig gennem dem alle: Forholdet
mellem Enhed og Mangfoldighed, Kontinuitet og Diskontinuitet
kommer under forskellige Former frem, hvad enten vi undersøge
Personligheden, Videnskaben, Livsførelsen eller Verdensanskuelsen.
De aandelige Kampe paa disse forskellige Omraader
frembyde analoge Træk: det gælder Muligheden af at
sammenfatte til en Enhed, hvad der foreløbigt fremtræder som
en Mangfoldighed, der gør større eller mindre Modstand mod
% --- p. 4 ---
\setcounter{footnote}{0}%
\opage{4}at opgive den sporadiske Karakter. Alle Hovedproblemer bestaa
i Opgaver for den psykiske Energi. Hvis en afsluttende
Løsning af dem er uopnaaelig, betyder det kun, at hvor langt
vi end naa i vor Forsken, danner der sig stadig nye Horizonter
for os, med nye Maal og nye Opgaver\footnote[1]{Smlg. Afhandlingen \textit{Filosofien og Livet} i mine „Mindre Arbejder“. Anden Række.}.

% --- p. 5 ---
\setcounter{footnote}{0}%

\closerule

\opage{5}\chapterhead{I.}{DEN NYERE FILOSOFIS TILBLIVEN.}

\sectionhead{1. Opdagelsen af det naturlige Menneske.}

a. \textsc{Burckardt} har i sit Værk \textit{Die Kultur der Renaissance
in Italien} karakteriseret den italienske Renaissance som en
Opdagelse af Mennesket. De historiske Forhold førte til Frigørelse
af Individerne. Mennesket blev ikke mere blot bedømt
efter sit Forhold til Kirken eller til sin Stand (Korporation),
men blev selvstændig Genstand for Interesse og Undersøgelse.
Herved medvirkede Opdagelsen af den antike Kunst og Literatur.
Man lærte, at der ogsaa udenfor Kirken gaves et ejendommeligt
Menneskeliv med egne Love og egne Formaal.
Synskredsen blev udvidet, og en sammenlignende Betragtning
blev mulig.

I Italien selv kan \textsc{Niccolo da Machiavelli} (1469---1527)
betragtes som karakteristisk Repræsentant for disse Bestræbelser.
Han tilhørte en gammel florentinsk Familie og traadte
i sin Fædrenestads Tjeneste som Diplomat, hvorved han fik
god Lejlighed til at iagttage Mennesker og Begivenheder. Efter
Republikens Fald (1512) traadte han i Forbindelse med Mediceerne,
men vakte derved sine Medborgeres Uvillie, og da
den republikanske Forfatning blev indført igen, tog man ikke
imod hans Tjeneste. Han døde i Mismod 1527. --- Det var
den politiske Interesse, der gjorde ham til Tænker. Italiens
ulykkelige Forhold vakte Ønsket hos ham om at se Oldtidens
Kraft og Tænkemaade vakt igen. Hvorfor skulde man kun
efterligne den antike Kunst og ikke de antike Bedrifter? Men
% --- p. 6 ---
\setcounter{footnote}{0}%
\opage{6}store Bedrifter ere kun mulige, naar man hensynsløst søger at
naa det Maal, man har sat sig! --- Det kan hos Machiavelli
undertiden se ud, somom det er ligegyldigt, hvilket Maal
man sætter sig, naar man kun forfølger det med hensynsløs
Energi. Men der var dog ét Maal, der stedse stod i Baggrunden
for hans Tænken: Italiens Enhed og Storhed. I dette
Maals Tjeneste vare alle Midler ham hellige. I sin berømte
Bog Il \textit{Principe} (1513) har han udtalt sig i denne Retning.
Naar han nu ikke fandt nogen Sans for det Store eller nogen
ret Energi hos sin Tids Italienere, udledte han det af, at Kirken
og Kristendommen havde svækket Lidenskaberne. I sine
\textit{Discorsi} (Afhandlinger over Livius' første ti Bøger) sammenligner
han Oldtidens og den nyere Tids Tænkemaade, og
Sammenligningen falder ikke ud til Fordel for denne sidste.
Man vurderer ikke Ære, Aandskraft og legemlig Styrke tilstrækkeligt.
I sin Beundring for disse Egenskaber er Machiavelli
en ægte Søn af Renaissancen.

I Frankrig fremtræder Renaissancens Aand hos \textsc{Michel} \textsc{de}
\textsc{Montaigne} (1533---1592), en Adelsmand, der levede et rigt
literært Liv paa sit Slot i Nærheden af Bordeaux, fjernt fra
Tidens politiske og religiøse Kampe. Hans Interesse var Studiet
af Menneskelivet, som han kendte det fra Rejser, Bøger
og Selviagttagelse. Hans \textit{Essais} (1580---1588) begynder med
den Udtalelse: Je suis moy mesme le sujet de mon livre.
Vor store Moder Naturen aabenbarer sig i Enhver paa ejendommelig
Maade. Jo mere vi undersøge, des flere Forskelligheder
opdage vi, uden at vi formaa at bringe dem under almindelige
Love. Derfor taler Montaigne mod al Overtro paa
Fornuften (Rationalisme) saavel som mod Dogmetro. Vi kende
kun Tingene saaledes, som vore Sanser vise os dem, og i alle
Forestillinger om Gud iklæde vi ham menneskelige Egenskaber,
ligesom Dyrene, hvis de kunde tænke, vilde forestille sig deres
Guder i Dyreskikkelse.

I rent teoretisk Form fremtræder Interessen for det Menneskelige
hos den lærde Spanier Luis Vives (f. i Valencia
1492, d. i Brügge 1540), der ved sin Bog \textit{De anima et vita}
(1538) er Forgænger for den nyere Erfaringspsykologi. Han
% --- p. 7 ---
\setcounter{footnote}{0}%
\opage{7}fordrer, at Erfaringen skal være Grundlag for al Erkendelse,
og i Kraft af dette Princip udtaler han, at vor Hovedinteresse
ikke er at vide, hvad Sjælen er, men \textit{hvorledes} den virker.
Hans Metode er beskrivende, ikke analyserende og forklarende;
men hans Skildring af de enkelte sjælelige Fænomener,
især af Sindsbevægelserne, har endnu Interesse. Sjælen og
Livskraften ere for ham Et, og han søger overalt at forbinde
Fysiologi og Psykologi.

b. Indenfor Protestantismen virkede \textsc{Melanchton} (1495
--- 1560) for Filosofien gennem en Række af Lærebøger. Men
den største Betydning fik Kalvinismen, indenfor hvilken Kampen
for religiøs Frihed hidførte Bestræbelser for en frisindet
Rets- og Statslære. Saaledes udviklede Borgermesteren i Emden
\textsc{Johannes} \textsc{Althaus} (Althusius) (1557---1638) Folkesuverænitetens
Ide i sin Bog \textit{Politica methodice digesta} (1603). Men
mest berømt er paa dette Omraade Hollænderen \textsc{Hugo} \textsc{de}
\textsc{Groot} (Grotius) (1583---1645). Tidligt vakte han Opmærksomhed
ved sin juridiske og teologiske Lærdom. Som Politiker
hørte han til Oldenbarnewelts aristokratiske og liberalteologiske
Parti. Efter Oldenbarnewelts Død blev han fængslet,
men undkom ved sin Hustrus List og levede derefter i Paris,
tilsidst som svensk Gesandt. I sit Værk \textit{De jure belli et pacis}
(1625) gaar han ud fra, at der er Krig i Verden --- 1) mellem
Stater, 2) mellem en Enkelt og Staten, 3) mellem Enkelte
indbyrdes, og 4) mellem Staten og en Enkelt. 1) Menneskerettigheder
og Humanitetspligter ophæves ikke i Krigen, der
skal føres paa menneskelig Maade og for Fredens Skyld. Ved
dette Princip er Grotius bleven Grundlægger af den moderne
Folkeret. 2) Den Enkeltes Kamp mod Staten er Oprør, som
Grotius forkaster (medens Althaus hævdede, at det kunde
være berettiget). 3) Krigen mellem Enkelte bør begrænses til
Nødværge. 4) Statens Krig mod den Enkelte fremtræder som
Straf. Denne er kun berettiget, naar den Lidelse, der paaføres
Lovovertræderen, opvejes ved et større Gode for ham selv
og for Samfundet. --- I alle disse Forhold gælder Retten uafhængigt
af teologiske Grunde, udspringer af Principer, der
ligge i Menneskets Natur (ex principiis homini internis), nem%
% --- p. 8 ---
\setcounter{footnote}{0}%
\opage{8}lig som Betingelse for det Samfundsliv, til hvilket Mennesket
har en uvilkaarlig Trang. Den vigtigste Betingelse er her, at
ethvert Løfte holdes; dertil have Menneskene forpligtet sig
ved en udtrykkelig eller stiltiende Kontrakt, saasnart de gaa
ind paa at leve i et Samfund.

Omtrent samtidigt med, at Grotius gav Natur- og Folkeretten
et rent menneskeligt Grundlag, hævdede den engelske
Diplomat \textsc{Herbert} \textsc{af} \textsc{Cherbury} (1582---1648) en naturlig,
af Menneskets eget Væsen udspringende Religion \textit{(De veritate}
1624). Han bekæmper saavel dem, der sætte en Kirketro
over Fornufterkendelsen, som dem, der kun have Tillid til
den undersøgende Fornuft, og dem, der ville udlede Alt af
Sanserne. Der gives, hævder han, en umiddelbar, instinktiv
Sans, som fører Menneskene til Anerkendelse af visse Sandheder
(notitiæ communes), en Sans, der er en Virkning af det
af Gud indplantede Selvopholdelsesinstinkt. Saadanne Sandheder
ere f. Ex., at to Sætninger, der stride mod hinanden,
ikke begge kunne være sande, --- at der er en første Aarsag
til Alt, --- at vi skulle gøre mod Andre, hvad man vil, at de
skulle gøre mod os. Den naturlige Religion beror ogsaa paa
et instinktivt Grundlag, en indre Aabenbaring i hver Menneskesjæl.
Et Vidnesbyrd herom finder han i, at vi have Evner
og føle en Trang, der ikke kunne fyldestgøres ved nogen
endelig Genstand. Den naturlige Religions Indhold er Troen
paa et højeste Væsen, der skal dyrkes ved Dyd og fromt
Sindelag, og paa Løn og Straf efter Døden. Spørgsmaal, som
gaa ud over dette, behøve ikke at besvares.

Paa Erkendelseslærens Omraade søgte samtidigt en anden
engelsk Adelsmand at lægge en ny Grund. \textsc{Francis} \textsc{Bacon}
af Verulam (1561---1626) ventede et vældigt Opsving i Kultur,
naar Menneskene ved Hjælp af Naturvidenskaben, som
Oldtid og Middelalder havde forsømt, gjorde sig til Herrer
over Naturkræfterne. Videnskabens Formaal er Berigelse af
Menneskelivet ved nye Opfindelser. Dog fremhæver Bacon
ogsaa selve Glæden ved Betragtningen af Naturen (contemplatio
rerum): Synet af Lyset er herligere end Lysets mangfoldige
Nytte! --- For at kunne udføre sine store videnskabe%
% --- p. 9 ---
\setcounter{footnote}{0}%
\opage{9}lige Planer ansaa han sig for berettiget til at bruge alle Midler.
Det gjaldt om Magt og Rigdom til de kostbare Undersøgelser.
Derfor styrtede han sig ind i det politiske Liv og
naaede --- ved Eftergiven mod Elisabeths og Jakob den Førstes
despotiske Luner --- til de højeste Stillinger. Som Lordkansler
blev han anklaget af Parlamentet for Bestikkelse og
anden Ulovlighed og afsat. Sine sidste Aar tilbragte han i
Stilhed med sine Arbejder.

Bacon kalder sig selv en Herold, der forkynder en ny
Æra, men ikke selv kæmper med. Han fordrer, at man skal
opgive ufrugtbare Spekulationer og anvende Erfaringsmetoden
(Induktionen) paa alle Omraader, de aandelige saavel som de
materielle. Mod den Induktion, man tidligere havde anvendt,
indvender Bacon, at man kun har holdt sig til de positive
Tilfælde, og fordrer, at man ogsaa skal undersøge de Tilfælde,
i hvilke et Fænomen ikke indtræder. Fremdeles fordrer han
en Undersøgelse af den Maade, paa hvilken et Fænomen
varierer under forskellige Forhold, og han lægger Vægt paa,
at man --- for ikke at overvældes af sit Stof --- skal prøve
at opstille Hypoteser og se, om de bekræftes ved nye Erfaringer.
Dog naar han --- fordi han undervurderer den kvantitative
Metode --- ikke til den nye Naturvidenskabs sande
Metode, som vi ville finde hos hans Samtidige Kepler og
Galilei. --- Det er i sit Værk \textit{Novum Organum} (1620), han
beskriver den induktive Metode, efter at han først har undersøgt
Grundene til Videnskabernes Ufuldkommenhed.

For ret at opfatte Naturen maa Mennesket gøre sig til en
ubeskreven Tavle; kun som Barn kan man komme ind i
Naturens Rige. Men der gives naturlige og tilvante Illusioner
(idola mentis), som hindre os deri. De kunne henføres til
fire Arter. Den første Art er fælles for hele Slægten, fordi
den ligger i den menneskelige Natur (idola tribus). Vi ere tilbøjelige til at opfatte Tingene i Lighed med os selv (ex
analogia hominis) istedetfor efter deres Plads i Verdenssammenhængen
(ex analogia universi), og vi ere tilbøjelige til at
finde Orden, Simpelhed og Regelmæssighed i Naturen. Den
anden Art Illusioner beror paa den Enkeltes Individualitet
% --- p. 10 ---
\setcounter{footnote}{0}%
\opage{10}(idola specus: Enhver ser ud paa Naturen fra sin egen Hule).
Nogle fremhæve mest Lighederne, Andre mest Forskellene;
Nogle fremhæve Helhederne, Andre Elementerne, hvoraf de
bestaa. Den tredie Art skyldes Sprogets Indflydelse paa Tænkningen
(idola fori). Ordene dannes efter det praktiske Livs
Behov og angive ofte Forskelle eller Sammenstillinger, Videnskaben
ikke kan godkende. Den fjerde Art (idola theatri)
skyldes overleverede Teoriers Indflydelse. --- Bacon undersøger
ikke, hvorledes det er muligt at aflægge disse Illusioner.
Især ligger der i idola tribus et stort Problem, der i en senere
Tid blev draget frem: vi maa stedse betragte Tilværelsen fra
et menneskeligt Synspunkt (ex analogia hominis), men hvorledes
kan vor Erkendelse saa have absolut Gyldighed?

\sectionhead{2. Det nye Verdensbillede.}

I sin Naturopfattelse, ligesom i sin Opfattelse af Mennesket,
havde Middelalderen, forsaavidt den ikke hentede sine Ideer fra
Bibelen og den kristelige Tradition, bygget paa den græske Oldtid.
Sin Medicin lærte den hos Galenos, sin Astronomi hos Ptolemaios,
sin Filosofi hos Aristoteles. Dens Verdensbillede skyldtes
Forbindelsen af den ptolemæiske og aristoteliske Lære med
den bibelske Opfattelse. Jorden stod fast i Verdens Midtpunkt.
Omkring den kredse, befæstede paa solide, men gennemsigtige
Kugleflader (Sfærer), Maanen, Solen, Planeterne og Fixstjernerne
(disse sidste alle paa en og samme Sfære). Jorden og
Rummet fra Jorden til Maanen, den sublunariske Verden, er
Foranderlighedens og Forgængelighedens Hjem; her ere de
fire Elementer (Jord, Vand, Luft og Ild) i stadig Bevægelse,
idet de søge deres „naturlige Steder“, det Tunge „nedad“, det
Lette „opad“. Ovenover Maanesfæren findes Æteren (det femte
Element), et Stof, der ikke har noget „naturligt Sted“ og derfor
% PRINTED AS IS (read at the image, structure pass s1/s2)
fortsætter sin Bevægelsee evigt og med absolut Regelmæssighed.
Himmellegemernes Bevægelser ere ved denne
Uforanderlighed et Billede af Guddommens Væsen. De bevæge
sig i Cirkler; Cirkelen er den fuldkomneste Figur, fordi
den stedse vender tilbage i sig selv. Verden er begrænset ved
% --- p. 11 ---
\setcounter{footnote}{0}%
\opage{11}Fixstjernesfæren, der umiddelbart bevæges af Guddommen,
medens de lavere Sfærer bevæges af underordnede Stjerneaander.
--- Dette Verdensbillede stemte ikke blot med, hvad
Tidens Autoriteter, Bibelen og Aristoteles lærte, men ogsaa
med, hvad Sanserne umiddelbart syntes at lære. Det kostede
derfor en haard Kamp at fortrænge det. Og denne Kamp
medførte en Omforming af hele Verdensanskuelsen.

Grundlæggeren af det nye Verdensbillede, Nicolaus Copernicus
(Coppernick) blev født i Thorn (1473), studerede i
Krakau og ved forskellige italienske Universiteter og levede
som Domherre i Frauenburg, beskæftiget dels med Studier,
dels med Administration. I Tidens store Kamp tog han ingen
Del, skønt han synes at have havt en vis Sympati med den
religiøse Bevægelse saavel som med Humanismen. Allerede
1506 begyndte han Udarbejdelsen af sin nye astronomiske
Teori, men betænkte sig længe paa at udgive sit Værk \textit{(De
revolutionibus orbium coelestium),} og først paa sin Dødsseng
(1543) fik han et trykt Exemplar i Hænde. --- Her interessere
os især de erkendelsesteoretiske Forudsætninger, paa hvilke
hans Værk hviler.

Naturen virker stedse paa den simpleste Maade. Herimod
vilde det stride, om hele Verden bevægede sig om et saa lille
Legeme som Jorden. Og dertil kom, at man for at forklare,
at Planeterne ikke, som Fixstjernerne, bevægede sig stedse i
samme Retning, men tilsyneladende frem og tilbage, havde
maattet antage, at de ikke bevægede sig i simple Kredse, men
i smaa Kredse, der (som Epicykler) bevægede sig paa de
større Kredse. Et langt simplere Verdensbillede faar man, naar
Solen tænkes som Verdens Midtpunkt, og Jorden og Planeterne
bevæge sig om den.

Foruden Simpelhedens Princip spiller Tanken om Bevægelsens
Relativitet en væsentlig Rolle hos Copernicus. En iagttagen
Bevægelse kan ikke blot forklares ved, at den iagttagne
Genstand bevæger sig, men ogsaa ved, at det iagttagende
Subjekt selv bevæger sig. Antages det nu, at Jorden, ud fra
hvilken Himmellegemernes Bevægelser iagttages, selv er i
Bevægelse (om sin Axe og omkring Solen), ville Iagttagelserne
% --- p. 12 ---
\setcounter{footnote}{0}%
\opage{12}kunne forklares ligesaa godt, men langt simplere, end efter
det gamle System.

Copernicus antog endnu Verden begrænset ved Fixstjernesfæren,
som han mente var Verdens ubevægelige Grænse, og
han troede paa solide Sfærer, i hvilke Stjernerne vare indfattede.
Men han havde dog foretaget en afgørende Ændring og
indskærpet, at vi ikke kunne slutte sikkert fra vor Opfattelse
af Tilværelsen til selve Tilværelsens Beskaffenhed.

Længe blev Copernicus' Lære betragtet som et ørkesløst
Tankeexperiment, og der hengik næsten et halvt Aarhundrede,
inden den vandt Tilslutning hos selvstændige Forskere.

\textsc{Giordano} \textsc{Bruno} staar her som den, der ikke blot drog
videre gaaende Konsekvenser af den nye Lære, men lagde
den til Grund for en hel Verdensanskuelse, i de store Træk
den nyere Tids ejendommelige Grundopfattelse og dens
første Forsøg paa en selvstændig Filosofi. Bruno er født i
Nola i Syditalien (1548) og traadte som ganske ung ind i
Dominikanerordenen. Da han gjorde sig skyldig i Kætterier,
og da hans Trang til at røre sig frit gjorde Klostertvangen
uudholdelig, flygtede han, aflagde Munkedragten og begyndte
nu et uroligt Studie- og Vandreliv, der førte ham til Schweiz,
Frankrig, England og Tyskland. Han optraadte som Lærer i
Toulouse, Paris, Oxford og Wittenberg, men fandt ikke et
blivende Sted, dels paa Grund af de gamle Skolers Modstand,
dels paa Grund af hans egen rastløse og lidenskabelige Aand.
Dog fandt han Tid til at affatte sine geniale Værker, af hvilke
navnlig de i London 1584 udgivne italienske Dialoger (der
nu foreligge i tyske Oversættelser) ere at fremhæve. Da han
ansaa en Forsoning med den katolske Kirke for mulig, og da
han havde større Uvillie mod den protestantiske Kirke end
mod den katolske, vendte han tilbage til Italien, hvor han vilde
leve i Ro for sine Studier. Men han blev fængslet af Inkvisitionen
og efter et langt Fangenskab brændt i Rom i Aaret
1600.

Bruno priser Copernicus for den høje og frie Aand, der
havde hævet ham over Sansernes Vidnesbyrd, ved hvilket den
store Mængde bliver hængende. Men han bestrider, at Verden
% --- p. 13 ---
\setcounter{footnote}{0}%
\opage{13}kan ende med Fixstjernesfæren, og støtter sig herved dels til
erkendelsesteoretiske, dels til religionsfilosofiske Grunde.

Kun tilsyneladende vidne Sanserne om et absolut Verdenscentrum
og et absolut Verdenshele. Hvergang vi skifte Standpunkt,
faa vi et nyt Verdenscentrum og en ny Verdensgrænse,
saa at hvert Sted i Verden snart kan betragtes som Midtpunkt,
snart som Grænsepunkt. Og det gaar med Tænkningen
som med Sansningen: Tal kan føjes til Tal og Tanke til Tanke
uden Grænse. Ethvert Sted er bestemt ved sit Forhold til
andre Steder, som ethvert Tal ved Forholdet til andre Tal og
enhver Tanke ved Forholdet til andre Tanker. Enhver Bevægelse
tager sig forskelligt ud efter Standpunktet, i Forhold
til hvilket den opfattes. Da Tiden maales ved Bevægelsen,
bliver Tidsbestemmelsen ogsaa stedse forholdsbestemt. Heller
ikke Tyngde og Lethed ere nu absolute Egenskaber, men
gælde i Forhold til de større Legemer, Elementerne høre til.
Ogsaa den gamle Modsætning mellem Himmel og Jord maa
bortfalde. Vi have ikke Ret til at antage, at Fixstjernerne hver
for sig ere ubevægede og alle i samme Afstand fra os; det
kunde jo gaa os her, som naar fjerne Skibe synes os at være
ubevægelige og alle lige langt borte. Det er i hvert Tilfælde
Noget, der maa nærmere undersøges. Og de faste Sfærer maa
forkastes, fordi Tycho Brahe har paavist, at Kometerne ikke
(som man før troede) ere Meteorer i Jordatmosfæren, men
gaa tvers gennem Himmelrummet og altsaa forlængst maatte
have knust Sfærerne. Stjernerne maa da ogsaa ligesaavel som
Jorden kunne bevæge sig uden Støtte.

Religionsfilosofisk gaar Bruno ud fra Guddommens Uendelighed.
Virkningen maa være uendelig, naar Aarsagen er det.
Den guddommelige Fylde kunde ikke komme til Udfoldelse
i en begrænset Verden.

Som det gaar med Modsætningen mellem Himmel og Jord,
gaar det ifølge Bruno med alle Modsætninger, hvormed vi
arbejde, fordi vi kun kunne tænke ved at stille det Ene overfor
det Andet. Hverken Modsætningen mellem Stof og Form,
eller den mellem Mulighed og Virkelighed, eller den mellem
Materie og Aand tør vi tillægge absolut Gyldighed. I det Væ%
% --- p. 14 ---
\setcounter{footnote}{0}%
\opage{14}sen, der ligger til Grund for Alt, men for hvilket ingen positiv
Tankebestemmelse kan findes, maa alle disse Modsætninger
falde bort. Guddommen kan ikke have sit Sæde i
„Himlen“, da jo Modsætningen mellem Himmel og Jord ikke
gælder, men ethvert Punkt kan ses baade som Midtpunkt og
som Grænsepunkt. Den virker i vort Inderste, i vor Stræben
efter Selvhævdelse; den er Sjælen i vor Sjæl, som den er
Sjælen i den hele Natur, og gennem Rummet bringer den alt
Existerende i Vexelvirkning. Men den æres bedre ved Taushed
end ved Tale.

Brunos Etik staar i nøje Sammenhæng med hans Verdensanskuelse.
Blandt Dyderne stiller han Sandhedskærlighed og
Arbejdsiver i første Linie. Intet Maal kan være det sidste,
men ethvert Maal føder ny Stræben. Hvor der er vakt en
heroisk Begejstring (furore eroico), lader Mennesket sig ikke
afskrække fra at stræbe efter et højt Maal, fordi der er Smerte
og Fare forbundne dermed; de ere kun Onder for den sanselige
Betragtning. Jo højere Maalet er, des større Smerte er
mulig. Men det heroiske Menneske glæder sig, naar den ædle
Ild er optændt i det, selvom Maalet ikke skulde naas.

\sectionhead{3. Den nye Videnskab.}

Kampen om det nye Verdensbilledes Gyldighed kunde kun
føres til Ende, naar det kunde godtgøres som en nødvendig
Følge af fysiske Naturlove. Dette skete ved \textsc{Isaac} \textsc{Newtons}
Paavisning af Tyngden som Verdenslov (1687). Men Forudsætning
herfor var igen en ny videnskabelig Metode; det er
en saadan, den nyere Naturvidenskab skylder sin Tilblivelse.
Istedetfor Spekulation og Konstruktion paa den ene Side og
blot Opsamling af Iagttagelser paa den anden Side traadte nu
nøjagtig Maaling, kvantitativ Bestemmelse af Fænomenerne og,
paa Grundlag deraf, streng Beregning og Slutten, hvis Resultater
igen kunde bekræftes ved nye maalelige Fænomener.
Den moderne Naturvidenskab blev til ved en saadan Vexelvirkning
mellem Induktion og Deduktion, Iagttagelse og matematisk
Beregning. Det er Kepler og Galilei, Indførelsen af
% --- p. 15 ---
\setcounter{footnote}{0}%
\opage{15}denne Metode skyldes, efterat allerede den berømte Kunstner
\textsc{Leonardo} \textsc{da} \textsc{Vinci} (1452---1519), Renaissancens mest alsidige
Geni, havde udtalt Tanker i samme Retning, som dog først
langt senere bleve bekendte.

Hos \textsc{Johannes} \textsc{Kepler} (1571---1630), den berømte Astronom,
foregik en interessant Udvikling fra en teologisk-poetisk
Verdensanskuelse til en exakt naturvidenskabelig Opfattelse.
Han gik først ud fra, at Universet som en Aabenbaring af
Guddommen maatte være ordnet efter simple geometriske
Forhold, som han endog mente at kunne konstruere paa Forhaand.
Senere fastholdt han kun den almindelige Overbevisning,
at der maatte gælde visse matematiske Love for Planeternes
Bevægelser, og paa Grundlag af Tycho Brahes
Iagttagelser fandt han disse Love (idet han bl. A. paaviste, at
Marsbanen var en Ellipse). Og medens han tidligere havde
forklaret Planeternes Bevægelse ved Antagelsen af særegne
Planetsjæle, opstillede han senere (i sit Skrift om Marsbanen)
det Princip, at de Aarsager, man antog, selv skulde kunne
paavises i Erfaringen; en saadan Aarsag kaldte han \textit{vera
causa.} Planetsjælene vare ikke veræ causæ. --- Foruden den
teologiske Begrundelse af den kvantitative Metode har han
ogsaa en psykologisk Begrundelse: den matematiske Erkendelse
er den klareste og sikreste, vi besidde, og bør derfor
anvendes saa vidt som muligt. I kvalitativ Henseende kunne
Fænomenerne stille sig forskelligt for forskellige Individer;
men ved kvantitativ Bestemmelse hører Uenigheden op. ---
Endeligt viser Erfaringen, at alle materielle Fænomener have
kvantitative, særligt geometriske Egenskaber: ubi materia, ibi
geometria.

Den egentlige Grundlægger af den moderne Naturvidenskab
er dog \textsc{Galileo} \textsc{Galilei} (1564---1642) ved sin klare
Bevidsthed om Vexelvirkningen mellem Induktion og Deduktion,
Iagttagelse og Beregning. Iagttagelsen kan ikke omfatte
alle mulige Tilfælde; vi maa derfor holde os til visse udvalgte
Fænomener, beskrive og maale dem saa nøje som muligt,
paa Grundlag deraf ved Slutning og Beregning opstille en
Hypotese og undersøge, om de Konsekvenser, der kunne
% --- p. 16 ---
\setcounter{footnote}{0}%
\opage{16}udledes af den, stemme med, hvad Iagttagelsen viser. Baade
ved Opstillingen og ved Prøvelsen af Hypotesen ere kvantitative
Bestemmelser nødvendige, og Galilei opstillede derfor det
Princip: maal Alt hvad maaleligt er, og gør det maaleligt, som
endnu ikke er det!

Allerede Kepler havde lært, at et Legeme ikke af sig selv
kan gaa over fra Hvile til Bevægelse. Galilei gaar et Skridt
videre. Det syntes ham det Simpleste (og ligesom Copernicus,
Bruno og Kepler betragtede han Simpelhedens Lov som en
Verdenslov), at en Ting bliver i den Tilstand, i hvilken den
er, naar ingen ydre Aarsager gribe ind. Et Legeme kan derfor
ikke af sig selv forandre sin Bevægelse. Kunde alle ydre
Paavirkninger holdes borte, vilde Bevægelsen stadigt fortsættes
med samme Hastighed. Ved Forsøg godtgjorde Galilei
derpaa, at Kugler, der trilledes ad Pergamentrender, fortsatte
deres Løb des længere, jo glattere Renderne vare, og jo glattere
og regelmæssigere de selv vare. Ogsaa Kredsbevægelsen,
som han ansaa for simpel, vilde efter hans Mening stadigt
fortsættes ved Fraværelse af ydre Hindringer, --- ikke blot
den retliniede Bevægelse.

Disse fysiske Teorier udviklede Galilei i \textit{Discorsi delle
nuove scienze} (1638). Nogle Aar forud havde han i sin
\textit{Dialogo sopra i due massimi sistemi del mondo} (1632) sammenlignet
det ptolemæiske og det kopernikanske Verdenssystem
indbyrdes. Han mente at have afvejet Mulighederne
for og imod uden at udtale sig afgørende; men hans egen
Opfattelse var tydelig, ikke mindst ved den Ironi, hvormed
Repræsentanten for den gamle Opfattelse blev behandlet. Dette
medførte hans Livs Katastrofe. Allerede tidligere (efter Opdagelsen
af Jupitermaanerne og af Solpletterne) havde han offentligt
udtalt sig til Gunst for det kopernikanske System. Og
da i Aaret 1616 Inkvisitionen satte Copernicus' Bog paa Listen
over forbudte Bøger, skal Galilei have lovet ikke at hylde
eller udbrede den nye Lære. Inkvisitionen saa i hans Dialog
et Brud paa dette Løfte og den halvfjerdsindstyveaarige Mand
maatte, under Trusel om Pinebænken, afsværge den ugudelige
Lære, at Jorden bevægede sig og ikke var Verdens Midtpunkt.
% --- p. 17 ---
\setcounter{footnote}{0}%
\opage{17}Indtil sin Død stod han under Inkvisitionens Opsigt, og hans
senere Værker maatte trykkes hemmeligt i Udlandet.

Galilei fremhæver stærkt, at Kopernikanismen indskærper,
at vi ikke maa forvexle vor Opfattelse med selve Tilværelsen:
„Tænk Jorden borte, saa staar Solen ikke op og gaar ikke
ned, og der gives hverken Horizont eller Meridian, hverken
Dag eller Nat!“ Denne Tanke udvider han til hele den sanselige
Naturopfattelse under Indflydelse af den Overbevisning,
at vi kun faa sand Erkendelse ved at holde os til Tingenes
kvantitative Egenskaber. Kun Størrelse, Figur og Bevægelse
eller Hvile kunne vi tilskrive selve Tingene; de ere de første
og reale Egenskaber. Derimod er det en sanselig Fordom at
betragte Smag, Lugt, Farve, Varme o. s. v. som absolute Egenskaber
ved Tingene. De ere Navne, som vi give Tingene,
naar de fremkalde visse Fornemmelser hos os; de have deres
Sæde i vort sanselige Legeme og bortfalde, naar dette falder
bort. Denne Lære, i hvilken den mekaniske Naturvidenskabs
Princip er udtalt, blev af stor Betydning for den følgende
Tids erkendelsesteoretiske Undersøgelser.

\closerule

\chapterhead{II.}{SYSTEMDANNELSER.}

De nye Interesser, Synspunkter og Opdagelser, som Renaissancetiden
havde bragt, maatte naturligt fremkalde en Bestræbelse
for paa det saaledes givne Grundlag at udvikle en
ny, i sig selv sammenhængende Verdensanskuelse. Der maatte
røre sig en Trang til at drage de sidste Konsekvenser af de
nye Anskuelser. Der rører sig til hver Tid, mere eller mindre
udpræget, en Trang i den menneskelige Aand til at sammenfatte
sin Viden i en Helhedsanskuelse. Særligt stillede sig nu
den Opgave at forene det nye Verdensbillede og den nye
% --- p. 18 ---
\setcounter{footnote}{0}%
\opage{18}Videnskab med Aandslivets øvrige Indhold. Bruno havde allerede
arbejdet i denne Retning; men den nye mekaniske Naturopfattelse,
ved hvilken mange Problemer stillede sig anderledes
end før, kunde ikke faa Indflydelse paa hans Problemstilling.

a. \textsc{Réné} \textsc{Descartes} (1596---1650) kan kaldes den nyere
Filosofis egentlige Grundlægger, fordi han først metodisk opsøger
Erkendelsens Forudsætninger, og fordi han var den
Første, hos hvem den mekaniske Naturforklaring fik en afgørende
Indflydelse paa Problemstillingen. Han var Søn af en
fransk Adelsmand, og hans økonomiske Uafhængighed gjorde
ham det muligt at hellige sit Liv til Tænken og Forsken. I
sin \textit{Discours de la méthode} (1637) har han givet en interessant
filosofisk Selvbiografi. Han fik sin første Undervisning
paa et Jesuiterkollegium. Skønt han havde Tidens bedste
Pædagoger til Lærere, var han dog ikke tilfreds med Resultatet
af sin Skoletid. Han vidste Mangt og Meget, men Sammenhæng
og klar Begrundelse manglede. Han kastede sig da
ind i Verdenslivet og besluttede senere at lære det virkelige
Liv at kende gennem Krigstjeneste og Rejser. Stedse blev
han dog dragen tilbage til Filosofien. Da han i Aaret 1619
laa i Vinterkvarter i Bayern, havde han en Art videnskabelig
Vækkelse. Den Vej, han skulde gaa som Forsker, stod pludseligt
for ham: Vore Forestillingers forvirrende Mangfoldighed
vil forsvinde, naar vi begynde med de simpleste og klareste
Tanker og derpaa Skridt for Skridt gaa over til de sammensatte
Spørgsmaal, stadigt ordnende vore Tanker saaledes, at
de følgende kunne afledes af de foregaaende. Denne Ide ledede
ham baade ved hans matematiske Arbejder (Grundlæggelsen
af den analytiske Geometri) og ved Udarbejdelsen af hans
filosofiske System. Efter nogle Studieaar i Paris drog han til
Holland for at leve mere i Ro, og fordi der i Frankrig blev
udstedt strenge Forbud mod Afvigelse fra den middelalderlige
Filosofi. Men da de hollandske Teologer ikke mindre end de
katolske betragtede hans Filosofi som farlig, modtog han en
% --- p. 19 ---
\setcounter{footnote}{0}%
\opage{19}Indbydelse fra Dronning Kristina og levede sine sidste Aar i
Stockholm.

For at finde de første Principer for Erkendelsen undersøger
Descartes (udførligst i \textit{Meditationes} 1641), om der er
Noget, som ikke kan drages i Tvivl. Om enhver Genstand,
ethvert Indhold af vore Tanker kunne vi tvivle. Enhver Iagttagelse
og enhver Antagelse kan være Illusion. Men selv den
mest radikale Tvivl er en Tænken, en Bevidsthedsakt; falder
al Tænken bort, saa falder ogsaa al Tvivl bort. Tvivler jeg,
saa véd jeg dog, at jeg tvivler. Selve Tanken kan jeg altsaa
ikke tvivle om, selvom jeg tvivler om alle dens Resultater.
Ordet Tænken tager Descartes i vid Betydning, om al Art af
Bevidsthed, Følen og Villen. Det, at jeg tænker, viser, at jeg
er til (som tænkende): je pense, donc je suis (cogito, ergo
sum)! Dette er dog ingen Slutning, men en umiddelbar Bevidsthed
(om at være bevidst). --- Den Klarhed og Tydelighed,
hvormed jeg véd, at jeg tænker, er nu for Descartes en
Maalestok: alle Sætninger, der ere ligesaa indlysende som
denne, maa være gyldige. Og der er to vigtige Sætninger, der
opfylde denne Fordring: den Sætning, at Alt har en Aarsag,
og den, at der ikke kan indeholdes Mere i Virkningen end i
Aarsagen. Af disse Sætninger udleder Descartes et Bevis for
Guds Tilværelse: medens alle Forestillinger om endelige Ting
og Væsener kunne finde deres Forklaring i os eller udenfor
os, er dette ikke muligt for Forestillingen om et uendeligt
Væsen; den kan kun skyldes en uendelig Aarsag. Desuden
vilde det være en Selvmodsigelse, om Guds Existens forkastedes:
ti Begrebet om Gud er Begrebet om et fuldkomment
Væsen, og til Fuldkommenhed hører ogsaa Existens.
Det følger altsaa af Guds Begreb, at han maa existere.

Først nu, da Descartes (ganske vist ved Aarsagssætningens
Hjælp) mener at have bevist Guds Tilværelse, er vor Erkendelses
Sandhed sikret for ham; ti et fuldkomment Væsen kan
ikke bedrage. Ligesom Kepler giver han altsaa sin Lære et
teologisk Grundlag. Dog er det Gudsbegreb, med hvilket han
her opererer, ikke et rent religiøst Begreb. Det udtrykker
Principet for den lovmæssige Sammenhæng i Naturen. Alt
% --- p. 20 ---
\setcounter{footnote}{0}%
\opage{20}hvad der skal anerkendes som sandt, maa finde sin Plads i
denne Sammenhæng, der gør det muligt for os at skelne mellem
Drøm og Virkelighed.

Ved Hjælp af Gudsbegrebet begrunder Descartes derpaa
de første Principer for Naturvidenskaben. \textit{(Principia} Philosophiæ.
1644). Kun tre Egenskaber have vi Ret til at tillægge
de materielle Ting: Udstrækning, Delelighed og Bevægelighed.
Og vi faa den klareste Erkendelse af Alt hvad der foregaar
i den materielle Verden, naar vi holde os til disse tre Egenskaber;
Sansekvaliteterne maa derimod skrives paa det opfattende
Subjekts Regning. Herved udtaler Descartes den
mekaniske Naturvidenskabs Grundtanke, som han synes at
være naaet til uafhængigt af Galilei, men sandsynligvis under
Indflydelse af Kepler. --- Af Guds Uforanderlighed (der er
en Del af hans Fuldkommenhed) følger nu, at et Legeme forbliver
i sin Tilstand, naar en ydre Aarsag ikke paavirker den.
Og ikke blot Inertisætningen begrundes saaledes teologisk.
Af den guddommelige Uforanderlighed følger fremdeles Sætningen
om Bevægelsens Bestaaen (en ufuldkommen Forløber
for Sætningen om Energiens Bestaaen): den ved Skabelsen i
Materien indlagte Sum af Bevægelse kan fordeles forskelligt
til forskellige Tider, men ingen Bevægelse gaar tabt, og ingen
helt ny Bevægelse opstaar.

Ud fra disse Forudsætninger kræver Descartes Alt i Naturen
forklaret paa mekanisk Maade. Selv har han kun fyldestgjort
denne Fordring ved Hypoteser, der ikke altid have
staaet deres Prøve. --- Som Verden idethele, saaledes er ogsaa
Organismen en Maskine. En Bekræftelse paa denne Opfattelse
fandt Descartes i Harvey's Opdagelse af Blodets
Kredsløb, der viste, at Blodets Bevægelse skyldtes mekaniske
Aarsager. Selv har han anvendt en mekanisk Forklaring paa
Nervesystemet, idet han beskriver den uvilkaarlige Bevægelse,
som nu kaldes Reflexbevægelse, og ved hvilken et Indtryk,
uden at nogen Bevidsthed behøver at være tilstede, umiddelbart
udløser Muskelvirksomheden. Bevidsthed antog han kun
hos Mennesket, medens han betragtede Dyrene som blotte
Maskiner.

% --- p. 21 ---
\setcounter{footnote}{0}%
\opage{21}Menneskets Sjæl staar i Vexelvirkning med Hjernen, særligt
med en bestemt Del (Koglekirtlen, glandula pinealis), der
efter hans Mening laa centralt i Hjernen, og der ikke, som
andre Hjernedele, var dobbelt. Paa denne Kirtel støder „Livsaanden“
(den fine Luft, der efter Tidens fra Oldtiden stammende
Fysiologi strømmede gennem Nerverne), og Stødet
forplanter sig derpaa til Sjælen, hvorved Sansefornemmelser
opstaa. Naar Sjælen igen støder til „Kirtlen“, frembringes der
en Ændring i „Livsaandens“ Retning, der igen foraarsager en
Muskelbevægelse, altsaa Udførelse af en Villiesakt. Descartes
kommer ved denne klare og konsekvente Gennemførelse af
den populære Forestilling om en Vexelvirkning mellem Sjæl
og Legeme i Strid med sin egen Lære om Bevægelsens Bestaaen:
der forsvinder Bevægelse, naar Koglekirtlen støder til
Sjælen, og der opstaar Bevægelse, naar Sjælen støder til
Koglekirtlen. Selvom Sjælen kun ændrer Bevægelsens Retning,
strider det mod Inertisætningen.

Sjæl og Legeme ere for Descartes to forskellige Væsener
(Substanser), fordi de frembyde forskellige Egenskaber, Tænkning
(Bevidsthed) og Udstrækning, der ikke kunne udledes
af hinanden. Hvorledes to saa forskellige Væsener kunne
paavirke hinanden, blev da ogsaa en Vanskelighed for Descartes
og hans Efterfølgere. ---

Kartesianismen var den første Form, i hvilken den nye
Tids Tænken blev tilgængelig for større Kredse. Ved den
strenge Fordring om mekanisk Naturforklaring betegnede han,
trods sine ofte uheldige Hypoteser, et stort Fremskridt, og
han vakte en levende Bevægelse paa Naturvidenskabens Omraade.
Ved sin Spiritualisme (Antagelsen af Sjælen som et fra
Legemet forskelligt Væsen) og ved sit Forsøg paa at forbinde
Teologi og Videnskab kunde han --- trods heftig Modstand ---
finde Indgang i religiøse Kredse. Ved sin store Klarhed kunde
hans Lære let fremstilles populært, og han vandt stor Tilslutning
i Frankrig, Holland og Tyskland.

b. \textsc{Thomas Hobbes} (1588---1679) har gjort det første selvstændige
Forsøg paa at betragte den mekaniske Naturlære
% --- p. 22 ---
\setcounter{footnote}{0}%
\opage{22}som den eneste Videnskab, --- dens Synspunkter som de
eneste, ud fra hvilke Tilværelsen kan og skal opfattes. Denne
som Tænker og Polemiker energiske, i sit Liv milde og frygtsomme
Mand var ligesom Descartes utilfreds med den middelalderlige
Skole, han havde gennemgaaet. Han beskæftigede
sig en Tid mest med literære Studier. De urolige Forhold i
England vakte hans Interesse for politiske og etiske Spørgsmaal,
og da han havde lært de nye naturvidenskabelige Synsmaader
at kende (især efterat han paa en Rejse i Italien var
bleven en personlig Ven af Galilei) førtes han over til almindelige
filosofiske Studier. Han levede som Privatmand hos
den adelige Familie Cavendish, først som Huslærer, derefter
som Husven.

Da Revolutionen brød ud, drog Hobbes til Frankrig, hvor
han senere en Tid var Lærer for den landflygtige Karl II.
Da Cromwell havde skaffet Ro i Landet og givet Ytringsfrihed,
vendte han tilbage. ---

Hobbes synes ad sin egen Vej at være kommen til den
Overbevisning, at al Forandring er Bevægelse. Da han var
tilstede ved en Diskussion om, hvad Fornemmelse er, kom
han paa den Tanke, at naar Alt i Naturen var ubevægeligt
eller i ensformig Bevægelse, vilde der ikke opstaa nogen Fornemmelse.
En Bevægelsesforskel (diversitas motuum) er derfor
Betingelse for Fornemmelse. Stadigt at fornemme Et og
det Samme og slet ikke at fornemme kommer ud paa Et.
Paa denne Sætning grundede Hobbes sin Psykologi. Da han
saa ogsaa tænkte over, hvad Forandring var, fandt han, at
den kun bestod i Bevægelse, og deraf fulgte igen, at Sansekvaliteterne
(Farve, Klang o. s. v.) kun have rent subjektiv Gyldighed.
Han beraaber sig herved ogsaa paa Erfaringen, f. Ex.
paa at der ved Slag paa Øjet opstaar Lys, skønt der faktisk
intet ydre Lys er.

Den nævnte Sætning, at al Forandring bestaar i Bevægelse,
er for Hobbes Videnskabens Hovedprincip. Som saadant kan
den ikke bevises, men den vil anerkendes, naar man blot aflægger
alle Fordomme. Paa lignende Maade opstiller han en
Række andre Grundsætninger: Aarsagssætningen, Inertisæt%
% --- p. 23 ---
\setcounter{footnote}{0}%
\opage{23}ningen, den Sætning, at Bevægelse kun kan opstaa af Bevægelse
og ikke frembringe Andet end Bevægelse. Naar nu disse
Grundsætninger skulle forklare Alt, saa maa Alt bestaa i Bevægelse.
I sin systematiske Fremstilling begynder Hobbes
med Læren om Legemet \textit{(corpus)} i Almindelighed, en Fremstilling
af de geometriske, mekaniske og fysiske Bevægelseslove,
gaar derefter over til Læren om Mennesket \textit{(homo),}
d. v. s. om de Bevægelser, der foregaa i Mennesket, de fysiologiske
og psykologiske Bevægelser, og ender med Læren om
Borgeren \textit{(civis),} d. v. s. om de Bevægelser i Menneskene, der
betinge deres gensidige Forhold og deres Samfundsliv. --- Den
friskeste Fremstilling af sine psykologiske og sociale Teorier
har han givet i et omtrent 1640 forfattet Skrift \textit{Elements of
Law} (ny Udgave ved Tønnies 1888). Senere udgav han dels
de statsfilosofiske Skrifter \textit{De cive} og \textit{Leviathan,} dels en systematisk
Fremstilling af Systemets to første Dele (De corpore
og De homine) under Navn af \textit{Elementa Philosophiæ} (1655).

De første Grundsætninger maa ifølge Hobbes opstilles
som vilkaarlige Definitioner. Men han er klar over, at Videnskaben
ikke kan nøjes med Deduktioner ud fra saadanne
Definitioner. Naar vi gaa over fra Geometri til Mekanik, derfra
til Fysik og derfra igen til Psykologi og Samfundslære optage
vi stadigt nye Forudsætninger og opstille Hypoteser, der gøre
Anvendelsen af de almindelige Grundsætninger paa de specielle
Omraader mulig. Men han tvivler ikke om, at Grundsætningerne
kunne anvendes helt igennem. Naar Alt er Bevægelse,
maa Alt være legemligt. En ulegemlig Ting er en Uting. Videnskaben
kan kun undersøge endelige Ting, da kun saadanne
kunne være i Bevægelse fra et Sted til et andet. Om Verden
som afsluttet Hele kunne vi ingen Viden have. Alle Spørgsmaal
angaaende Verden som Helhed føre over til det Ubegribelige
og kunne kun afgøres ved Tro, ikke ved Viden.
Hverken om Verdens Opstaaen, Grænse eller Ende kan Videnskaben
lære Noget. Den højeste Videnskab, Visdommens Førstegrøde
(primitiæ sapientiæ) er forbeholdt Teologerne, ligesom
i Israel Høstens Førstegrøde var forbeholdt Præsterne. ---

Naar Alt er Bevægelse, maa Fornemmelsen, og Bevidst%
% --- p. 24 ---
\setcounter{footnote}{0}%
\opage{24}heden idethele, ogsaa være Bevægelse: „Fornemmelsen er
\textit{ikke Andet end} en Bevægelse i det fornemmende Legemes
Smaadele“; „Bevidsthed (mens) er \textit{ikke Andet end} en Bevægelse
i nogle Dele af et organisk Legeme“. Bevægelsen er
det Virkelige; Bevidstheden er kun Bevægelsens Fremtræden
(apparition) for os. Saaledes er Lystfølelse i Virkeligheden
(really) kun en Bevægelse i Hjertet, Tænken kun en Bevægelse
i Hovedet. Psykologien er altsaa en Del af den almindelige
Bevægelseslære. Dog skelner Hobbes jo udtrykkeligt
selve Bevægelsen fra dens Fremtræden eller Apparition for
os. Han udtaler endog, at af alle Fænomener er det allermærkeligste
det, at Noget overhovedet kan blive Fænomen, altsaa
træde frem for et Subjekt. Der maa altsaa dog gives Andet i
Verden end Bevægelse.

I sin specielle Psykologi fremhæver Hobbes den store
Betydning, Sammenligning har, allerede for selve Sansningens
Vedkommende, og han viser, hvorledes Forestillinger opstaa
og forbindes efter Associationslove. Han har lagt Grunden til
den engelske psykologiske Skole. --- Af særlig Interesse er
hans Lære om Følelser og Lidenskaber. Lyst eller Ulyst opstaa,
naar en Forandring fremmer eller hæmmer vor Stræben
efter Selvopholdelse. Et absolut Maal, der skulde kunne naas
en Gang for alle, er utænkeligt; al Stræben og Følelse vilde
da falde bort. Det højeste Gode maa bestaa i uhindret Fremgang
efter stedse højere Formaal. Ti Liv er jo Bevægelse! ---
De specielle Former for Stræben og Følelse udvikle sig af
Magt- og Afmagtfølelsen. Forestillingen om Fremtiden spiller
nemlig en stor Rolle: mener jeg nu at have Midler til fremtidig
Tilfredsstillelse og til Overvindelse af Hindringer, opstaar
Magtfølelsen, i modsat Tilfælde Afmagtfølelsen. Til de Midler,
jeg regner paa, kan ogsaa høre Hjælp fra Venner eller fra
Gud. Da min Stræben efter Selvopholdelsen let krydses af
Andres Stræben, spiller Sammenligning med dem en stor
Rolle. Livet er et Kapløb. Vi hovere, naar vi komme forud
for Andre, føle Ydmygelse, naar vi blive tilbage, --- Haab,
naar vi ere i Aande, Fortvivlelse, naar vi ere trætte, Vrede,
naar vi møde en pludselig Hindring, Stolthed, naar Hindrin%
% --- p. 25 ---
\setcounter{footnote}{0}%
\opage{25}ger overvindes; --- vi le, naar vi se Andre falde, græde, naar
vi selv falde; --- vi føle Medlidenhed, naar vi se En, vi ønske
vel, blive tilbage, Kærlighed, naar vi hjælpe Andre. Først med
Døden ender Kapløbet.

Under dette store Kapløb erfares det, at de forskellige
Menneskers Selvopholdelsesbestræbelser ikke ere i gensidig
Harmoni. Der opstaar Strid, og der kan stadigt frygtes Overgreb.
Naturtilstanden, d. v. s. Menneskelivets Tilstand som den
vilde være uden Statsordning, er en Krig af Alle mod Alle
(bellum omnium contra omnes). Her gælder kun de Enkeltes
Trang og Magt, saavidt de kunne brede sig. Frygt, Had og
uroligt Magtbegær raade. Men i rolige Øjeblikke (sedato
animo) ville Menneskene indse, at de kunne opnaa større
Fordele ved Sammenslutning og Fællesskab end ved Strid.
Derved opdage de den moralske Grundlov: Freden skal
hævdes, og kan den ikke naas, skal Krigen føres i Fællesskab!
Af denne Grundsætning kunne de enkelte Dyder og
Pligter udledes, ti Ordholdenhed, Taknemmelighed, Tjenstvillighed,
Overbærenhed, Billighed og Selvbeherskelse ere nødvendige
Egenskaber, hvis Freden skal raade og Samfundet
bestaa. Ligeledes følger af hin Grundsætning den Regel, at
man ikke maa gøre mod Andre, hvad man ikke vil, at de
skulle gøre mod os.

Men til praktisk Gennemførelse og Hævdelse af disse Love
og Regler er en stærk Statsmagt nødvendig, Naturtilstandens
Frihed maa opgives. Dette sker ved en udtrykkelig eller stiltiende
Kontrakt, ved hvilken alle de Enkelte paa en Gang
give Afkald paa deres Frihed og forpligte sig til ubetinget
Lydighed mod en Autoritet (en Fyrste eller en Forsamling).
Kun ved en saadan absolut Underkastelse kan Freden sikres.
Fra nu af bliver Regeringens Villie Et med Folkets Villie.
Hobbes er den absolute Suverænitets naturalistiske Teoretiker.
Enhver Begrænsning (ved Stænder eller Kirke) og enhver
Deling (mellem Fyrste og Parlament) vilde indeholde Spiren
til Krig og kunne føre tilbage til Naturtilstanden. Ogsaa paa
det moralske og religiøse Omraade skal Suverænens Villie
være afgørende. Fremfor Alt skal Suverænen afgøre, hvorledes
% --- p. 26 ---
\setcounter{footnote}{0}%
\opage{26}Gud skal dyrkes; ellers vilde det, der for den Ene er Tilbedelse,
være Gudsbespottelse for den Anden, og der vilde
være en stadig Kilde til Strid og Opløsning. Af samme Grund
maa de første Definitioner af Godt og Ondt fastslaas af Suverænen.
Hobbes underforstaar herved stadigt, at Suverænen
udtrykker Folkets Villie til at holde Naturtilstanden nede og
hævde Freden. Han er en Forløber for det attende Aarhundredes
oplyste Despotisme. Hvad han bekæmpede, var først
og fremmest Klasseherredømmet og Kirkevælden, og idet han
haaber paa Folkets fremskridende Oplysning (paulatim eruditur
vulgus), venter han ogsaa, at en oplyst Suveræn engang vil
hævde det oplyste Folks Villie.

c. \textsc{Baruch} \textsc{Spinoza} (1632---1677) har i sin \textit{Ethica more
geometrico demonstrata} (1677) gjort det dybsindigste Forsøg
paa at indarbejde den nye Naturopfattelses Grundtanker i en
almindelig Verdensanskuelse. Dette Værk har trods sin abstrakte
Form ikke en upersonlig og rent teoretisk Karakter.
Hos Spinoza gaa Liv og Tænken i Et. Den klare Tænken
var for ham Vejen til aandelig Frihed, til den højeste Form
for personligt Liv. Det var hans Bestræbelse at se alle Tilværelsens
Sider og Former i indre Sammenhæng. Det Store
i hans Tænken er den fulde Gennemførelse af forskellige
Tankegange hver for sig, uden Afbrydelse og uden at gaa
paa Akkord, og dernæst Paavisningen af, at enhver af de i
sig sluttede Tankerækker kun betyder en Side af eller en
Grundform for den uendelige Tilværelse. Saaledes forsøger
han at paavise, at Enhed og Mangfoldighed, Aand og Materie,
Evighed og Tid, Værdi og Virkelighed hver for sig paa sin
Vis udtrykker den hele Tilværelse, saa at der er en indre
Identitet imellem dem.

Spinoza, der var født i Amsterdam i en jødisk, fra Spanien
flygtet Købmandsfamilie, begyndte som jødisk Teolog, til hvem
hans Trosfæller knyttede store Forhaabninger. Men han stillede
sig mere og mere kritisk overfor de fædrene Trosforestillinger
og blev tilsidst højtideligt udstødt af Synagogen. Han levede nu
en Tid paa Landet, derefter i Nærheden af Leyden og tilsidst i
% --- p. 27 ---
\setcounter{footnote}{0}%
\opage{27}Haag, beskæftiget med Studier og med Udarbejdelsen af sine
Skrifter. Sit tarvelige Livsophold skaffede han sig ved at slibe
optiske Glas. En Kreds af unge Venner stod i livlig Tankesamkvem
med ham og studerede hans „Etik“ i Manuskript.
Hans Liv er et ophøjet Exempel paa frejdig Resignation og
inderlig Hengivelse til intellektuelt Arbejde.

Den deduktive („geometriske“) Metode, Spinoza anvender
i sit Hovedværk, er ikke den, ved hvis Hjælp han er kommen
til sin Lære. Af en ufuldendt Afhandling „om Forstandens
Forbedring“, saavel som af Antydninger i Hovedværket,
ser man, hvorledes han har opfattet Erkendelsens Udvikling.
Denne begynder med spredt og tilfældig Erfaring
(experientia vaga), i hvilken de forskellige Fænomener sammenstilles
efter Regler, der uvilkaarligt frembyde sig og blot
bero paa, at man ingen Undtagelser har fundet. Videnskaben
(ratio) anstiller derimod en nøje Sammenligning af de givne
Fænomener. Den gaar ud fra Erfaringen, men søger Noget,
der ytrer sig saavel i Naturen idethele som i de enkelte Dele:
almindelige Love, der gælde overalt. Som Exempler nævner
Spinoza Bevægelsens Love paa det materielle Omraade og
Forestillingsassociationens Love paa det aandelige Omraade.
Kun i saadanne Love kan man finde en Afslutning for sine
Tanker, hvorimod de enkelte Fænomeners Rækker stadigt
kunne fortsættes, idet det, der i en Henseende er Aarsag, i
en anden Henseende er Virkning, og omvendt. Et Absolut, i
hvilket Tanken kan hvile, er kun at finde i Aarsagsrækkens
Lov, ikke i dens formentlige Begyndelse eller Ende. Et saadant
Absolut kalder Spinoza Substans: det, som bestaar i sig
selv og forstaas ved sig selv, idet dets Begreb ikke forudsætter
noget andet Begreb. Spinozas Substans, Afslutningen for hans
Tænken, er altsaa Principet for Tilværelsens Lovmæssighed.
Den højeste Erkendelsesart er ikke selve Erkendelsen af Naturlovene.
Over experientia vaga, der kun giver enkelte Tilfælde,
og ratio, der giver almindelige Love, staar for Spinoza
den anskuende Viden (scientia intuitiva), ved hvilken det enkelte
Fænomen umiddelbart fremstilles som ejendommeligt
Led i den store Natursammenhæng, --- det enkelte Væsen
% --- p. 28 ---
\setcounter{footnote}{0}%
\opage{28}som Del af Substansen, der aabenbarer sig i denne Natursammenhæng.
Denne højere Anskuelse kan kun naas, naar
man har gennemgaaet Erfaringens og Videnskabens Skole.
Den synes snarere at være en kunstnerisk Skuen end en
rent videnskabelig Opfattelse. Den giver et Billede af det
Enkelte set i Lyset af det Hele. Baade ved den anden Erkendelsesmaade
(ratio) og ved den tredie (scientia intuitiva)
betragte vi Tingene fra Evighedens Synspunkt (sub specie
æterni), d. v. s. ikke i deres Isolation og Tilfældighed, men
som Led i en omfattende Ordning.

Ligesom Descartes og Hobbes bygger Spinoza hele sin
Filosofi paa Aarsagssætningen, hvis Gyldighed for ham ligesom
for dem er selvfølgelig. Ganske særlig Vægt lægger han
paa, at Aarsag og Virkning ikke kunne være forskellige Ting:
„Naar to Ting Intet have fælles med hinanden, kan den ene
ikke være Aarsag til den anden; ti der vilde da ikke være
Noget i Virkningen, som ogsaa var i Aarsagen, og Alt i Virkningen
vilde da være opstaaet af Intet.“ At to Ting forholde
sig som Virkning og Aarsag, vil efter Spinoza sige, at den
enes Begreb kan udledes af den andens. Han skelner ikke
mellem Grund og Aarsag. Forholdet mellem Aarsag og Virkning
er for ham det samme som mellem Præmisser og Konklusion
i en Slutning. At Aarsagen \textit{i Tiden,} ikke blot i Tanken
gaar forud for Virkningen, lægger han ikke Vægt paa.

„Fra Evighedens Synspunkt“ forsvinder Tiden.

Aarsagen til en Ting kan enten ligge i den selv eller i
noget Andet. Hvad der har sin Aarsag i sig selv, er Substans
(Grundvæsen). Det bestaar i sig selv og kan forstaas ved sig
selv. Vi have allerede set, at Grundvæsenet for Spinoza aabenbarer
sig i Naturens indre Lovmæssighed. Det ligger til Grund
for alle Tings Bestaaen og Virken. Der kan kun gives én
Substans; ti gaves der flere, maatte de begrænse hverandre,
og ingen af dem kunde da siges at bestaa i sig selv og forstaas
ved sig selv. Dette Begreb, der for Spinoza danner Afslutningen
af hans Tænken, er for ham identisk med Gudsbegrebet
saavel som med Begrebet Natur. Men disse Begreber
maa da opfattes paa en anden Maade end den sædvanlige.
% --- p. 29 ---
\setcounter{footnote}{0}%
\opage{29}Naturen betyder en indre Kraft, der rører sig i Alt, ikke den
blotte Sum af alt Existerende. Og Gud kan ikke tænkes som
et personligt Væsen, da psykologiske Egenskaber som Forstand
og Villie ikke kunne anvendes paa et absolut og uendeligt
Væsen. Forstand forudsætter en fra den forskellig Genstand,
og Villie et Maal, hvorefter der stræbes; men saadanne
Forskelle kunne ikke gælde for det absolute Væsen. --- Ting,
der ikke have deres Aarsag i sig selv, kalder Spinoza \textit{modi}
(Enkeltvæsener, Fænomener). Modi have deres Aarsag i Substansen,
hvis Aabenbaringsformer de ere. De staa i ydre Aarsagsforhold
til hverandre indbyrdes, men hele denne i Erfaringen
givne Række af Aarsager og Virkninger har sin Aarsag
i Substansen, som er det indre Baand, den sammenholdende
Kraft i Fænomenernes Række.

Alt hvad der existerer, modi saavel som Substansen, fremtræder
i vor Erfaring med to \textit{Attributer} (Grundegenskaber
eller Grundformer): Aand og Materie (eller, som Spinoza
udtrykker sig, Tænken og Udstrækning). Da Substansen er
et uendeligt og fuldkomment Væsen, maa den have uendeligt
mange Attributer; men vi kende kun to af disse, da Erfaringen
ikke viser os flere. Ved Attribut forstaar Spinoza „det,
som Forstanden opfatter om Substansen som udgørende (konstituerende)
dens Væsen“. Ethvert Attribut maa derfor forstaas
ved sig selv ligesom Substansen, hvis Væsen det skal udtrykke.
Man kan derfor ikke forklare et Attribut ved et andet, Aand
ved Materie eller Materie ved Aand. Er man begyndt at betragte
Tilværelsen under Materiens Attribut, maa man vedblive
dermed og ikke springe over til et andet Attribut. Tilværelsens
Væsen (Substans og modi) lægger sig for Dagen i Materiens
Form saavel som i Aandens Form, og den ene Form kan
ikke udledes af den anden. Det følger allerede af den Maade,
% PRINTED AS IS (read at the image, odd spots o6)
paa hviken Spinoza opfatter Aarsagssætningen: Aarsag og
Virkning maa være ensartede, altsaa høre ind under samme
Attribut. At udlede en materiel Forandring af Sjælens Indvirken,
er derfor ifølge Spinoza egentligt at indrømme, at man
ikke forstaar det. --- Ved denne Opfattelse har Spinoza grundlagt
en ny Hypotese, der fører ud over Modsætningen mellem
% --- p. 30 ---
\setcounter{footnote}{0}%
\opage{30}Spiritualisme og Materialisme, Identitetshypotesen. Mod Descartes
indvender han, at fordi to Egenskaber ikke kunne udledes
den ene af den anden, er det ikke berettiget at henføre
dem til to forskellige Væsener (Substanser); de kunne være
irreduktible Egenskaber ved samme Væsen. Fra Hobbes adskiller
han sig ved den klare Indsigt, at det Aandelige ikke
kan udledes af det Materielle, og at den naturvidenskabelige
Opfattelse af Tilværelsen derfor ikke kan være udtømmende.
--- I Descartes, Hobbes og Spinoza have vi Hovedrepræsentanter
for de tre Hovedhypoteser om Forholdet mellem Aand
og Materie. ---

Af Spinozas specielle Psykologi er især hans Fremstilling
af Følelser og Drifter af blivende Betydning. Ligesom Hobbes,
og under Paavirkning af ham, gaar han ud fra Selvopholdelsestrangen,
der i hans System viser sig i Sammenhæng med
hele Verdensanskuelsen, idet det jo er den uendelige Substans,
der virker i den enkelte modus, hvis Stræben efter
Selvopholdelse altsaa er en Del af den guddommelige Virken.
Naar denne Stræben tilfredsstilles, opstaar der Lyst; naar den
hæmmes, opstaar der Ulyst. Dog sker dette kun ved Overgangen
til Tilfredsstillelse eller Hæmning; en stadig og uforandret
Tilstand vilde ikke være forbunden med Lyst eller
Ulyst. De specielle Arter af Følelse udvikle sig under Indflydelse
af Forestillingers Association, hvorved en Følelse ved
en Genstand efterhaanden kan blive knyttet til helt andre
Genstande. Medfølelse kan f. Ex. opstaa, naar vi forestille os
et Væsen, der ligner os, føle Lyst eller Ulyst.

Paa Selvopholdelsestrangen bygger Spinoza, ligesom Hobbes,
sin Etik. Det er Modsætningen mellem Aktivitet og Passivitet,
mellem at være Aarsag og at være Virkning, der her
bliver afgørende. Aktive ere vi, naar Aarsagen til vor Tilstand
og Handlinger ligge i os selv, passive, naar den ligger udenfor
os. Som Enkeltvæsen (modus) er Mennesket fra først af
overvejende bestemt ved ydre Forhold, ved Forholdet til andre
modi; dets Følelser og Drifter fremkaldes udefra og gøre
det afhængigt. Men i Kraft af Selvopholdelsestrangen kan der
vækkes en Stræben efter saavidt som muligt at have Aarsagen
% --- p. 31 ---
\setcounter{footnote}{0}%
\opage{31}til sine Følelser og Drifter i sig selv. Der danner sig da et
Ideal af Menneskeliv (idea hominis, tanqvam naturæ humanæ
exemplar), saaledes som det vilde være, naar det udfoldede
sig i fuld Frihed og Uafhængighed. Dette Maal afgiver en
Maalestok for, hvad der kaldes godt og ondt. Hvad der fører
os nærmere til Idealet, er godt; hvad der fjerner os fra det,
ondt. Disse Prædikater kunne kun anvendes paa endelige Væsener,
der staa midt i Kampen med ydre Forhold. --- Da en
Drift kun kan overvindes ved en anden Drift, maa Idealet selv
vække en Drift, der voxer sig stærk, saa den kan komme til
at herske. Gennem Livet i Samfundet styrkes denne Drift.
Naar Menneskene forene deres Kræfter, kunne de lettere
overvinde ydre Hindringer og sørge for deres fysiske Selvopholdelse,
saa at de aandelige Goder, især sand Erkendelse,
kunne komme til at optage Sindet. Selve de aandelige Goder
kunne være Genstand for fælles Stræben; de vække ikke den
Strid som materielle Goder, der hvert for sig kun kunne blive
Enkelte til Del; tvertimod, kan paa det aandelige Omraade
den Ene hjælpe den Anden uden selv at miste Noget. Naar
den fulde Kraft og Frihed er naaet, den stærke Selvhævdelse
(fortitudo), hvori Dyden bestaar, saa vil den Enkelte ikke blot
med Livsmod (animositas) gøre sin egen Personlighed gældende,
men ogsaa have Kraft nok til med Højmod (generositas) at
yde Andre aandelig og materiel Hjælp. Den højeste Grad af
Frihed og Selvhævdelse naas dog kun gennem den fulde Forstaaelse
af os selv i vor Enhed med Tilværelsens Substans,
der rører sig i os saavel som i den ydre Natur. Vi faa da en
scientia intuitiva af os selv, se os selv sub specie æterni, og
vi fyldes af en Erkendelsesglæde, der afføder Kærlighed til
Guddommen (amor intellectualis dei).

I sin Statslære hævder Spinoza ligesom Hobbes en skarp
Modsætning mellem Naturtilstanden og Statslivet. Men paa
Grund af den Betydning, Samfundslivet ifølge ham har for
den Enkeltes Udvikling, er det saa langt fra, at den Enkelte,
som hos Hobbes, opgiver sin Frihed ved at træde ind i
Staten, at den fulde Frihed, særligt forsaavidt den betinges ved
aandelig Udvikling, først bliver mulig i Staten. Staten skal
% --- p. 32 ---
\setcounter{footnote}{0}%
\opage{32}ikke gøre Menneskene til Dyr eller Maskiner, men skaffe
Betingelser for Udviklingen af Menneskenes aandelige og legemlige
Selvvirksomhed. Den vilde da handle mod sit eget
Maal, naar den ikke hævdede Tænke-, Tale- og Religionsfrihed.
--- Det er i sin \textit{Tractatus theologico-politicus} (1670), at
Spinoza tager Ordet for Religionsfrihed. I dette Værk hævder
han tillige Berettigelsen og Betydningen af en rent historisk
% PRINTED AS IS (read at the image, odd spots o6)
Metode ved Studiet af de bibelske Skrifter

d. \textsc{Gottfried Wilhelm Leibniz} (1646---1716) er ligesom
sine tre Forgængere, Descartes, Hobbes og Spinoza, gennemtrængt
af den mekaniske Naturopfattelses Betydning. Som
Dreng havde han i sin Faders, en Leipzigerprofessors Bibliotek
studeret middelalderlige Værker, og det var først efter en
Kamp, han sluttede sig til den nye Naturvidenskab og Filosofi.
Han følte sig da som „hensat i en anden Verden“. Han saa,
at de nye Tanker ikke kunde afvises; men han kunde dog
ikke opgive den Overbevisning, at Naturen tilsidst lededes af
et Forsyn, og at den mekaniske Lovmæssighed tilsidst var
teologisk begrundet, d. v. s. havde sin sidste Grund i et guddommeligt
Formaal. Under et Ophold i Paris havde Fysikeren
Huyghens stor Indflydelse paa hans matematiske og fysiske
Opfattelse og senere gjorde han personligt Bekendtskab med
Spinoza. Fra 1676 af levede han i Hannover som kurfyrstelig
Hofraad og Bibliotekar, rastløst beskæftiget med filosofiske,
matematiske, historiske og juridiske Arbejder. Hans universelle
Geni gjorde ham det muligt at gribe skabende og ordnende
ind paa meget forskellige Omraader. Overalt lededes han af
én Grundtanke, Kontinuiteten, den ubrudte Sammenhæng
gennem en Mangfoldighed af Overgange --- i Tankens, i Tallets,
i Naturens og i Formaalenes Verden.

Naar Descartes og Spinoza lærte, at Bevægelsens Sum i
Verden stedse er den samme, fandt Leibniz heri den Vanskelighed,
at det var uforklarligt, hvorledes Bevægelse og
Hvile kunde afløse hinanden paa de enkelte Steder i Verden:
de staa jo som modsatte Tilstande! Kontinuiteten kan kun
hævdes ved Begrebet Kraft, d. v. s. det i en Tilstand, der
% --- p. 33 ---
\setcounter{footnote}{0}%
\opage{33}indeholder Grunden til en følgende Tilstand. Det er Kraften,
ikke Bevægelsen, der bestaar. Naar Bevægelsen hører op paa
et Sted i Verden, afløses den ikke af absolut Hvile, men af
en Ligevægtstilstand, i hvilken modsatte Kræfter holde hinanden
i Ligevægt. Overgangen til Bevægelse sker da ved, at
Ligevægten ophæves. Leibniz har her lagt Grunden til den
moderne Forskel mellem potentiel og aktuel Energi, ledet af
sin store Interesse for Kontinuitet, der igen hænger sammen
med den Forudsætning, at Alt har sin tilstrækkelige Grund.
Naar man igen vilde spørge, hvorfor da Kraften bestaar, svarer
Leibniz, at det Modsatte, en Forsvinden af Kraften, vilde
være en Ufuldkommenhed og en Tilbagegang, der strider mod
Forsynet. Han fordrer den mekaniske Naturforklaring gennemført
i det Enkelte; men denne sidste Forudsætning, Kraftens
Bestaaen, er teleologisk begrundet.

Ogsaa paa et andet Punkt afviger Leibniz fra sine Forgængere.
Descartes og Spinoza gik ud fra Udstrækningen som
en absolut Egenskab (skabt af Gud eller guddommeligt Attribut).
Men det Udstrakte er stedse sammensat, og det egentlige
Reale maa være de Elementer, hvoraf det er sammensat, og
som derfor (da de skulle ligge til Grund for Udstrækningen)
ikke selv kunne være udstrakte, men maa være Enheder
(Monader). Vil man danne sig en Forestilling om disse Monader,
staar der ingen anden Vej aaben end Analogien med
vore egne Sjæle. Enhver Monade er en lille Verden, indenfor
hvilken de følgende Tilstande have deres Grund i de foregaaende,
ligesom i den menneskelige Sjæl sjælelige Tilstande
udvikle sig af hverandre gennem Stræben efter Klarhed.
Leibniz stiller sig overfor Modsætningen mellem Ubevidsthed
og Bevidsthed ligesom overfor Modsætningen mellem Hvile
og Bevægelse. Han fastholder Kontinuiteten i Tilværelsen fra
dens aandelige Side set, saavel som fra dens materielle Side
set. Intet i Verden begynder med absolut Pludselighed, ved
et Spring, Bevidstheden saalidt som Bevægelsen. Ubevidsthed
er kun en lavere Grad af Sjæleliv.

I sin Opfattelse af Forholdet mellem det Aandelige og det
Materielle staar Leibniz paa Spinozas Standpunkt, forsaavidt
% --- p. 34 ---
\setcounter{footnote}{0}%
\opage{34}som han hylder Identitetshypotesen. Men medens for Spinoza
Aand og Materie stode som sideordnede Attributer, er det
ifølge Leibniz den aandelige Side, der er den fundamentale;
den udtrykker Tilværelsens inderste Væsen, medens Materien
er dens ydre, sanselige Form. Materielle Genstande maa vi
tyde som sanselige Fremtrædelser af Noget, der danner en
Analogi til vort eget Sjæleliv. --- Det er en betydningsfuld
Opdagelse, Leibniz har gjort, at al Metafysik, d. v. s. al afsluttende
Verdensanskuelse, grunder sig paa Analogi. Alle
Mytologier, Religioner og metafysiske Systemer vare gaaede
denne Vej, men Leibniz er den Første, der er bleven sig det
bevidst. Hans eget System er det første originale Forsøg paa
metafysisk Idealisme (d. v. s. den Opfattelse, at det Aandelige
udgør Tilværelsens inderste Kærne) siden Platon, og det staar
som Forbillede for alle senere Forsøg i lignende Retning. Naar
han selv har saa stor Tillid til Analogien, er det, fordi den
anviser den eneste Vej til afsluttende Forstaaelse af Verden,
og han gaar, i Kraft af den tilstrækkelige Grunds Princip, ud
fra, at Verden maa være forstaaelig. (\textit{Système nouveau} 1695.
\textit{Monadologie} 1714).

Kontinuitetens Lov fører Leibniz til overalt i Naturen og
i Sjælelivet at lægge Vægt paa de smaa Overgange, de forsvindende
Differenser i Modsætning til de Spring, Erfaringen
synes at vise. Den samme Tanke har ledet ham som Matematiker
og førte ham her til Opfindelsen af Differentialregningen.
Hans „Differentialer“ ere forsvindende smaa Størrelser
eller Størrelsesændringer, der dog ved Summering (Integration)
give en endelig Størrelse. Differentialerne udtrykke
kontinuerlige Variationer. --- Paa mange forskellige Omraader
har Leibniz' Tænkning arbejdet; men dens Type er overalt
den samme.

Ogsaa mellem de uendeligt mange Monader indbyrdes er
der Kontinuitet. De ere alle individuelt forskellige. Der findes,
paastaar Leibniz, ikke to Ting i Verden, der ere absolut ens.
Der vilde jo saa (tilføjer han teleologisk) ikke være nogen
Grund til, at de skulde existere. I Tilværelsen er den størst
mulige Mangfoldighed virkeliggjort sammen med den størst
% --- p. 35 ---
\setcounter{footnote}{0}%
\opage{35}mulige Kontinuitet. Og den sidste Grund til denne store
Harmoni finder Leibniz (idet han gaar over fra Filosofien til
Teologien) deri, at alle Monader ere Udstraalinger („Fulgurationer“)
af en guddommelig Urmonade, der blandt alle Verdensmuligheder
har valgt den, der vilde udvise den størst
mulige Harmoni. Han paatager sig endog at bevise (i sin
\textit{Théodicée} 1710), at denne Verden er den bedste af alle mulige
Verdener. Enhver endelig Verden maa frembyde Ufuldkommenheder,
dels Begrænsninger (metafysiske Onder), dels
Lidelser (fysiske Onder), dels Synder (moralske Onder). Ved
Guds Valg ere disse Ufuldkommenheder indskrænkede til det
mindst Mulige. Men ud over Mulighedernes Omraade kan selv
den guddommelige Villie ikke gaa. Og naar vi finde saa meget
Ondt i vor Erfaring, maa vi ikke glemme, at den Del af Verden,
som vi kende, og som maaske er den værste Del, kan
være forsvindende i Forhold til Uendeligheden, --- og desuden,
at det Onde ikke ophæver Verdensharmonien, men netop
tjener til at fremhæve denne, ligesom en mørk Skygge forøger
Virkningen af de lyse Farver i et Maleri, eller som en Klang,
der for sig alene vilde lyde ilde, kan passe udmærket ind i
et helt Musikstykke! ---

Foruden de anførte Værker af Leibniz (kun faa af den
store Mængde Optegnelser, han har efterladt sig) maa mærkes
\textit{Nouveaux Essais,} der først udkom længe efter hans Død,
og i hvilken han imødegik John Lockes Erkendelseslære.

\closerule

\chapterhead{III.}{ERFARINGSFILOSOFIEN.}

Systemdannelserne i det 17. Aarhundrede havde til Grundlag
Begreber og Forudsætninger, der ikke bleve nærmere
undersøgte. Ganske særligt gælder dette om Aarsagssætningen
% --- p. 36 ---
\setcounter{footnote}{0}%
\opage{36}og nogle naturvidenskabelige Sætninger, der ligesom den bleve
ansete for selvfølgelige. Det er den engelske Filosofis Fortjeneste
at have undersøgt saadanne Begrebers og Forudsætningers
psykologiske Oprindelse og erkendelsesteoretiske Gyldighed.
Denne Undersøgelse fik Betydning ud over det rent
filosofiske Omraade, idet man nu ogsaa prøvede de Principer,
der gjaldt i Staten, i Religionen og i Opdragelsen. Autoriteter,
der hidtil havde vundet uvilkaarlig Tilslutning, skulde nu aflægge
Regnskab for deres Herkomst og deres Berettigelse.
Sammen med Psykologien og Erkendelsesteorien kom da nu
Etiken i Forgrunden.

a. \textsc{John Locke} (1632---1704) skal have faaet Ideen til sit
banebrydende Værk \textit{Essay on human understanding} (1690)
ved en Diskussion om religiøse og moralske Emner; da det
viste sig vanskeligt at opnaa Enighed, indsaa Locke, at man
først maatte undersøge vor Erkendelsesevne for at se, hvad
den formaar og hvad den ikke formaar. Han opstiller det
Princip, at alle vore Forestillinger skyldes Erfaring, Iagttagelse
enten af ydre Ting eller af vore egne indre Tilstande og
Virksomheder. Dette hvad Oprindelsen angaar; hvad Gyldigheden
angaar, finder han det selvfølgeligt, at hvad vi faa fra
Iagttagelse, altsaa overvejende passivt, maa være gyldigt, hvorimod
der kommer en Vanskelighed ved Begreber og Sætninger,
der skyldes vor egen Virksomhed.

Locke var ligesaalidt som Descartes og Hobbes tilfreds med
sin middelalderlige Undervisning. De to nævnte Mænds Skrifter
førte ham ind i Filosofien. Da han ikke kunde underskrive
de 39 Artikler, maatte han opgive sin oprindelige Plan at
blive Præst og studerede derefter Medicin. Senere traadte han
i Jarlen af Shaftesburys Tjeneste og vedblev at være knyttet
til hans Familie, som Lærer, Sekretær og Ven. Da Jarlen blev
styrtet, drog Locke over til Holland, og han har her forfattet
sine vigtigste Skrifter. Han vendte tilbage til England med
Vilhelm af Oranien og fik en ikke ringe Indflydelse under
den nye Regering. Sine sidste Aar tilbragte han i landlig
Ensomhed.

% --- p. 37 ---
\setcounter{footnote}{0}%
\opage{37}Da alle Forestillinger tilsidst stamme fra Erfaring, kan
der ikke være Tale om „medfødte“ Forestillinger. Antagelsen
af saadanne skyldes den menneskelige Ladhed, der skyr den
Møje, der kan være forbunden med at udfinde vore Forestillingers
Oprindelse. Men af de simple Forestillinger vi optage
overvejende (for the most part) passivt, danne vi selv nye
Forestillinger. Vi \textit{forbinde} simple Forestillinger til sammensatte;
saaledes danne vi Forestillinger om Ting (Væsener,
Substanser) ved Forbindelse af de forskellige Egenskaber, der
hver for sig svare til Fornemmelse. Vi \textit{sammenstille} forskellige
Forestillinger og faa derved nye Forestillinger, f. Ex. om
Lighed og Ulighed, om Aarsagsforhold. Og vi \textit{udskille} (abstrahere)
Forestillinger af deres oprindelige Sammenhæng, f.
Ex. naar vi danne Forestillingen om en Farve i al Almindelighed,
eller om Rummet afset fra ethvert Indhold. Hvorledes
kunne nu saadanne af os selv frembragte Forestillinger være
gyldige for selve Tilværelsen? De ere jo ikke Kopier eller
Aftryk! Gyldigheden kan kun bestaa i, at vi kunne bruge dem
som Mønstre eller Forbilleder, efter hvilke vi maale det Iagttagne.
Paa den Maade bruge vi i Matematiken de geometriske
Figurer og i Moralfilosofien de etiske Idealer. Men nu Begrebet
Ting (Væsen, Substans), der jo udtrykkeligt skal angaa
Noget, der existerer i sig selv? Hvad vi opfatte, er jo aldrig
Tingene selv, men kun deres Egenskaber! Det nytter ikke at
sige, at der maa ligge „Noget“ til Grund for Egenskaberne;
ti saa maa der jo ligge „Noget“ til Grund for dette „Noget“,
og saa fremdeles. Alligevel vil Locke ikke nægte Existensen
af Ting eller Væsener. --- Karakteristisk er ogsaa hans Stilling
overfor Aarsagssætningen. Han anser den for selvindlysende,
ligesom Systematikerne: Intet kan ikke frembringe Noget! Og
han benytter den trygt til at bevise Guds Tilværelse (da Verden
maa have en Aarsag). Aarsagsforestillingen derimod svarer,
mener han, til den Iagttagelse, at Ting og Egenskaber opstaa
ved andre Tings Virken. Men han har jo selv regnet Aarsagsforestillingen
til dem, vi selv frembringe (ved Sammenstilling)!
--- Locke efterlod da Problemer nok til sine Efterfølgere.

% --- p. 38 ---
\setcounter{footnote}{0}%
\opage{38}b. \textsc{George} \textsc{Berkeley} (1685---1753) ønskede at fjerne de
for Religionen farlige Konsekvenser af den nye Videnskab,
og dette Maal mente han at naa ved Kritik af de abstrakte
Begreber (især Begreberne Rum og Materie) og ved at gaa
tilbage til den umiddelbare Erfaring og Anskuelse. Sjældent
have barnlig Fromhed og skarp kritisk Analyse været forbundne
paa saa nøje en Maade som hos denne klare og fine Aand.
Paa Universitetet i Dublin havde han studeret Locke, Newton
og Kemikeren Boyle, og han forfattede sine Hovedskrifter
som ung Mand. Derpaa blev han anglikansk Gejstlig og polemiserede
ivrigt mod Fritænkerne. Hans Yndlingsplan var at
stifte en Højskole i Amerika, dels for at omvende Indianerne,
dels for at forplante Kunst og Videnskab til ny Blomstring i
den nye Verden. Efter tre Aars Ophold derovre maatte han
opgive denne Plan. Han virkede derefter som Biskop i Cloyne
i Irland, lige ivrig som Sjælesørger, som Filantrop og som
Patriot.

I sit Hovedskrift \textit{(Principles of Knowledge.} 1710) viser
Berkeley, at vi i Ordets strenge Betydning ingen Almenforestillinger
have, saaledes som endnu Locke havde antaget. Vi
kunne forestille os en Del eller en Egenskab sondret fra andre
Dele og Egenskaber. Men vi kunne ikke forestille os
noget for flere Egenskaber „Fælles“, som f. Ex. Farve overhovedet.
En Forestilling om Farve maa enten være en Forestilling
om Rødt eller om Gult eller om Grønt o. s. v. Naar
jeg vil have en Forestilling, der skal gælde for flere kvalitativt
forskellige Ting, maa jeg enten betragte en enkelt af disse
Ting som Exempel, eller bruge et Ord, der passer paa enhver
af dem. Forestillingen om Materie som noget Udstrakt
og Bevægeligt er nu en blot Almenforestilling; ti vi kunne
ikke iagttage noget Udstrakt uden Farve, Haardhed o. s. v.
Og Rummet (Udstrækningen) kende vi (som Berkeley havde
vist i \textit{Theory of Vision.} 1709) fra to forskellige Sanser, Syn
og Berøringssans. Rummet i Almindelighed (det matematiske
Rum) skal være „fælles“ for Syn og Berøring; men noget
Saadant kan ikke forestilles. Vor Rumsforestilling maa enten
angaa Noget, der kan ses, eller Noget, der kan berøres.

% --- p. 39 ---
\setcounter{footnote}{0}%
\opage{39}Berkeley vil kun bekæmpe den Antagelse, at der bag
Sansekvaliteterne skulde ligge en materiel Substans. Han vil
ikke bekæmpe selve Naturvidenskaben, men ser dennes Opgave
i at udfinde, efter hvilke Regler vi, naar en Sansekvalitet
er fremtraadt, kunne vente en anden. At forklare Naturen vil
sige at paavise de lovmæssige Forhold mellem Sansekvaliteterne.
--- Som Teolog erklærer Berkeley det for en Omvej
at lade Gud først skabe „Materien“ og derpaa at lade denne
frembringe Sansefornemmelserne hos os. Simplest er det at
tænke sig Gud (hvis Virken vi forestille os i Analogi med
vor Villen, der udgør vor Sjæls inderste Væsen) umiddelbart
virkende i Frembringelsen af vore Sansefornemmelser efter
bestemte Love.

c. \textsc{David Hume} (1711---1776) bragte den erfaringsfilosofiske
Undersøgelse af den menneskelige Erkendelse til en foreløbig
Afslutning, især ved sin Undersøgelse af Aarsagsbegrebet, dette
Begreb, hvis Gyldighed er en stiltiende Forudsætning for
naturvidenskabelig, psykologisk og historisk Forsken saavel
som for de religiøse Trosforestillinger og for det praktiske
Menneskeliv. Hume's Problemer ligge derfor ved Roden af
al menneskelig Tænken. Og han er en Problemstiller af første
Rang, foruden at hans psykologiske Talent gør ham det
muligt at kaste Lys over de menneskelige Tankers faktiske
Udvikling selv der, hvor deres objektive Gyldighed er problematisk.

Hume var Søn af en Godsejer i det sydlige Skotland og
viste tidligt Lyst og Anlæg til at tænke og lære. Efter forgæves
at have søgt at bane sig en Vej i det praktiske Liv
trak han sig tilbage i Ensomhed og skrev under et Ophold i
Frankrig sit Hovedværk \textit{Treatise on Human Nature} (1739
---40). Senere gik han over til historiske og økonomiske
Studier og forfattede en Englandshistorie, et af de første Værker,
i hvilke der tages Hensyn til alle Sider af Kulturudviklingen.
I hans \textit{Essays} (1748 ff.) findes foruden betydningsfulde
økonomiske Afhandlinger to filosofiske Afhandlinger (om den
menneskelige Forstand og om Moralens Principer), i hvilke
% --- p. 40 ---
\setcounter{footnote}{0}%
\opage{40}han giver kortere Fremstillinger af Hovedværkets vigtigste
Problemer. I \textit{Natural History of Religion} (1757) behandlede
han det religiøse Problem fra psykologisk og historisk Synspunkt;
i de 1751 forfattede, men først efter hans Død udgivne
\textit{Dialogues on Natural Religion} behandlede han det samme
Problem fra den kritiske Side. --- Efterat have beklædt forskellige
offentlige Embedsstillinger henlevede han sine sidste
Aar i Edinburgh i videnskabeligt Otium.

Medens Locke havde skelnet mellem Forestillingernes Oprindelse
og deres Gyldighed, falde for Hume begge Dele
sammen. En Forestilling (idea) er kun gyldig, naar den stammer
fra en Fornemmelse (impression). Hvorfra Fornemmelserne
igen komme, undersøger han ikke, fordi det er erkendelsesteoretisk
ligegyldigt; med Hensyn til Gyldigheden kommer
det ud paa Et, om de stamme fra os selv, eller fra ydre
Ting, eller fra Gud. Ogsaa den rene Logik og Matematik, der
kun angaa Forhold mellem vore Forestillinger, ikke Forestillingerne
selv, udelukker Hume fra sin Undersøgelse. Han vil
kun drøfte Gyldigheden af de Forestillinger, ved hvis Hjælp
vi mene at kunne gaa ud over hvad der er givet i vore
Fornemmelser. En Forestilling af denne Art er (som Berkeley
allerede havde paavist for det Materielles Vedkommende)
Forestillingen om Ting (Væsener eller Substanser). Vi fornemme
kun enkelte Egenskaber, ikke det, der „har“ dem. I
os selv iagttage vi skiftende Fornemmelser, Forestillinger og
Følelser, men ikke nogen Sjæl eller noget Jeg, ingen konstante
Elementer, der kunne fortjene Navn af „os selv“. Og
paa tilsvarende Maade gaar det med Aarsagsbegrebet. Vi iagttage
enkelte Fænomener, der følge efter hverandre i Tiden,
men et indre Baand, en nødvendig Sammenhæng iagttage vi
ikke. Kausalitet kan ikke iagttages. Og alle Fænomener, alle
Ting optræde for vor Iagttagelse som selvstændige, hver med
sin Beskaffenhed; vi kunne meget vel forestille os dem hver
for sig, uden at være nødte til at gaa over fra den ene til
den anden. Naar en Kugle bevæger sig, kunne vi forestille os
denne Bevægelse uden at nødes til at forestille os nogen anden
Bevægelse.

% --- p. 41 ---
\setcounter{footnote}{0}%
\opage{41}Men vi bruge jo dog Begreberne Ting og Aarsag! Hvoraf
kan det komme? --- Hume svarer: Begreberne Ting eller
Substans danne vi under Indflydelse af den Tilbøjelighed, vor
Bevidsthed har til at fortsætte med en Virksomhed, selv efter
at Anledningen er ophørt. Det er derved, at vi danne ideale
Forestillinger: gennem vor Indbildningskraft forestille vi som
fuldkomment, hvad der for Iagttagelsen kun er ufuldkomment.
Saaledes danne vi f. Ex. Forestillinger om absolut regelmæssige
Figurer. Og saaledes danne vi Forestillinger om konstante
Ting eller Substanser, skønt vi kun kende forskellige Grader
af Konstans. For Aarsagsbegrebets Vedkommende virker her
Vanen med. Vi ere vante til, at der ikke sker Noget, uden at
der forud er sket noget Andet, og bagefter sker noget Tredie.
Derfor vente vi, naar Noget sker, dels en „Aarsag“, dels en
„Virkning“. Men deri ligger intet Bevis. Det er en blot Ideassociation.
Det gaar med „Nødvendigheden“, som det gaar
med Sansekvaliteterne; den er ligesom disse subjektiv, ikke
objektiv eller real. Vanen kan begrunde en Tro, en instinktiv
Forvisning, men kan ikke afføde Bevis. Hume mener naturligvis
ikke, at vi kunne undlade at bruge Aarsagsbegrebet;
men det var for ham det Vidunderlige, at Gyldigheden af det
Begreb, hvoraf al vor reale Videnskab afhænger, ikke selv
kunde begrundes.

Ligesom Hume i sin Erkendelsesteori ender med at finde
det sidste Grundlag i en Trang eller et Instinkt, hvis Berettigelse
ikke kan bevises, saaledes grunder han ogsaa sin
Etik paa Følelse, ikke paa Fornuft. Fornuften kan kun sammenligne
og analysere; den kan overtyde os om Kendsgerninger
og Forhold, men ikke om Værdi. Godt og Ondt
ere Kvaliteter, vor Følelse driver os til at tillægge Tingene,
ligesom vore Sansefornemmelser fører os til at tillægge
dem Sansekvaliteterne. Da vi nu kalde Handlinger og Egenskaber
gode, selvom de ikke komme os selv til Gode, kan
den Følelse, der ligger til Grund for Billigelsen, ikke være af
egoistisk Art. Vi indtage i vor Billigelse eller Misbilligelse et
fælles, ikke et privat Standpunkt. Naar vi f. Ex. betragte Retfærdigheden
som en god Egenskab, kan dette kun udspringe
% --- p. 42 ---
\setcounter{footnote}{0}%
\opage{42}af, at vi omfatte Menneskelivet, hvorsomhelst det rører sig,
med Sympati. Maaske var det fra først af kun Interessen for
vor egen Sikkerhed, der lod os billige Retfærdigheden; der
maa da være foregaaet en Ændring i vort Følelsesliv, naar vi
kalde den god, ogsaa hvor vor egen Velfærd ikke berøres af
dens Overholdelse. Grundmotivet for etisk Vurdering er da
Sympatien eller Kammeratfølelsen (fellow-feeling).

Ligesom i Etiken, saaledes opponerer Hume ogsaa i Religionsfilosofien
mod Overvurdering af det Intellektuelles Indflydelse.
Og ligesom i Erkendelsesteorien anstiller han ogsaa
her en dobbelt Undersøgelse, idet han dels drøfter Religionens
psykologiske Udspring, dels Gyldigheden af de religiøse Forestillinger.
Den første fremstilles i „Religionens Naturhistorie“,
den anden i „Dialogerne om naturlig Religion“. --- Religion
opstaar ikke af rent intellektuelle Motiver, men af Frygt og
Haab og af Tilbøjeligheden til at forestille sig alle Væsener i
Lighed med Mennesket. I Begyndelsen forestiller Mennesket
sig de Væsener, til hvilke det tager sin Tilflugt under Livets
Angst, i meget ufuldkomne Former. Men ogsaa her virker
Indbildningskraftens Tendens til at gaa ud over det Begrænsede
og Endelige, og det gaar op for Mennesket, at dets Gud
maa være et uendeligt Væsen, og at der kun kan være én
Gud. Ved Siden af denne Tendens, der fører til at forestille
sig Guddommen som hævet over alt Menneskeligt og fjernet
fra den endelige Verden, gives der dog en modsat Tendens,
en Trang til at forestille sig den som nær og paa anskuelig
Maade, i menneskelig Skikkelse og med menneskelige Egenskaber,
og vi finde i Religionen en stadig Svingen mellem
disse to Tendenser. Derved forstaa vi de Modsætninger og
Modsigelser, der findes dels mellem de forskellige Religioner,
dels indenfor samme Religion. --- I sine „Dialoger“ indfører
Hume en Supranaturalist (der hævder Troen paa et Overnaturligt),
en Rationalist (der mener ad Fornuftens Vej at
kunne bevise Guds Tilværelse) og en Naturalist (der hævder,
at vi hverken i Tro eller Fornuft kunne komme ud over
Erfaringen). Skønt Naturalisten tilsidst gør et høfligt Tilbagetog,
er det dog tydeligt, at hans Grunde ere afgørende i Hume's
% --- p. 43 ---
\setcounter{footnote}{0}%
\opage{43}egne Øjne. Han bestrider Retten til at slutte til en Gud fra
Naturens Orden og Hensigtsmæssighed. Hvorfor skulde det
Hensigtsmæssige i Verden ikke kunne opstaa af naturlige
Aarsager? Gennem en fremskridende Udvikling kunde de
ordnede Helheder og Systemer i Verden tænkes at have
dannet sig. Og i hvert Tilfælde frembyder den virkelige Verden
saa meget Ufuldkomment, at man ikke kan slutte til et
% PRINTED AS IS (read at the image, odd spots o6)
Væsen, der paa én Gang skulde være algodt og almægtigt
Endeligt --- selvom man vil finde Verdens Udspring i en
guddommelig Tanke, maa man igen spørge, hvorfra denne
Tanke har sit Udspring. --- Naar Naturalisten dog ender med
at finde Vanskeligheder i alle Standpunkter, maa vi vistnok
ikke deri se blot Forsigtighed, men et Udtryk for Hume's
stadige Stræben efter at stille Problemerne skarpt og at holde
dem aabne.

\closerule

\chapterhead{IV.}{OPLYSNINGSFILOSOFIEN.}

Systemerne og Erfaringsfilosofien virkede kun paa en
snever Kreds. Men ved Midten af det attende Aarhundrede
vaagnede der en Stræben efter at udbrede de stille Tænkeres
Ideer i større Kredse. I England antog denne Bevægelse, som
man plejer at kalde Oplysningen, ikke saa karakteristiske
Former som i Frankrig og Tyskland. I Frankrig, hvor den
gamle Tingenes Orden i Kirke og Stat stod mere afsluttet og
afvisende overfor al Forandring end i England og Tyskland,
antog Oplysningsbevægelsen en afgjort revolutionær Karakter.
Man benyttede Lockes Fordring om Paavisning af de Erfaringer,
fra hvilke Forestillingerne stamme, til at kræve den
bestaaende Orden til Regnskab, og man mente selv at være
i Besiddelse af Forudsætninger, der vare tilstrækkelige til
% --- p. 44 ---
\setcounter{footnote}{0}%
\opage{44}Løsning af Tankens og Livets Problemer, saa at hvad der
ikke stemmede med dem, maatte omstyrtes. I Tyskland forenede
Eftervirkningen af Leibniz' Filosofi sig med Lockes
Indflydelse til at frembringe en Populærfilosofi, der foretog
en Udrensning af de overleverede Forestillinger. Men baade
i Frankrig og i Tyskland optraadte Mænd, der hverken tilfredsstilledes
ved den nye Oplysning eller ved den gamle
Overlevering, fordi deres Blik var rettet mod dybere liggende
Livsbetingelser. Saadanne Mænd vare Rousseau i Frankrig og
Lessing i Tyskland.

1. Ved et Besøg i England lærte \textsc{Voltaire} (1694---1778)
engelske Forhold og engelske Tanker at kende. Af hans
\textit{Lettres sur les Anglais} (1733, brændt 1734) lærte Frankrig
Locke's og Newton's Tanker, saavel som engelske religiøse,
literære og politiske Forhold at kende. Voltaires digteriske
Virksomhed og hans utrættelige Arbejde for Tolerance og
Retfærdighed tilhører Literatur- og Kulturhistorien. Her skal
kun anføres hans vigtigste filosofiske Værker: \textit{Dictionnaire
philosophique portatif} (1764, brændt næste Aar) og \textit{Le philosophe
ignorant} (1766). Snart med Grunde, snart med Haan
og Spot, snart med dyb Indignation bekæmpede han Fordomme
og aandelig og legemlig Undertrykkelse. Sine Grunde
hentede han fra Locke's Lære om alle Forestillingers Udspring
af Erfaringen og fra Newton's Lære om den lovmæssige
Orden i Naturen, og kritiserede ud fra dem Kirkens
Lære. Hans eget religiøse Standpunkt var den saakaldte naturlige
Religions (smlg. Herbert af Cherbury). Han ansaa den
Religion for den bedste, der indeholder megen Moral og faa
Dogmer. Fornuften kan ud fra Hensigtsmæssigheden i Naturen
bevise Guds Tilværelse, men ender iøvrigt i uløselige Problemer;
vi vide hverken, hvad Materien er, eller hvad Sjælen
er. Gud har givet os Fornuften til at lede os i vor Handlen,
% PRINTED AS IS (read at the image, odd spots o6)
ikke til at trænge ind i Tingenes Væsen Hvad der ligger ud
over den naturlige Religion, skyldes Dumhed eller Bedrag;
Dumheden gør Begyndelsen, og Bedraget tager den saa i sin
Tjeneste. Dannede Folk behøver ingen anden Religion end
den naturlige; de Andre maa helst blive ved at gaa i Messe.
% --- p. 45 ---
\setcounter{footnote}{0}%
\opage{45}--- Skønt Hensigtsmæssigheden i Naturen for Voltaire beviser
en Guds Tilværelse, er der dog --- mener han --- for meget
Ondt og for megen Lidelse i Verden, til at denne Gud kan
være almægtig. I sit Digt om Lissabons Undergang og i Romanen
„Candide“ har Voltaire udtalt sine Tanker om det
Onde i Verden. Man kan have Haab om Fremgang, men
Verden er langt fra at være „den bedste Verden“. Voltaire
haaner Leibniz' Optimisme.

I afgjort Modsætning til Voltaire fremstillede \textit{Holbach}
(1723---1789), en tysk, i Paris levende Baron, Materialismen
som Naturvidenskabens sande Konsekvens. \textit{(Système de la
nature.} 1770). Bevægelsen er en oprindelig Egenskab ved
Materien, og materielle Fænomener maa kun forklares ved
materielle Aarsager. Det er da ikke nødvendigt at antage
aandelige, fra Materien forskellige Væsener. Tænken eller
Bevidsthed er kun Bevægelse i Materiens Smaadele. Der gives
kun én Videnskab, Fysiken, Læren om Bevægelsen. Antagelsen
af en dobbelt Art Væsener, aandelige og materielle, er
dog ikke blot unyttig, den er skadelig. Den begunstiger Overtroen
og derved Præsternes Herredømme. Allerede den saakaldte
naturlige Religion er farlig; ti Religion kan ikke bestaa
uden Kultus, og dermed er strax Præsternes Herredømme
givet. Gudsforestillingen skyldes snedige Præster (fabricateurs
de la divinité), der ved dens Hjælp vilde herske over Menneskene!
---

Trods al Modsætning vare Voltaire og Holbach enige om,
at Oplysningens Tid var kommen, og at det nu kun gjaldt
om at udbrede den saa langt som muligt. I Diderots og
d'Alemberts store Encyklopædi, ved hvilken Tidens mest fremragende
Mænd vare Medarbejdere, fandt de dristigste Tanker
og den skarpeste Kritik Indpas under Forklædninger og Omskrivninger.
Efter den have de franske Oplysningsfilosofer
faaet Navn af Encyklopædister. Til deres Kreds hørte en Tid
ogsaa \textsc{Jean Jacques Rousseau} (1712---1778), og hans Brud
med dem --- ved hvilket foruden de principielle Modsætninger
ogsaa mange personlige Motiver medvirkede --- var en
kulturhistorisk Begivenhed, et Tegn paa, at et nyt Problem
% --- p. 46 ---
\setcounter{footnote}{0}%
\opage{46}var ved at bryde igennem. Som Hume's Problem gjaldt Videnskabens
Mulighed, saaledes gjaldt Rousseaus Problem Kulturens
Værdi.

Rousseau var født i Genève. Hans urolige, ud over det
sædvanlige Livs Grænser higende Aand kastede ham ind i et
romantisk Vandre- og Eventyrliv, der i Aaret 1741 førte ham
til Paris, hvor han blev Ven med Diderot og Holbach. Da
den store Tanke om Modsætningen mellem Natur og Kultur
var gaaet op for ham, forlod han Paris for at leve paa Landet,
og snart efter brød han med sine encyklopædiske Venner.
Formedelst hans Skrifter (der bleve brændte ligesom Voltaire's,
men ikke vare anonyme som dennes) blev han landflygtig;
selv i Schweiz kunde han ikke finde noget blivende Sted. I
hans senere Aar havde hans Mistænksomhed og Forfølgelsesfrygt
afgjort Karakteren af Sindssyge. Sine sidste Aar levede
han i Frankrig i Tilbagetrukkenhed.

Ved Natur forstaar Rousseau Livets umiddelbare og samlede
Kraft, den uvilkaarlige Udfoldelse, i Modsætning til den
Delthed og Begrænsning, som Kulturen volder, især formedelst
Arbejdets Deling, der ikke lader den Enkelte naa til alsidig
Udvikling, men fremelsker en enkelt Evne paa de andres
Bekostning, hvorved saa Menneskene blive mere afhængige
af hverandre. Den sunde, uvilkaarlige Selvhævdelse (amour
de soi) staar ikke i Modsætning til Medfølelse og Hengivelse;
tvertimod behøver den, hos hvem det indre Liv rører sig
kraftigt, ingen Forskel mellem sig og de Andre at gøre. Det
er kun, naar en usund Sammenligning og Kappestrid gør sig
gældende, at Selvhævdelsen bliver til Egoisme (amour propre).
At vende tilbage til Naturen betyder for Rousseau at vende
tilbage til Livets og Kræfternes Enhed og Fylde, hvor hverken
indre eller ydre Modsætninger gøre sig gældende. I Modsætning
til Kulturens splittende Virkning udraaber han: Rendez
l'homme un! I den frie Natur, som Rousseau ret havde lært
at kende under sin Ungdoms Omstrejfen, rørte denne Enhedsfølelse
sig især levende hos ham. Han er den første Forfatter,
der ret har givet Udtryk for Glæden over Naturensomheden
og Naturskønheden. --- Da Følelsen --- i Modsætning
% --- p. 47 ---
\setcounter{footnote}{0}%
\opage{47}til Erkendelsen og Villien, altsaa til Sammenlignen og Foretrækken
--- ganske særligt har Karakteren af Umiddelbarhed,
Enhed og Fylde, saa han Menneskets inderste Væsen og Værd
udtrykt i den, og det skyldes især ham, at den nu bliver betragtet
som en selvstændig Side af Sjælelivet.

Den uforsonlige Modsætning mellem Natur og Kultur, som
Rousseau hævdede i sine første Skrifter (især \textit{Discours sur
l'origine et les fondements de l'inégalité parmi les hommes,}
1755), blev, efterhaanden som hans Tanker uddybedes og
klaredes, afløst af en relativ Modsætning. I sine senere Skrifter
\textit{(Émile} og \textit{Contrat social,} 1762) skelner han mellem en
naturlig og en unaturlig Kultur. Kultur bestaar i Kræfternes
Udvikling, og en Kultur er naturlig, naar den svarer til det Udviklingstrin,
paa hvilket Mennesket staar, saa at Selvvirksomhed
bliver mulig og Livets samlede Kraft ikke kommer til at lide
under Arbejdets Deling. Hvad Rousseau protesterede mod,
var, at det, Samtiden priste som Kultur, skulde betegne et
ubetinget Fremskridt, og dernæst, at Kulturen uden videre
skulde kunne overføres fra et Folk til et andet. Baade paa
det pædagogiske, det religiøse og det sociale Omraade hævder
han den uvilkaarlige Livsudfoldelse som det ene Fornødne.
--- Al Paavirkning udefra skal kun tjene til at fremlokke det
Ejendommelige hos hver Enkelt; Opdragelsen skal faa Barnet
til at udvikle sine Evner paa egen Haand og har især Betydning
ved at holde Hindringer borte. Barnealderen er en
selvstændig Livsperiode og maa ikke blot behandles som
Middel eller Forberedelse for den voxne Alder. --- I religiøs
Henseende stemmer Rousseau overens med Voltaire, hvad
Forestillingerne angaar. Han tror paa en god Gud, der dog
ikke er almægtig, da det Onde i Verden kun kan forklares
af den Modstand, Materien gør mod den guddommelige Virken.
Dog gør han ikke blot Front mod Materialismen, men
ogsaa mod Voltaire. Han følte en stor Antipati mod Voltaires
Standpunkt, der havde en rent forstandsmæssig Karakter. For
ham selv havde Religionen sin Rod i Følelseslivet. Trangen
til Selvudfoldelse kan være saa stor, at den søger ud over
enhver Grænse og kun finder Hvile i en uendelig Genstand.
% --- p. 48 ---
\setcounter{footnote}{0}%
\opage{48}Den religiøse Følelse naar sit Højdepunkt i en Extase af
Kærlighed, Beundring og overstrømmende Hengivelse. At Religion
udspringer af Naturen, betyder hos Rousseau, at den
udspringer af den indre Følelse. I et Brev siger han: J'ai
laissé là la raison, et j'ai consulté la nature, c'est-à-dire le
sentiment intérieur qui dirige ma croyance. --- I sin Statslære
hævder han Folkets Suverænitet. Den almindelige Villie (la
% PRINTED AS IS (read at the image, odd spots o6)
volonté générale, forskellig fra la volonté de tous) skal raade
Den er Folkets inderste Trang og herskende Tendens, rettet
mod de fælles Interesser, det Heles og de Enkeltes Velfærd
gennem de skiftende Generationer. At lyde den er ikke Afhængighed,
ti hver Enkelt er membre du souverain og kan
finde sin Villie fyldestgjort i den.

2. I Tydskland var det især \textsc{Chr. Wolffs} (1679---1754)
Popularisering af Leibniz's Filosofi, der kom til at ligge til
Grund for Oplysningen. Hans brede, systematiske Skrifter
fortrængte den middelalderlige Filosofi og Teologi fra de tyske
Universiteter. Det var ikke, som hos Leibniz, den metafysiske
Idealisme og Monadelæren, der stod i Forgrunden, men Læren
om den af Gud bevirkede Harmoni i Verden. Hans faste Tro
paa, at der var en tilstrækkelig Grund til Alt, gjorde ham til
Martyr, idet Kong Friedrich Wilhelm I af Preussen paa Grund
af en Anklage for Fatalisme afsatte ham fra hans Professorat
i Halle og --- med Trusel om Galgen --- bød ham forlade
Landet inden 48 Timer. Wolff gik da til Marburg, hvorfra
Friedrich II kaldte ham tilbage strax efter sin Tronbestigelse.
--- Ved Siden af Wolffs Popularisering af det 17. Aarhundredes
Systemer virkede ogsaa Locke's Indflydelse i Tyskland.
Den tyske Oplysningsfilosofi blev da en ofte meget udvortes
Kombination af Wolff og Locke. Efterhaanden traadte den
psykologiske Interesse afgjort i Forgrunden for den metafysiske,
og man ledtes derved til at se, at Modsætningen mellem
Forestillingernes Dunkelhed og Klarhed ikke var den eneste
afgørende. Under Indflydelse af den engelske Literatur og
Rousseau begyndte Følelsen at blive anerkendt som en selvstændig
Side af Sjælelivet. Skyggesiden heraf var den Sentimentalitet,
der prægede Tiden ved Siden af Rationalismen.

% --- p. 49 ---
\setcounter{footnote}{0}%
\opage{49}Under den Strid mellem Rationalismen og Ortodoxien, der
blev Følgen af Oplysningsfilosofiens Optræden i Tyskland,
indtog \textsc{Gotthold} \textsc{Ephraim} \textsc{Lessing} (1729---1781) et Stade,
der pegede ud over Modsætningerne. Skønt han i Karakter
og Begavelse er meget forskellig fra Rousseau, indtager han
dog en analog Stilling i det aandelige Livs Historie. Han
havde en stor og fin Sans for det Oprindelige paa det aandelige
Omraade, det religiøse og det poetiske saavel som det intellektuelle,
og for en indre Sammenhæng, der ligger bagved de
bestemt udprægede Former for Tro og Tænken. Derfor indtog
han en kritisk Stilling baade overfor Rationalister og Ortodoxe.
Som teologisk Kritiker viste han hen til den oprindelige
Kristendom, der var ældre end den saa meget omstridte
Bibel (\textit{Über den Beweis des Geistes und der Kraft}). Og han
satte som Forsker den stadige Stræben efter Sandhed, under
hvilken Kræfterne komme til levende Udfoldelse, over den
blotte Besiddelse af Sandheden, der gør lad og passiv. \textit{(Duplik).}
Sammenhængen i den aandelige Udvikling fandt han ikke i
Dogmer eller Resultater, men i den indre Stræben, de skylde deres
Opstaaen. Sansen for det Oprindelige og Ejendommelige ledte
ham ogsaa som Æstetiker. Han henviste til Shakespeare som det
store Forbillede for dramatisk Kunst, i Modsætning til den herskende
franske Smag, der var endt i det Kunstlede og Indholdsløse.

Lessings eget religiøse Standpunkt er udtalt i den mod de
Ortodoxe rettede Sætning, at man af historiske Beretninger
ikke kan grunde en Erkendelse af Tilværelsens evige Sammenhæng.
De forskellige positive Religioner stode for ham
som Trin i menneskelig Aandsudvikling; hvad Opdragelse er
for den Enkelte, ere de positive Religioner for Slægten. Det
gamle og det nye Testamente ere „Menneskeslægtens Elementarbøger“.
Den Tid vil komme, da saadanne Bøger ville være
unødvendige; da vil „det tredie Rige“, det nye evige Evangelium
komme. \textit{(Erziehung des Menschengeschlechts.} --- \textit{Gespräche
über die Freimaurer.)} --- I sine almindelige filosofiske
Ideer stod han Spinoza nær og erklærede, at hvis han skulde
kalde sig efter Nogen, vilde han kalde sig efter ham.

% --- p. 50 ---
\setcounter{footnote}{0}%
\closerule

\opage{50}\chapterhead{V.}{IMMANUEL KANT OG DEN KRITISKE FILOSOFI.}

Som Lessing og Rousseau tilhørte ogsaa Kant „Oplysningens“
Aarhundrede, men ligesom de pegede han ved de Synspunkter,
han anlagde, ud over det. Som saa mange store
Aander har han modtaget Paavirkninger fra mange Sider.
Wolff, Newton, Hume og Rousseau have især faaet Indflydelse
paa hans Filosofi, saaledes som han efter en lang Periode
af Selvtænkning fremstillede den i sit Hovedværk \textit{Kritik
der reinen Vernunft} (1781). Dette Værks Opgave er en Prøvelse
af den menneskelige Tankes Evner til at naa Videnskab
uden at støtte sig til Erfaring; ved „den rene Fornuft“ menes
netop Tanken, saaledes som den virker uden Erfaringens
Hjælp. At al Erkendelse \textit{begynder med} Erfaring, derom tvivler
Kant ikke; men derfor var det jo ikke sagt, at den altid
\textit{maa grunde sig paa} Erfaring. Det 17. Aarhundredes Filosofer
havde antaget, at der kunde gives Videnskab ud over Erfaringens
Omraade, men paa den anden Side havde Hume
hævdet, at selve det Princip, der ligger til Grund for al Videnskab
om Tilværelsen, Aarsagssætningen, kun grundede sig paa
en Vane og ikke kunde videnskabeligt begrundes. Det var
netop Hume, der, som Kant selv har udtrykt det, vakte ham
af hans dogmatiske Slummer. Ti i hans første Periode, som
Tilhænger af Wolffs Filosofi, havde han trygt stolet paa Videnskabens
Grundsætninger og paa Berettigelsen af at anvende
dem ud over Erfaringen. Paavirkningen fra Hume indledede
for Kant det lange Arbejde, der endte med den definitive
Grundlæggelse af den kritiske Filosofi, hvis Hovedopgave
er en Undersøgelse af den menneskelige Videnskabs Forudsætninger
og deres Gyldighed. Kant er klar over, at den
kritiske Filosofi ikke kan være Menneskets første Filosofi.
Fra først af tvivler Mennesket ikke om sin Erkendelses Gyl%
% --- p. 51 ---
\setcounter{footnote}{0}%
\opage{51}dighed og antager ikke, at den har Grænser. Den første Filosofi
er derfor stedse \textit{dogmatisk,} d. v. s. den anstiller ingen
Prøvelse af de Grundbegreber og Grundsætninger, hvorfra
den gaar ud, men stoler trygt paa dem og mener ved deres
Hjælp at kunne løse alle Gaader. Naar Dogmatikerne komme
i Strid med hverandre, kommer der naturligt en \textit{skeptisk} Periode,
i hvilken man tvivler, om der idethele kan gives gyldig
Viden. Men den \textit{kritiske} Filosofi sætter sig saa den Opgave
at dømme Dogmatikere og Skeptikere imellem ved metodisk
Prøvelse af Betingelserne for Videnskab og Undersøgelse af,
paa hvilke Omraader disse Betingelser er tilstede. --- Det er,
som Titelen paa Kants Hovedværk viser, ikke al Videnskab,
han vil undersøge; den, der blot grunder sig paa Erfaring, f.
Ex. beskrivende Naturvidenskab og Psykologi, drager han
ikke ind med.

Foruden Wolff og Hume have, som allerede sagt, Newton
og Rousseau øvet en dybtgaaende Indflydelse paa Kants Udvikling.
--- \textsc{Isaac} \textsc{Newton} (1642---1727) stod for ham som
Mønster paa den strenge naturvidenskabelige Opfattelse af
Verden. I sine \textit{Principia philosophiæ naturalis mathematica}
(1687) havde Newton givet det største Exempel paa exakt,
Induktion og Deduktion forbindende Videnskab. Udgaaende
fra den Erfaring, at Tyngden er større i Dalene end paa
Bjerghøjderne, og at Legemer, der kastes i Vejret, falde ned,
formodede han, at ogsaa Planeterne ere tunge, og at de kun
afvige saa meget fra den Retning, de ifølge Inertisætningen
vilde bevæge sig i, som den paa Jordens Overflade af Galilei
paaviste Faldlov angiver. Han drog saa de matematiske Konsekvenser
af denne Formodning og viste endeligt, at Resultaterne
af denne Deduktion stemmede med Iagttagelserne.
Planeterne bevæge sig altsaa efter samme Lov som Legemer,
der kastes i Vejret fra Jorden. Og af denne Antagelse kunde
saa de af Kepler paa Grundlag af Tycho Brahes Iagttagelser
udfundne Love (bl. a. Loven om Planetbanernes elliptiske
Form) udledes. Det blev nu Kants Opgave at finde, hvilke
Forudsætninger Newton's Videnskab hvilede paa, og undersøge,
hvorvidt disse Forudsætninger kunde bevises. Men han
% --- p. 52 ---
\setcounter{footnote}{0}%
\opage{52}prøvede tillige paa at fortsætte Newton's Værk. Denne havde
erklæret, at den vidunderlige Harmoni i vort Solsystem ---
idet Planeter og Drabanter bevæge sig i samme Retning, omtrent
i samme Plan og i næsten kredsformige Baner --- ikke
kunde forklares ved Naturaarsager. Kant viste (i Allgemeine
\textit{Naturgeschichte und Theorie des Himmels} 1755) at naar
man tænkte sig, at alle Kloder i Solsystemet engang havde
udgjort en stor roterende Taagemasse, vilde man kunne forklare
den nu bestaaende Ordning efter almindelige Naturlove.
Under Rotationen vilde der formedelst midtpunktflyende Kraft
udskille sig Ringe af Taagens Ækvator; disse vilde igen spaltes
i Kloder i forskellige Afstande, men alle roterende i samme
Retning og i omtrent samme Plan. Det var den samme Hypotese,
der (uafhængigt af Kant) senere blev opstillet af Laplace,
men som nu er omtvistet. Det var en Hovedtanke hos
Kant, at alle Fænomener underligge Aarsagsloven, og at selv
den Orden og Hensigtsmæssighed, der findes i Naturen, maa
være forklarlige efter Naturlove. Des større et Problem opstod
der for ham, da han senere ved Humes Kritik kom til
at tvivle om, at Aarsagssætningen overhovedet kunde bevises.
--- Rousseau's Indflydelse førte Kant til en ny Menneskevurdering:
forhen havde han betragtet Erkendelse og Oplysning
som det Højeste og foragtet Pøbelen, der Intet véd;
Studiet af Rousseau bragte ham til at finde en indre Kærne
i Menneskets Følelse og Villie, der kan bestaa med en vis
Uafhængighed af Oplysning og Kundskab. Ogsaa paa den
senere Udvikling af Kants Etik har Rousseau øvet Indflydelse.
---

Kants Liv (1724---1804) bevægede sig, hvad det Ydre angaar,
i en snever Kreds. Hans Fader var en fattig Haandværksmand
i Kønigsberg. Paa Universitetet i Kønigsberg studerede
han Wolffs Filosofi og Newton's Fysik. Efterat han i
nogle Aar havde været Huslærer i forskellige adelige Familier
i Østpreussen, virkede han i en Række af Aar som Privatdocent
ved Universitetet i Kønigsberg med stor Tilslutning;
først 1770 blev han Professor. Han har aldrig været udenfor
Østpreussen; hans Tid var tagen i Beslag af Udarbejdelsen
% --- p. 53 ---
\setcounter{footnote}{0}%
\opage{53}af hans Værker og af hans akademiske Undervisning. Dog
tog han Del i det selskabelige Liv og havde Ry som aandrig
og zirlig Selskabsmand. Han hørte Oplysningens Tid til, men
forstod ved „Oplysning“ en Stræben, ikke Noget, man en
Gang for alle kunde opnaa, end sige nu havde naat. Og medens
hans kritiske Tanke førte ham i Dybder, den sædvanlige
Oplysning ikke kendte, havde han en Sans for det Ophøjede,
der var Grundtonen i hans Livsanskuelse, næret ved
Naturbetragtning og ved hans Overbevisning om de etiske
Idealers ubetingede Gyldighed. I sin Alderdom kom han til
at lide under den klerikale Reaktion efter Friedrich IIs Død.
Da han i Aaret 1793 havde udgivet et religionsfilosofisk Skrift,
rettedes der et kongeligt Reskript til ham, der paatalte hans
frie Opfattelse og truede med videregaaende Forholdsregler,
naar han fortsatte i samme Retning. Kant ansaa det for sin
Pligt at lyde, men erklærede, at han nu slet ikke vilde skrive
eller tale om religiøse Emner. Da der kom en ny, mere frisindet Regering, genoptog han sin religionsfilosofiske Virksomhed
og offentliggjorde den forargelige Episode. Hans sidste
Aar frembyde Billedet af en stor Aands Opløsning; sandsynligvis
har han lidt af en Hjernesygdom; kun enkelte Glimt
erindrede endnu om, hvad han havde været. ---

Kants Filosofi kan fremstilles i tre Afsnit, efter hans tre
vigtigste Værker, af hvilke det ene giver en Erkendelsesteori,
det andet en Etik, og det tredie søger at naa ud over begges
Resultater i en afsluttende Opfattelse.

1. \textit{Erkendelsesteori.} (Kritik der reinen Vernunft).

a. Idet Kant vil undersøge, hvor langt vi kunne komme ad
den rene Fornufts Vej, maa han først udfinde, hvad den rene
Fornuft er. Han analyserer da vor Erkendelse og finder, at der
er Noget, der stadigt gør sig gældende, hvad vi saa erkende eller
erfare, og som vi selv maa medbringe, nemlig de Former,
i hvilke vi samle og ordne vore Iagttagelser. Erkendelse er
en aandelig Virksomhed, hvis Love og Virkeformer maa
kunne spores i al Erfaring og al Erkendelse. I Modsætning
til disse Former kalder Kant de blotte Fornemmelser og Iagttagelser
for Stoffet i vor Erkendelse. Det modtage vi passivt.
% --- p. 54 ---
\setcounter{footnote}{0}%
\opage{54}Det er kun Formerne, Kant efter sin Opgave har med at
gøre, og af dem finder han, at der faktisk gør sig tre Arter
gældende, der dog alle frembyde det samme Grundpræg og
ere Ytringer af en og samme Grundform.

I enhver Sansning er der Noget, som er uafhængigt af de
specielle Sansekvaliteter, jeg opfatter: jeg kan i et og samme
Rumsbillede sammenfatte forskelligt farvede Ting til forskellige
Tider, og jeg kan i den samme Tidslængde opleve højst
forskellige Begivenheder. Alt hvad jeg sanser, hvor forskelligt
og skiftende det end er, sammenfatter jeg i Rum eller i Tid
eller i begge Former. Rum og Tid ere \textit{Anskuelsesformer,}
ved hvis Anvendelse Fornemmelsernes Kaos sammenfattes
og ordnes til Sansebilleder.

Men vi blive ikke staaende ved Sansebillederne. Disse
sammenfattes og ordnes igen af Eftertanken („Forstanden“)
og stilles i bestemte Forhold til hverandre. Ogsaa herved gør
sig visse Former gældende: Forstandsformerne eller \textit{Kategorierne}
(Grundbegreberne). Ad en meget kunstig Vej, som vi
her ikke kunne gaa ind paa, søger Kant at vise, at der er 12
saadanne Kategorier; men han antyder, at de kunne reduceres
til to: Størrelses- og Aarsagskategorien. \textit{Størrelsesbegrebet}
anvende vi, naar vi gaa ud fra, at Alt, hvad vi erfare, maa
have en vis Udstrækning i Rummet og en vis Styrkegrad, og
at der mellem alle Punkter i Rummet og mellem alle Grader
paa Styrkeskalaen er en kontinuerlig Overgang mulig.
\textit{Aarsagsbegrebet} anvende vi, naar vi gaa ud fra, at de Begivenheder,
vi iagttage, maa kunne forbindes saaledes, at der
ikke blot er et Tidsforhold mellem dem, men ogsaa et saadant
Forhold, at den ene Begivenhed for os indeholder
Grunden til den anden, saa at der er en Analogi mellem
Forholdet Grund---Følge paa den ene Side og Tidsforholdet
mellem Begivenhederne paa den anden Side. Begge Kategorier
udtrykke Fordringen om \textit{kontinuerlig Sammenhæng} og
virke sammen til, at vor Erfaring kan udgøre en Enhed. I
Kraft af Størrelsesbegrebet gaa vi kontinuerligt fra Sted til
Sted og fra Grad til Grad, og i Kraft af Aarsagsbegrebet gaa
vi kontinuerligt fra Begivenhed til Begivenhed. Der skal in%
% --- p. 55 ---
\setcounter{footnote}{0}%
\opage{55}gen Spring være i Naturen, hverken i Udstrækning eller i
Grad eller i Aarsagernes Række.

Baade Anskuelsesformer og Kategorier ere Former for
vor erkendende Bevidsthedsvirksomhed, der paa alle Trin
viser sig som en Sammenfatten, en \textit{Syntese.} De ere forskellige
Maader at sammenfatte det givne Stof paa. Ingen Erfaring
er mulig uden derved, at Fornemmelserne forbindes til Sansebilleder,
og Sansebillederne forbindes til kontinuerlige Rækker
i Tid, Rum, Grad og Kausalitet. Saadan Syntese behøver ikke
at foregaa med fuld Bevidsthed. Den er, siger Kant, Virkning
af en blind, men uundværlig sjælelig Funktion, uden hvilken
vi slet ikke vilde have nogen Erkendelse, men som vi kun
sjeldent blive os bevidst. Al Analyse, vi kunne foretage, forudsætter,
at der er forbundet Noget i vor Bevidsthed; ti vi
kunne kun opløse, hvad der er forbundet.

Anskuelsesformer og Kategorier udtrykke Forbindelser af
det givne Stof. Men den uvilkaarlige, i Bevidstheden dybt
grundede Trang til at forbinde og sammenfatte er saa stærk,
at den vil omfatte Alt til en afsluttende Helhed. En Afslutning
af alle extensive og intensive Rækker og af Aarsagsrækken
kan ikke tænkes indenfor Erfaringen, der stedse er begrænset.
Forsaavidt Syntesen fører til Begreber, der kræve Afslutning
ved Dannelse af absolute Helheder, danner den \textit{Ideer,}
en tredie Gruppe af Former. En saadan Ide er Begrebet
„Verden“, der skal omfatte alt Existerende, skønt ingen Grænse
kan sættes for Rummet, Tiden og Aarsagsrækken. Ligeledes
Begrebet „Gud“ i Betydning af en første Aarsag, der ikke
igen har nogen Aarsag, og Begrebet „Sjæl“ i Betydning af en
Substans, der ligger bag Bevidstheden og ikke kan opdages
gennem Selviagttagelse. Det er i og for sig naturligt, at saadanne
Ideer ere dannede; vor Stræben efter at sammenfatte
er Et med vor inderste Natur.

b. Den Undersøgelse, ved hvilken Erkendelsesvirksomhedens
Former udfindes, er beskrivende. Den kan kun paavise,
paa hvilken Maade vor Erkendelse faktisk virker, men
ikke Gyldigheden og Berettigelsen af den Maade, den virker
paa. Kant gør da Overgangen fra Psykologi til Erkendelses%
% --- p. 56 ---
\setcounter{footnote}{0}%
\opage{56}teori („Fornuftkritik“) ved at spørge, hvorledes en Erkendelse,
der blot støtter sig til Formerne, ikke til Stoffet, --- altsaa en
Erkendelse af „ren Fornuft“, kan være Videnskab. Problemet
opstaar for Kant (ligesom allerede for Locke) derved, at Erkendelse
er en Virksomhed fra vor Side; med hvilken Ret
tillægge vi de Former, i hvilke denne Virksomhed fremtræder,
Gyldighed for Tingene? --- Kants Svar er, at hvis det viser
sig, at Anvendelsen af vore Erkendelsesformer er \textit{Betingelse
for, at vi kunne have virkelig Erfaring} (i Modsætning til blot
Forestilling, Fantasi eller Drøm), saa maa Sætninger, der
bygge paa disse Former, ogsaa gælde for \textit{al mulig Erfaring}
--- men heller ikke videre. Han undersøger da i denne Henseende
de forskellige Grupper af Former.

Alt Erfaringsindhold opfattes \textit{i Rum eller Tid eller baade
i Rum og Tid.} Da nu den rene Matematik egentligt kun udvikler
Rummets og Tidens Love, maa den gælde for Alt, hvad
vi nogensinde kunne erfare, for alle mulige Iagttagelser. Heraf
følger Gyldigheden af den anvendte Matematik, uagtet Matematiken
i sig selv er en ren Fornuftvidenskab. Vi kunne paa
Forhaand (apriori) vide, at alle Erfaringer maa stemme med
matematiske Love. Men da Rum og Tid ere menneskelige
Opfattelsesformer, tør vi ikke sige, at den Erkendelse, vi her
naa, er den eneste mulige. Vi kende kun Tilværelsen fra
Menneskets Standpunkt. Hvad Kopernikus hævdede, at vi
ved Opfattelsen af Himmellegemernes Bevægelser ikke maa
glemme, at vi se dem fra det menneskelige Stade, det gælder
for al vor Erkendelse, al vor Videnskab, fordi den skyldes
et Arbejde, der foretages i menneskelige Former.

Paa lignende Maade forholder det sig med \textit{Kategorierne.}
Da de ere Former, uden hvilke vi ikke forstaa Tingene, maa
Alt, hvad vi skulle erfare, falde ind under dem. Ethvert Fænomen
eller enhver Begivenhed maa fyldestgøre Størrelsesbegrebets
Fordringer: deres Opstaaen og Udvikling maa foregaa
kontinuerligt, saaledes at en højere Styrkegrad stedse
naas ved Gennemløben af alle lavere Styrkegrader, og at et
fjernere Punkt i Rummet stedse naas ved Gennemløben af
alle Punkter, som ligge mellem Udgangspunktet og det. Hvis
% --- p. 57 ---
\setcounter{footnote}{0}%
\opage{57}nemlig en saadan kontinuerlig Overgang ikke kan paavises,
kunne vi ikke skelne mellem Erfaring og Indbildning, Virkelighed
og Drøm. Det Samme gælder for Aarsagsbegrebets
Vedkommende. Hvis jeg ikke kan ordne Begivenhederne i
Rækker, indenfor hvilke hvert Led forholder sig til det følgende
paa analog Maade, som i vor Tænkning Grund forholder
sig til Følge, kan jeg ikke have virkelig Erfaring. Naar
jeg kan forestille mig Begivenhederne i hvilkensomhelst Orden,
vilkaarligt lade dem bytte Plads, kan jeg ikke være vis
paa, om jeg staar overfor noget Virkeligt. I min Betragtning,
i min Erindring og min Fantasi kunne Forestillingerne komme
i forskellig Orden; men naar mine Forestillinger skulle udtrykke
en i Erfaringen given Virkelighed, maa de danne en
Række, i hvilken Leddene ikke kunne byttes om. I min Betragtning
eller min Fantasi kan jeg ligesaa godt lade Huset
blive til for mig ovenfra nedad som nedefra opad; men skal
Husets Tilblivelse være en virkelig Begivenhed, er kun én
eneste Rækkefølge mulig. (Smlgn. Psykologi V D, 2). Alt hvad
vi nogensinde skulle kunne erfare, maa underligge Aarsagssætningen.
Skønt denne Sætning er en ren Fornuftsætning,
idet den bygger paa en af Formerne for vor sammenfattende
Virken, maa den altsaa paa Forhaand (apriori) erklæres gyldig
for alle Erfaringer. --- Men Kategoriernes Gyldighed er,
ligesom Anskuelsesformernes, begrænset til mulig Erfaring.
Længere gaar Begrundelsen ikke. Erfaring er, siger Kant, den
Syntese, der giver al anden Syntese Gyldighed.
Ved Ideerne kan den Begrundelse, som ved Anskuelsesformer
og Kategorier førte til Maalet, ikke anvendes. Ti det
Ejendommelige ved Ideerne var, at de krævede en afsluttende
Syntese, en ubetinget Helhed, et Første eller Sidste. Men et
saadant kan Erfaring ikke frembyde, da Erfaringen altid er
begrænset. Ideerne kunne altsaa ikke siges at angive Betingelser
for at skelne mellem Erfaring og Indbildning. Verden
f. Ex. kan ikke fremtræde i nogen Erfaring; ud over enhver
Helhed, vi danne ved Syntese, kan der stedse tænkes en
mere omfattende Helhed. Paa Ideerne kan der derfor ingen
Videnskab bygges. Der er ingen Videnskab om Gud, om
% --- p. 58 ---
\setcounter{footnote}{0}%
\opage{58}Verden som absolut Helhed eller om Sjælen som Substans.
Kant anerkender en Erfaringspsykologi; men den er ingen ren
Fornuftvidenskab. --- Ideernes Betydning ligger derimod i, at
de pege ud over enhver enkelt Erfaring, hindre os i at slaa os
til Ro ved nogetsomhelst Resultat; de stille stedse nye Opgaver.
Kant mødes her med Lessing i Tanken om den stadige
Stræben efter Sandhed i Modsætning til den blotte Besiddelse.
---

Naar Kant grunder Fornufterkendelsens reale Betydning
derpaa, at den udtrykker Betingelserne for at have Erfaring,
tager han Begrebet Erfaring i streng Betydning som fast og
nødvendig Sammenhæng af Iagttagelser. Han har ikke set, at
Erfaring i denne Betydning selv er en „Ide“ eller et Ideal,
da Iagttagelserne kun ufuldkomment og tilnærmelsesvis fyldestgøre
Anskuelsesformernes og Kategoriernes Fordringer. Kant
har derfor ikke ført noget afsluttende Bevis for Aarsagssætningens
Gyldighed og ikke løst Humes Problem.

Den Begrænsning, han paaviste for den rene Fornuft i
dens forskellige Former, udledte han af, at vi ikke have Ret
til at betragte disse Former som de eneste, i hvilke Tilværelsen
kunde opfattes og forstaas. Vi maa skelne mellem den
Maade, paa hvilken vi erkende, og selve Tingenes Væsen. Vi
erkende kun Tingene, som de fremtræde for os, efter vore
Former, ikke som de ere i sig selv, eller som Kant udtrykker
det, vi erkende kun Fænomener (Erscheinungen) ikke
Tingen i sig selv (das Ding an sich). Dette begrunder han
ogsaa ved, at medens Formerne ere vore Maader at virke
paa i vor Erkendelse, modtage vi Stoffet (Sansekvaliteterne)
rent passivt. Medens Formerne høre til vor Natur, kan Stoffet
ikke udledes af denne. Begrebet „Ding an sich“ tjener da
til at markere en Grænse. Det er et \textit{Grænsebegreb,} der slaar
fast, dels at vore Former \textit{ikke} ere de eneste mulige, dels at
vor Erkendelses Stof \textit{ikke} kan udledes af dens Form. Det er
altsaa et negativt Begreb. Men uvilkaarligt kommer Kant til
at tilskrive det positiv Betydning, idet han kalder „Ding an
sich“ Aarsag til Fænomenerne (især hvad Stoffet angaar, dog
ogsaa hvad vore Former angaar). Herved kommer han i Mod%
% --- p. 59 ---
\setcounter{footnote}{0}%
\opage{59}sigelse med sin egen Lære om, at man ikke har Ret at anvende
Aarsagsbegrebet ud over Erfaring; ti „Ding an sich“
kan ikke erfares. Og desuden kommer det Spørgsmaal frem,
om vi kunne være sikre paa, at „Ding an sich“ stedse vil
give os Stof, der passer til vore Former, og stedse vil lade
de samme Former bestaa hos os. Ogsaa paa dette Punkt,
ikke blot ved Begrebet Erfaring, viser det sig, at Kant ikke
har frigjort sig fra al Dogmatisme.

2. \textit{Etik.} (Kritik der praktischen Vernunft. 1788).

Under stærk Paavirkning af Rousseau, saavel som af
Hume og andre engelske Filosofer lagde Kant i sine første
Udtalelser om etiske Spørgsmaal megen Vægt paa Følelsen af
„den menneskelige Naturs Skønhed og Værdighed“ som
Grundlag for de etiske Principer \textit{(Betrachtungen} über das
Schöne \textit{und Erhabene} 1764). Senere lagde han mere Vægt
paa Villien, Selvvirksomheden og paa Nødvendigheden af, at
vi stedse i vor Stræben ere i Overensstemmelse med os selv.
Hans Etik fik derved en overvejende formal Karakter. Da nu
historiske Studier førte ham til at betragte de enkelte Individers
Stræben som Led i hele Slægtens Stræben efter at
hæve sig fra den blotte Naturtilstand til sand Kultur, maatte
det etiske Princip ikke blot blive den Enkeltes Overensstemmelse
med sig selv, men maatte blive et saadant, som alle
Individer i samme Situation konsekvent kunde følge. Jeg skal
da handle saaledes, at Principet for min Handling kan opstilles
som Regel for Alle. Jeg skal f. Ex. give betroet Gods
tilbage, selvom jeg uden at blive opdaget kunde beholde det;
ti det modsatte Princip kan ikke opstilles som Regel, da saa
Ingen vilde betro Andre Noget. Ethvert Menneske skal behandles
som Formaal, aldrig blot som Middel; kun paa den
Maade bliver et menneskeligt Samfund muligt, i hvilket den
Enkeltes Frihed kan bestaa med Andres Frihed. Ligesom vi
kun faa gyldig Erkendelse, naar alle Sansefornemmelser (alt
„Stof“) ordnes efter Love, der tilsidst udspringe af vor Aands
Former, saaledes er vor Villen og Handlen kun gyldig, naar
den ikke følger Følelser og Drifter (hvad enten disse ere
sympatiske eller egoistiske), men underkaster sig en Fornuft%
% --- p. 60 ---
\setcounter{footnote}{0}%
\opage{60}lov, der kan gælde for Alle. En saadan Lov kundgør sig som
et ubetinget Bud (et kategorisk Imperativ) overfor alle andre
Tilskyndelser, overfor alle Autoriteter, og overfor Alt hvad
Erfaringen lærer eller synes at lære. Aller tydeligst fremtræder
det Etiske i saadanne Tilfælde, hvor Pligt og Tilbøjelighed
staa i den største Modsætning til hinanden. Der vil altid,
mener Kant, være et Element af Ulystfølelse i Pligtbevidstheden
paa Grund af den Modstand, de sanselige Dele af vor
Natur gøre mod at underordne sig under Lovens Form.

I Evnen til saaledes trods al indre og ydre Modstand at
underlægge sin Handlen under en almengyldig Lov ser Kant
Menneskets Adelsmærke. Loven er jo ikke en fremmed Lov,
men grundet i Menneskets Natur som Fornuftvæsen, Et med
hans inderste Villie, saa ophøjet dens Myndighed end er.
Mennesket er gennem den etiske Lov Borger i en anden
Verden end Erfaringens. Inderlighed (ved Udspringet af Menneskets
inderste Natur) og Ophøjethed (ved Modsætningen til
den i Erfaringen fremtrædende Menneskenatur) ere de to
karakteristiske Træk i den etiske Lov efter Kants Opfattelse.
Han antager denne Lov for rent formel og apriorisk, udsprungen
af den rene Fornuft. Men den forudsætter dog tydeligt,
at der skal bestaa et menneskeligt Samfund, og formulerer
egentligt kun dettes vigtigste Betingelser. Det var ogsaa
hos Kant selv historiske Overvejelser, der havde ført ham
til hans sluttelige etiske Opfattelse. ---

Der er efter Kant et indre Forhold mellem Religion og
Etik. Overgangen fra det Etiske til det Religiøse finder han
deri, at Mennesket skal hævde det Gode, det, som den ubetingede
etiske Lov kræver, under indre og ydre Modstand.
Deraf følger en skarp Modsætning mellem Pligt og Lykke,
mellem Ideal og Virkelighed, Frihed og Natur. Denne Modsætning
vil igen afføde en Trang til at tro paa en anden Tilværelse,
i hvilken den Stræben, der afbrødes i Erfaringens
Verden, kan fortsættes i stedse højere Former, og til at tro paa
et Væsen, der mægter at bringe Pligt og Lykke, Ideal og
Virkelighed i Harmoni. Troen paa Gud og paa Udødelighed,
den saakaldte naturlige Religions Indhold, kan dog ikke viden%
% --- p. 61 ---
\setcounter{footnote}{0}%
\opage{61}skabeligt begrundes. Det er praktiske Postulater, der føre til
den, Postulater, som ikke kunne modbevises, da Videnskaben
selv har sin Grænse (i „Ding an sich“). Deres Betydning
bestaar i, at de holde os oppe i vor Stræben efter at hævde
den etiske Lovs Gyldighed. Men de udvide ikke vor Erkendelse;
det følger allerede deraf, at vor Erkendelses Former
kun gælde indenfor Erfaringen. Alle religiøse Forestillinger
grunde sig derfor kun paa Analogi; de ere symbolske og
billedlige. --- I sit Skrift \textit{Religion innerhalb der Grenzen der
blossen Vernunft} (1793), --- det Skrift, der bragte Kant i
Konflikt med Statsmyndigheden, --- har han gennemført en
symbolsk Opfattelse af de kristelige Dogmer, i hvilke han
finder dybe etiske Sandheder indeholdte.

3. \textit{Enhedslære paa æstetisk og biologisk Grundlag.} (Kritik
der Urtheilskraft. 1790).

I sin Erkendelsesteori og i sin Etik havde Kant hævdet
en skarp Modsætning mellem Erfaringsverdenen og „Tingene
i sig selv“. Ja, hans hele Filosofi er karakteriseret ved skarpe
Distinktioner, der efter hans Mening vare nødvendige, naar
Erkendelsens Gyldighed og det etiske Ideals ubetingede Karakter
skulde hævdes. Men netop ifølge den kritiske Filosofis
Principer, som jo indskærpe Uberettigelsen af at overføre
vore Tanker paa Tilværelsen i sig selv, maatte den
Tanke opstaa, om ikke ogsaa selve hine Distinktioner og
Modsætninger maatte være at skrive paa vor menneskelige
Opfattelses Regning. Kant fandt to Grupper af Kendsgerninger,
som kunde pege i Retning af en Enhed, den ene Gruppe
af æstetisk, den anden af biologisk Art.

a. I de Fænomener, vi kalde skønne og ophøjede, frembringer
Tilværelsen et uinteresseret Velbehag hos os. Det er
ikke vor Sanselighed, der tilfredsstilles, naar vi staa overfor
noget Skønt. Det Skønne beror paa, at vor Anskuelsesevne
og vor Forstand let og uvilkaarligt kunne forbinde den givne
Mangfoldighed til en Enhed, ret somom Fænomenet var afpasset
efter vor Opfattelsesevne. Kant tænker her især paa
Formskønheden, som den viser sig i Harmoni og Symmetri,
i Naturens og i Kunstens Værker. Ved det Ophøjede over%
% --- p. 62 ---
\setcounter{footnote}{0}%
\opage{62}vældes vor Opfattelsesevne og vor Selvfølelse, idet vi staa
overfor det i Udstrækning og Kraft Umaalelige; men netop
derved vækkes Bevidstheden om en i os selv boende, over
alle Grænser stræbende Kraft, den, der ytrer sig i den etiske
Lov. Baade ved det Skønne og det Ophøjede viser det sig
altsaa, at der er Noget i Tilværelsen, som stemmer med det
i os, der er hævet over Sanselighed og Egoisme. --- I Geniet,
i den geniale Frembringen, er der et Træk, som peger i
samme Retning. Ti i Geniet virker en uvilkaarlig Naturkraft
og Naturtrang, og dog frembringes derved Værker, der staa
som Forbilleder. I Geniet virker Naturen, somom den var
bevidst Kunst. Heller ikke her gør da Modsætningen mellem
Mennesket og Tilværelsen sig gældende.

b. I det organiske Liv have vi en Kendsgerning, der kan
tydes paa analog Maade som det Skønne, det Ophøjede og
Geniet. Naturen virker i det Organiske paa en Maade, for
hvilken vi egentligt ikke have noget Udtryk. Vi se her et
Hele, der ikke opstaar ved ydre Sammenføjning af Delene,
og ved hvilket vi videnskabeligt set heller ikke have Ret til
at antage en forudgaaende Plan eller et Formaal (som ved
Tilblivelsen af en menneskelig Bygning). Naturen virker ogsaa
her, somom den var bevidst Kunst. Hverken den teleologiske
eller den mekaniske Naturforklaring kan tyde Organismen.
Men Modsætningen mellem Teleologi og Mekanik beror maaske
kun paa vor Forstands Natur: et Hele kunne vi kun
forstaa, naar enten Delene føjes sammen ved ydre Aarsager,
eller der gaar en Plan forud. I Tilværelsen selv kunde derimod
samme Naturorden, der gør Aarsagssammenhængen
mulig, ogsaa muliggøre Opstaaen af Organismer, hensigtsmæssigt
byggede Livsformer. I Naturens ubekendte Grund kunde
Teleologi og Mekanisme være Et. Det var en Tanke, Kant
allerede i sin Ungdom, da han udarbejdede sin Hypotese om
Solsystemets Udvikling, havde lagt saa stor Vægt paa.

Og paa samme Maade kunde det jo forholde sig med den
største Modsætning: mellem de Naturlove, Videnskaben paaviser,
og de Idealer, den etiske Lov fører til at anerkende.
Det vilde da være et og samme Princip, som kundgjorde sig
% --- p. 63 ---
\setcounter{footnote}{0}%
\opage{63}i Naturlovene og i de etiske Idealer. Det vilde da ikke blot
være gennem Postulater, at hin Modsætning kunde overvindes.

Kant er imidlertid klar over, at en saadan Enhedslære ikke
kan videnskabeligt gennemføres.

\closerule

\chapterhead{VI.}{ROMANTIKENS FILOSOFI.}

a. Den nyere Filosofi fik, som vi have set, sin Opgave og
sin Karakter bestemt dels ved det nye Verdensbillede og den
nye Videnskab, dels ved den friere Opfattelse af menneskelige
Forhold, der kom med Renaissancen. Senere bestemtes
Kants Filosofi ikke blot ved Modsætningen mellem Wolff og
Hume, men ogsaa ved Newton's Naturvidenskab og Rousseau's
Kulturproblem. Den tyske Filosofis Udvikling i den
nærmeste Tid efter Kant bestemtes ved Indflydelsen af den
æstetiske og religiøse Retning, der i Literaturhistorien kaldes
Romantik, og som betegner en stærk Reaktion mod det
attende Aarhundredes Oplysning og Revolution. Man var træt
af Kritik og Analyse og følte Trang til umiddelbart at fordybe
sig i Naturen og i Historien, føle sig Et med Tilværelsens
inderste Kræfter. Spaltningen mellem Poesi, Religion og
Videnskab, denne Frugt af Arbejdsdelingen paa det aandelige
Omraade, skulde ophøre; det aandelige Liv skulde udgøre en
Enhed, og Poesi, Religion og Videnskab skulde kun være
forskellige Former for denne Enhed. Man vilde overhovedet
ud over alle Modsætninger, og derfor ogsaa ud over Kant's
Filosofi, der saa stærkt havde fremhævet Modsætningerne.
Man tog sit Udgangspunkt fra den Tanke, Kant saa forsigtigt
havde antydet, om en Enhed bag de Modsætninger, vi stanse
ved, --- en Tanke, der for Kant betød Tankens Grænse. Og
% --- p. 64 ---
\setcounter{footnote}{0}%
\opage{64}det blev da Idealet at udlede Alt af ét Princip. I sin Begejstring
for dette Ideal undersøgte man ikke dets Berettigelse
og Muligheden af at virkeliggøre det og oversaa, at enhver
Slutning jo dog maa bygges idetmindste paa to Præmisser.
Tænkningens og Videnskabens hidtil gældende Metoder foragtede
man og fordrede en spekulativ Tænken, der kunde
løse de højeste Problemer. Og man brød den Sammenhæng
med Naturvidenskaben, Filosofien havde havt siden Brunos
og Descartes' Dage. Ud fra Anelser om en Enhed i Alt vilde
man nu selv konstruere Naturens Former som en Række af
Trin op til Aandens Verden. Det blev undertiden Poesi, ofte
Fantasteri, ikke Videnskab, man naaede til ad denne Vej. Og
dog har Romantikens Filosofi ikke blot udtalt dybsindige Tanker
af blivende Værdi, men ogsaa banet Vejen for en Forstaaelse
af Sammenhængen mellem Naturens Kræfter, som
senere kunde gennemføres paa mere exakt Maade, ligesom
den ogsaa bidrog til at vække og nære historisk Sans.

De tre vigtigste Repræsentanter for Romantikens Filosofi
ere Fichte, Schelling og Hegel.

J. G. \textsc{Fichte} (1762---1814) en saksisk Bondesøn, fulgte i
sin første Ungdom med levende Interesse Lessings Optræden
og blev senere, efter heltemodigt at have arbejdet sig gennem
trange Studieaar, ført til Filosofien ved Studiet af Kants Skrifter.
Paa Universitetet i Jena virkede han som Lærer med
stor Tilslutning, ikke blot paa Grund af sin Tankekraft og
Veltalenhed, men ved sin hele energiske Personlighed. Da
han paa Grund af sine religiøse Anskuelser blev afsat, gik
han til Berlin, hvor han blev en af det nye Universitets første
Lærere. Han var en af de Mænd, der virkede for Preussens
nationale Genrejsning efter Jenaslaget, især ved sine
under den franske Okkupation holdte „Reden an die deutsche
Nation“.

Fichte tog sit Udgangspunkt fra Kants Modsætning mellem
Tænkningen og „Ding an sich“. Tænkningen er en Virksomhed,
en Handling, og Tingene ere dels Genstande for denne
Handling, dels dens Produkter. Men hvorledes er det da at
forklare, at der existerer Ting som noget fra vor Aandsvirk%
% --- p. 65 ---
\setcounter{footnote}{0}%
\opage{65}somhed Forskelligt? Tingene som Genstande kunne ikke udledes
af selve det tænkende og handlende Jeg; dette kan ikke
selv fremkalde sine Skranker, sin Modstand. Den eneste Løsning
er, at Ting (Genstande, Ikke-Jeg) kun ere til, for at
vi kunne have Noget at virke paa, Noget, der kan udløse
vore Kræfter. Den frie Brug af Kræfterne er det Højeste;
men den er kun mulig ved, at der er Skranker, Modstand at
overvinde. --- Det er da for Fichte en etisk Betragtning, der
giver Løsningen. Naturen er for ham Materiale for vor pligtmæssige
Handlen, vor Selvudvikling til sand Frihed.

\textsc{Fr}. W. \textsc{Schelling} (1775---1854) er Romantikens egentlige
Repræsentant i Filosofien. Uden nogen kritisk Betænkelighed
proklamerede han i sine Ungdomsskrifter, som er det
Betydeligste, han har frembragt, en ny Videnskab, der skulde
løse alle Gaader. Ogsaa han virkede en Tid i Jena, som var
Romantikens Samlingssted, senere ved forskellige bayerske
Universiteter. I sin Alderdom blev han af den romantiske
Kong Friedrich Wilhelm IV kaldt til Berlin for at modarbejde
de radikale Hegelianere; men hans Optræden her var en
Skuffelse af store Forventninger.

Schelling virkede paa Tiden ved det nittende Aarhundredes
Begyndelse gennem sin „Naturfilosofi“, hvorved han forstod
en aandfuld Skuen, for hvilken Naturen staar som ubevidst
Aand, der ad forskellige Trin gennem den uorganiske og den
organiske Verden nærmer sig til det Punkt, hvor Bevidstheden
vaagner til fremskridende Klarhed. Vi kunne kun erkende
Naturen, fordi den er selve Aanden paa et lavere Trin.
Naturen maa forstaas indefra, ikke, som den sædvanlige Naturvidenskab
forsøger, gennem udvortes Vexelvirkning mellem
de forskellige Stoffer og Kræfter. --- Det var dog ikke nogen
real Udvikling, Schelling antog i Naturen. De forskellige Trin
existere samtidigt; det er kun for Tanken, deres Slægtskab
viser sig. Og det var en  symbolsk Tydning, ikke en videnskabelig
Forklaring, der blev givet; Principet vera causa blev
skudt tilside. Den Begejstring, denne „Naturfilosofi“ vakte,
blev tildels af Betydning for Poesiens Udvikling (saaledes hos
% --- p. 66 ---
\setcounter{footnote}{0}%
\opage{66}os ved de Forelæsninger, Schellings Discipel holdt i København
1902); men for Tænkningen var det en Tilbagegang.

Ogsaa \textsc{G. W. F. Hegel} (1770---1831), Romantikens Systematiker,
virkede en Tid i Jena. Efter Jenaslaget drog han
til Bayern og kom senere til Berlineruniversitetet, hvor han
grundlagde en talrig og indflydelsesrig filosofisk Skole.

Hverken Fichtes etiske Postulat eller Schellings Skuen
tilfredsstillede Hegel. Han stillede sig den Opgave at udvikle
de Tanker, der udtrykke Tilværelsens inderste Natur, saaledes
at de danne en Række, i hvilken hvert Led med Nødvendighed
fører over til et andet. Ingen enkelt Tanke kan udtømme
Tilværelsen; enhver Tanke viser sig ved konsekvent Gennemførelse
at være begrænset, men fører netop ved at drives til
det Yderste over i en anden Tanke; forbindes nu disse to
Tanker, faa vi en højere Enhed, i hvilken de begge ere Led.
Men denne nye Tanke har selv sin Modsætning, sammen
med hvilken den danner en ny højere Enhed, og saaledes
fremdeles --- en fremskridende Proces, stedse gennem Treheder.
Dette er, hvad Hegel kalder den dialektiske Metode.
Men det er ikke blot i vor subjektive Tænken, Dialektiken
fremtræder. Der er ogsaa Dialektik i Tilværelsen. Udviklingen
skrider frem gennem Modsætninger og Overvindelse af Modsætninger.
Saaledes udfoldes den organiske Spire under Vexelvirkning
med Omverdenen, og i den fuldtudviklede Organisme
have vi den højere Enhed af hvad der var anlagt i Spiren,
og af Modsætningen til Omverdenen. Barndommens Umiddelbarhed
naar gennem Ungdommens Tvivl og Strid til Manddommens
Modenhed.

Selve Hegels System bestaar af tre Led. I den rene Logik
fremstilles Tankernes indre Forhold. I Naturfilosofien
træde ydre Forhold i Rum og Tid istedetfor de rene Tankeforhold.
Gennem denne Udvorteshed udvikler Tanken sig til
% PRINTED AS IS (read at the image, odd spots o6)
Aandsfilosofi, i hvilken Modsætningen mellem det Indre og
og det Ydre falder bort. I Kunst, Videnskab og Religion,
Aandslivets højeste Former, naar Tilværelsen sin Fuldendelse.
Og det er gennem alle disse Trin og Former den absolute
Verdensaand, der udfolder sig.

% --- p. 67 ---
\setcounter{footnote}{0}%
\opage{67}I Hegels Etik er Statens Ide Hovedtanken. Hans Etik er
en social Etik. Det Gode har sin Bestaaen i Samfundet, i
hvilket der rører sig Kræfter, som ligge ud over den Enkeltes
Bevidsthed. Sin højeste Udvikling naar Individet kun ved
at leve i og for Samfundet --- i Familien og det borgerlige
Samfund, men især i Staten, der optager begge i sig. Staten
kan ikke forklares ved Sammenslutning af Individer, men er
en højere Tingenes Orden, i Forhold til hvilken det enkelte
Individ er af forsvindende Værdi. Overfor alle individuelle
Idealer og Utopier fremhæver Hegel, at den etiske Ide ikke
er saa afmægtig, at den ikke skulde kunne realiseres; den
har fundet sin reale Form i den historisk udviklede Stat. Det
er ikke Filosofiens Sag at konstruere Idealer, men at paavise
Idealiteten i den historiske Udviklings Resultater. --- Paa
analog Maade forholder efter Hegel Filosofien sig til Religionen.
Der er kun en Formforskel mellem Filosofien og
Religionen; Indholdet er det samme. Filosofien lærer Enheden
i Tilværelsen gennem alle Modsætninger og paa alle
Trin; Religionen lærer, at Alt har sit Udspring fra Gud. Hin
taler Tankens, denne Billedets Sprog. Hvad Religionen forkynder
i store Billeder og paa konkret Maade, udtrykker
Filosofien i Tankens Form. Især i Treenighedslæren og Læren
om Mennesketilblivelsen fandt Hegel samme Sandheder
udtalte, som han dialektisk havde udviklet i sit System: Treenigheden
udtrykker, at det Højeste ikke kan være en abstrakt
Enhed, men en Enhed af Modsætninger; Mennesketilblivelsen
vidner om, at Guddommen netop gennem at gaa
over i sin Modsætning, i Fornedrelse og Lidelse, viser sig
som den højeste Verdensmagt. Det var især i Teologien, at
Hegels Filosofi blev modtagen som en Forbundsfælle; man
saa i den en Form, i hvilken det teologiske Læreindhold
kunde faa en ny Begrundelse. Som Hegel mente at have forsonet
alle Modsætninger, mente han ogsaa at have forsonet
Tro og Viden.

b. De romantiske Tankeretninger herskede i Tyskland og
Nabolandene til henimod Midten af Aarhundredet, og de gøre
sig endnu stadigt gældende i ændrede Former. Det er Inter%
% --- p. 68 ---
\setcounter{footnote}{0}%
\opage{68}essen for at hævde en idealistisk Opfattelse af Tilværelsen,
der endnu fører Mange til i dem at se et Højdepunkt i Filosofiens
Historie. Kants kritiske Tankegang forsvandt dog
ikke; i selve Romantikens Periode holdt den sig som en
Understrøm, repræsenteret af dygtige Forskere, som \textsc{Fries},
\textsc{Herbart} og \textsc{Beneke}, hvis Betydning især ligger paa Psykologiens
Omraade. Desuden var der Mænd, som vel vare
stærkt bevægede af Romantikens Tendenser, men som dog
hævdede en Erkendelsesteori i Kants Aand, medens de ad
forskellige Veje søgte en Livsanskuelse, der kunde føre ud
over Videnskabens Skranker. Saadanne Mænd vare, foruden
Teologen \textsc{Schleiermacher}, \textsc{Schopenhauer} i Tyskland og
\textsc{Kierkegaard} i Danmark. Endeligt var der Mænd, der ud
fra Hegels Skole, men i skarp Kontrast til Mesterens retsfilosofiske
og religionsfilosofiske Konservatisme udviklede radikale
Anskuelser. David Strauss hævdede, at man efter
Hegels Principer ikke kunde gøre Forskel paa Teologi og
Mytologi. \textsc{Karl} \textsc{Marx} og \textsc{Ferdinand} Lassalle udviklede
Hegels Statslære i socialistisk Retning. \textsc{Feuerbach} førte Religionsfilosofien
tilbage til ren Psykologi og antydede en filosofisk
Opfattelse, der er beslægtet med den i Frankrig opstaaede
Positivisme, som vi i næste Afsnit skulle fremdrage.

% PRINTED AS IS (read at the image, sweep u4)
Kun ved enkelte af disse Tænkeré kunne vi her dvæle:
Schopenhauer, Kierkegaard og Feuerbach.

α. \textsc{Arthur} \textsc{Schopenhauer} (1788---1860) var Søn af en rig
Købmand i Danzig, fik en alsidig Opdragelse og lærte tidligt
Verden at kende paa Rejser og ved at bevæge sig mellem
Mennesker. Uafhængig som han var, kunde han ganske hellige
sig til sine Studier og til Udviklingen af sin Livsanskuelse.
Efter et mislykket Forsøg som akademisk Lærer i Berlin
levede han som Privatmand i Frankfurt am Main. Meget tidligt
havde han gennem sin egen lidenskabelige Natur lært
Livets dunkle og disharmoniske Kræfter og Drifter at kende,
og hvad han saa omkring sig, vakte dels hans Harme, dels
hans Medlidenhed. Over disse Erfaringer, der lærte ham, at
ikke Forundring (som Platon mente), men Bestyrtelse og
% --- p. 69 ---
\setcounter{footnote}{0}%
\opage{69}Sorg var Begyndelsen til Filosofien, søgte han at hæve sig
ved Tankearbejde og ved kunstnerisk Skuen.

Schopenhauer protesterer imod, at Fichte, Schelling og
Hegel skulde være Kants rette Efterfølgere, og at den kritiske
Filosofi skulde være et tilbagelagt Stadium. I første Bog af
sit Hovedværk \textit{(Die Welt als Wille und Vorstellung,} første
Del 1819, andet Bind 1844) fremstiller han en Erkendelsesteori,
der slutter sig nær til Kants. Han er da klar over, at
Erkendelsen (die Vorstellung) ikke kan lære os Tilværelsens
Inderste (das Ding an sich) at kende; men han mener at
have fundet en anden Vej til Gaadens Løsning. Vi høre jo
selv med til Tilværelsen; denne rører sig i vort eget Indre.
Og nu rører der sig, langt dybere nede end al Erkendelse
(Vorstellung) en Tragten og Higen, en Selvopholdelsestrang,
der kundgør sig i al vor Lyst og Ulyst, i Haab og Frygt, i
Kærlighed og Had, --- en Villie, der udgør vort inderste Væsen,
Urfænomenet i os. Alle Forestillinger ere fra først af
denne Villies Tjenere. Hvad der er Urfænomenet i os, maa
--- mener Schopenhauer --- ogsaa være Urfænomenet i Tilværelsen
idethele; ved det Inderste i os forstaa vi det Inderste
i Verden. Saaledes gør han Overgangen fra Erkendelsesteori
til Metafysik ved Hjælp af en Analogi, hvis Berettigelse han
i sin romantiske Stemning ikke tvivler om. --- Overfor den
Vægt, Hegel lagde paa Tankelivet, var det af stor Betydning,
at Villieslivet blev fremhævet paa saa energisk, og i det Enkelte
ofte saa genial Maade. En Voluntarisme blev her stillet
overfor Intellektualismen. Allerede Fichte, af hvem Schopenhauer
er paavirket, havde jo iøvrigt givet Villien en fremherskende
Plads i sin Filosofi.

Spørges, hvad Villien gaar ud paa, svarer Schopenhauer:
paa at existere for Existensens egen Skyld; den er en Villie
til Livet (Wille zum Leben). Der gives her ikke mere noget
Hvorfor. Sætningen om den tilstrækkelige Grund gælder kun
for Erkendelsen (Vorstellung), ikke for Villien. Naturens mangfoldige
Former og Kræfter i deres indbyrdes Brydning og
Kamp, de stedse fornyede Bevægelsesprocesser og den stadige
Uro i Verden vidne om den blinde Selvopholdelsestrangs al%
% --- p. 70 ---
\setcounter{footnote}{0}%
\opage{70}drig hvilende Energi. Denne dunkle Trang hilder os i den
Illusion, at Livet er godt og værdifuldt. Men Livet staar kun
for os som godt, fordi Villien til Livet for enhver Pris raader
i os. I sig selv er Tilværelsen Smerte, netop fordi dens inderste
Væsen bestaar i en aldrig hvilende og aldrig mættet
Trang, og kun som Kontrastfænomen opstaar der Lyst eller
Tilfredsstillelse, nemlig naar den underjordiske Ild for en Tid
dæmpes. Al Lyst er Illusion, er et Nulpunkt, der kun paa
Grund af Kontrasten til den stadige Brynde og Uro kommer
til at staa som noget Positivt (Smlg. herom Psykologi VI E, 6).
I livlige Skildringer af Menneske- og Dyrelivet beskriver
Schopenhauer Eksistensens Kvaler, de levende Væseners indbyrdes
Kamp. Disse Kvaler mærker den store Hob ikke,
eller glemmer dem let igen. De dybere skuende, især de geniale
Mennesker derimod erfare Disharmonierne og sønderrive
Illusionens Slør.

Villien til Livet kan dog ophæves ad tre forskellige Veje.
Man kan stille sig overfor Livet som blot Tilskuer, idet man
ganske gaar op i æstetisk eller intellektuel Skuen. Naar vi
f. Ex fordybe os i Betragtningen af et Kunstværk, beroliges
vor Villie, og vi glemme, at vi høre med til Verden. Det er
snarest denne Vej, Schopenhauer selv er gaaet. Men ogsaa i
Menneskekærligheden (der paa Grund af Livets Kval maa
fremtræde som Medlidenhed) er Villien ophævet, idet jeg gør
mig ganske til Et med Andre. Paa denne Tanke har Schopenhauer
bygget sin Etik. Men den fuldkomne Villiesophævelse
lykkes dog kun for den Hellige, Asketen, for hvem der
ikke mere gives noget Motiv til Handlen. I Buddhismen, i
Urkristendommen og i Mystiken finder Schopenhauer den
højeste praktiske Løsning af Livsgaaden og Livskvalen, og
han staar med dyb Beundring overfor Askesens store Repræsentanter
--- saa meget desto mere, som han var sig bevidst,
ikke selv at være nogen Helgen.

β. I Danmark havde Hegels Filosofi vundet ikke ringe
Indflydelse, især hos Æstetikere og Teologer (J. L. \textsc{Heiberg},
og H. L. \textsc{Martensen}). F. C. \textsc{Sibbern} (1785---1872) gav en
træffende Kritik af den og fremsatte selv en Verdensanskuelse,
% --- p. 71 ---
\setcounter{footnote}{0}%
\opage{71}der lagde stor Vægt paa, at Udviklingen i Naturen og i Menneskelivet
udgaar fra spredte Punkter (sporadisk) og for en
stor Del gennem den Kamp og Brydning, der opstaar, naar
forskellige Udviklingstendenser mødes --- i Menneskets Indre
eller i Naturen og i Samfundet. Da ethvert Menneske selv er
et af de sporadiske Udgangspunkter, kan det ikke overse Tilværelsen;
deraf Debatten mellem forskellige Standpunkter i
Tænkningens Historie. I det Enkelte har Sibbern især beskæftiget
sig med Psykologi, hvortil hans Iagttagelsesevne og
hans friske Sympati med det menneskelige Liv gjorde ham
vel skikket.

Den religiøse Tænker \textsc{Søren} \textsc{Kierkegaard} (1813---1855)
fornyede den kritiske Filosofi og bestred Hegels Løsning af
Erkendelsesproblemet. Vi kunne have System, d. v. s. Orden
og Sammenhæng i vore Tanker, men vi kunne ikke sammenfatte
Tilværelsen i et afsluttet Tankesystem. Vi forstaa kun,
hvad der ligger færdigt og afsluttet for os; men vi leve midt
i Tiden, og vor Tanke kommer stedse efter Oplevelsen, idet
der stedse komme nye Erfaringer. Vi leve forlængs, men vi
forstaa baglængs. Medens vi tænke, gaar Livet videre, i os og
udenfor os. Og hvad der gælder videnskabelig Tænken, gælder
endnu mere den Tænken, der medvirker i vor personlige
Livsanskuelse. Her gælder det først og fremmest om, at Tankerne
ere vore egne, erhvervede eller tilegnede paa første
Haand; kun da svare de ganske til vor Personlighed, have
personlig Sandhed. Den, der med Inderlighed beder til en
Afgud, beder dog til den sande Gud, medens den, der i Ligegyldighed
og af Vane beder til den sande Gud, beder til en
Afgud. Personligheden (Subjektiviteten) er Sandheden. Kierkegaard
sympatiserer med Lessings Ord om den stadige Stræben
efter Sandhed. --- Det er i sit filosofiske Hovedskrift
Uvidenskabelig \textit{Efterskrift} (1846), at Kierkegaard udvikler sin
Opfattelse af Erkendelsen og opstiller Begrebet personlig Sandhed.
Han fremstiller tillige her sin sammenlignende Livsfilosofi,
Læren om de forskellige „Stadier“ eller Typer for personlig
Livsanskuelse og Livsførelse. Tre af disse, det æstetiske, det
etiske og det religiøse, har han i poetisk Form skildret i
% --- p. 72 ---
\setcounter{footnote}{0}%
\opage{72}\textit{Stadier paa Livets Vej} (1845) (i „Enten-Eller“ skildrer han
de to første af disse Stadier). Maalestokken for „Stadierne“
henter Kierkegaard fra Graden af de Modsætninger, det
aandelige Liv formaar at omspænde. Æstetikeren lever i Øjeblikkets
Nydelser; hos Etikeren er der en indre Kerne, han
trofast bevarer under Livets Skiften, og han finder sin Opgave
i Samfundet (Familie, Stat); den Religiøse lever i Bevidstheden
om den skærende Modsætning mellem Evigheden
og Livet i Tiden. Overgangen fra et af disse Stadier til et
andet kan ifølge Kierkegaard ikke ske ved kontinuerlig Udvikling,
men kun ved et Spring, der ikke kan psykologisk
motiveres. (Smlgn. om Overgangen fra det Æstetiske til det
Etiske Psykologi VI E, 4). --- Mere og mere blev for Kierkegaard
ikke selve Syntesens Energi, men den Lidelse, den
stærke Spænding mellem Modsætningerne voldte, den egentlige
Maalestok, og han kom derved i Modsætning til Romantikens
Proklamation af Harmonien af alle Modsætninger, saavel
som til den sædvanlige Opfattelse af Kristendommen.
Derved fremkaldtes den voldsomme Kamp mod Folkekirken,
der opfyldte hans sidste Aar, og som han tilsidst førte gennem
Flyvebladet \textit{Øjeblikket} (1855).

Kierkegaard levede som Privatmand, optagen af den store
filosofiske, æstetiske og religiøse Produktion, i hvilken hans
skarpe Reflexion og rige Stemningsliv, hans Tungsind og hans
Lidenskabelighed fik Afløb. Da han i „Øjeblikket“ havde faaet
sagt sit sidste Ord, bukkede han under for Sygdom.

γ. \textsc{Ludwig} \textsc{Feuerbach} (1804---1872) var under Hegels Indflydelse
gaaet over fra Teologi til Filosofi. Men han kunde
ikke blive staaende ved Hegels Religionsfilosofi, der søgte at
finde filosofiske Ideer i de religiøse Forestillinger. Han stillede
sig derimod den Opgave at finde den psykologiske Kilde til
disse Forestillinger, at give en psykologisk Forklaring af Dogmernes
Opstaaen. Han genoptager Hume's Religionspsykologi.
Hans Hovedværk er \textit{Das Wesen des Christenthums} (1841).
--- Feuerbach levede, efterat have forsøgt en Virksomhed ved
Universitetet i Erlangen, som Privatmand i landlig Ensomhed,
% --- p. 73 ---
\setcounter{footnote}{0}%
\opage{73}beskæftiget med Studier og Forfattervirksomhed. I sine sidste
Aar maatte han kæmpe med Fattigdom og Sygdom.

Religionerne opstaa ifølge Feuerbach ved, at den inderlige
Tragten efter Opfyldelsen af Hjertets Ønsker sprænger Fornuftens
Skranker. Derved forklares, at Religionen paa sit
Højdepunkt faar en antirationel Karakter. Ønsket er det teogoniske
(gudeskabende) Princip. De religiøse Prædikater ere
menneskelige Ønskers Indhold; Menneskenes Forestillinger
om Himmelens Beskaffenhed og om Gudernes Egenskaber
ere Udtryk for, hvad der bevæger sig i deres Hjerters Inderste.
At Gud er personlig, betyder, at det personlige Liv er af
højeste Værdi (er guddommeligt); at Gud er Kærlighed, betyder,
at Kærlighed er af højeste Værdi (er guddommelig).
Vil man forstaa Religionen, maa man gøre de religiøse Prædikater
til Subjekter.

Fra Begyndelsen af har Mennesket ingen Grund til at
sætte Grænser for sine Ønsker og for de Forestillinger, Ønskerne
fremkalde. Og det ligger i Følelsens Natur at uendeliggøre
sin Genstand og dog fastholde dens Virkelighed. Tvivlen
opstaar først, naar Mennesket lærer sine Skranker at kende,
og naar det begynder at skelne mellem sine Forestillinger og
Virkeligheden.

Mellem Religion og Kultur bestaar der et omvendt Forhold.
De ere to forskellige Veje, ad hvilke Mennesket haaber
at naa sine Formaal, og jo mere Tillid det har til den ene,
des mindre Brug har det for den anden. Paa lignende Maade
forholder det sig med Religion og Etik. Jo mere Forskellen
mellem Gud og Menneske indskærpes, des mere maa Egenskaber
som Kærlighed og Retfærdighed tages i en ganske
anden Betydning, naar de anvendes om Gud, og naar de anvendes
om Mennesker, og der kan komme en skarp Strid
mellem Guds Villie og det, menneskelig Kærlighed og Retfærdighed
kræver.

Ingen virkelige Værdier forsvinde med Religionen. De
Egenskaber, som ere værdifulde, kunne vedvarende hævdes,
selvom de ikke henføres til Guder. Feuerbach grunder Etiken
paa Menneskekærligheden, i hvilken det enkelte Individs En%
% --- p. 74 ---
\setcounter{footnote}{0}%
\opage{74}hed med hele Slægten kundgør sig. Han lægger Vægt paa, at
den Enkelte maa kende Lykke og Ulykke af egen Erfaring
for at kunne have Medfølelse med andre, og hævder derfor
Lykketrangens Berettigelse. --- I sin almindelige Opfattelse af
Filosofien (udtalt i \textit{Grundsätze der Philosophie der Zukunft}
1843) finder han dens Genstand i Mennesket som Erfaringsvæsen
og i Naturen som Grundlag for dets Eksistens. Overfor
Materialismen fremhæver han (ligesom Schopenhauer), at
Materien er et af Mennesket dannet Begreb, og at derfor den
Opfattelse, der betragter Mennesket som Produkt af Materien
ikke kan være afsluttende.

\closerule

\chapterhead{VII.}{POSITIVISMEN.}

De to store Tankeretninger i det nittende Aarhundrede
ere Romantiken og Positivismen. Hin gaar ud fra Tankens
Former og Idealer, denne fra det faktisk Givne; „positiv“ betyder
fremfor Alt „virkelig, sat, given.“ Saa forskellige, ja modsatte
de ere, betyde de dog begge, hver paa sin Maade, en
Modsætning til Oplysningens, Kritikens og Revolutionens Aarhundrede.
Positivisterne ville, ligesom paa deres Vis Romantikerne,
fordybe sig i Naturen og i Historien. Begge Retninger
udsprang ved Aarhundredets Begyndelse, og Positivismen
er ikke at opfatte som en Reaktion mod Romantiken, skønt
den først kom i Forgrunden, da Romantikens Herredømme
var brudt. --- Medens Romantikens Filosofi havde sit Hovedsæde
i Tyskland, hørte Positivismen væsentligt hjemme i
Frankrig og England.

1. Positivismen har sit Udspring i den af \textsc{Saint Simon}
(1760---1825) grundede sociale Skole. Hans Opgave var at
forberede en social Reformation, der efter hans Opfattelse
% --- p. 75 ---
\setcounter{footnote}{0}%
\opage{75}kun var mulig, naar man grundede en ny Verdensanskuelse,
der for Nutiden kunde være det Samme, som Kristendommen
havde været for Middelalderen. En saadan ny Verdensanskuelse
maatte bygges paa de positive Videnskabers Grundlag,
ligesom Menneskene ogsaa kun paa dette Grundlag kunde
udbytte Naturen i fælles Arbejde, istedetfor at de forhen
havde udbyttet hverandre. Hertil udfordredes, at de teologiske
Forudsætninger, der forhen lagdes til Grund i Videnskaben,
udskiltes, saa at der byggedes paa rent naturlige Forudsætninger.
Naar dette var sket, kunde den positive Filosofi grundlægges.

\textsc{Auguste Comte} (1798---1857), Positivismens egentlige
Grundlægger, studerede ved den polytekniske Skole i Paris;
men da denne paa Grund af Elevernes revolutionære Sindelag
blev lukket af Bourbonerne, fortsatte han sine Studier
paa egen Haand og gav dem en mere encyklopædisk Karakter,
hvortil ogsaa hans Samarbejde med Saint-Simon bidrog.
Dette Samarbejde hørte op, fordi Saint-Simon efter Comte's
Mening for meget underordnede Videnskaben under sine Ideer
om den sociale Reformation og Fremtidsreligionen. Med stor
Arbejdskraft og Koncentrationsevne gennemførte Comte sin
Fremstilling af den positive Filosofi som en stor Encyklopædi
af de positive Videnskaber \textit{(Cours de Philosophie positive.}
1830---1842). I hans senere Aar antog hans Tænken en mere
subjektiv og mystisk Karakter, idet han betragtede sig som
Grundlægger af en Humanitetsreligion og endog indrettede en
Art Kultus.

Ifølge Comte gennemløber den menneskelige Viden (idethele
og paa de enkelte Omraader) tre Stadier. Paa det teologiske
Stadium raades der kun over en lille Kreds af Erfaringer,
og Indbildningskraften spiller derfor en stor Rolle. Det
Baand, der her forener Kendsgerningerne til et Hele, er Forestillingen
om Aander og Guder. Teologiens kulturhistoriske
Betydning beror paa, at den paa dette Udviklingstrin har tjent
Videnskaben som intellektuelt Baand, ligesom den ogsaa i de
religiøse Autoriteter ydede et Grundlag for Moralen. Paa det
metafysiske Stadium findes Forklaringen ikke mere i person%
% --- p. 76 ---
\setcounter{footnote}{0}%
\opage{76}lige Væsener, men i almindelige Kræfter eller Ideer. Man taler
om Tiltrækningskraft, kemisk Kraft, Livskraft o. s. v. Dette
er et Fremskridt, da saadanne Kræfter eller Ideer antages at
virke paa regelmæssig Maade. Men dette Stadium har dog
for Comte en overvejende negativ Karakter. Det er Individualismens
og Revolutionens Stadium. Paa det positive Stadium
underordnes Indbildningskraft og Spekulation under Erfaringen.
Nu søges Fænomenernes Love, d. v. s. de konstante
Forhold mellem dem. De forskellige Love kunne ikke føres
tilbage til en eneste; Videnskabens Enhed bestaar blot i, at
det er den samme Hovedmetode, der anvendes paa alle Omraader.
Denne Enhed i Metoden gør et aandeligt Fællesskab
muligt, som ikke havde været kendt siden Middelalderen.

Hovedmetoden fremtræder dog paa noget forskellig Maade
i de forskellige Videnskaber. Disse kunne ordnes i en Række,
saaledes at der begyndes med dem, der angaa de simpleste
Emner og opstille de mest omfattende Love; dette er tillige
dem, der hurtigst ere naaede til det positive Stadium, og i
hvilke Deduktionen kan anvendes i stor Udstrækning. Jo mere
vi skride frem i Rækken, des mere indviklede blive Emnerne,
des mere specielle blive Lovene, des mere maa Induktion
anvendes, og des senere naas det positive Stadium. Rækkefølgen
er: Matematik, Astronomi, Fysik, Kemi, Biologi og Sociologi
(Samfundslære). Der gives ingen kontinuerlig Overgang
fra den ene Videnskab til den anden. --- Vi savne i denne
Klassifikation Logiken, der vilde have sin Plads foran Matematiken,
og Psykologien, der hos Comte er delt mellem Biologien
og Sociologien.

Comte skelner mellem de positive Kendsgerninger og deres
Forbindelse. Denne Forbindelse sker paa det positive
Stadium dels ved Ordning efter Lighed (par similitude), dels
ved Udledning af det Følgende fra det Foregaaende (par
filiation), men i begge Tilfælde gælder det en Sammenfatten
(tout se reduit toujours à lier). (\textit{Discours sur l'esprit positif.}
1844). Comte regnede da ogsaa Kant blandt sine Forgængere.
--- I sin Etik anlægger han et lignende Synspunkt som Grotius
og Hume. Pligtbegrebet udspringer af den Enkeltes Be%
% --- p. 77 ---
\setcounter{footnote}{0}%
\opage{77}vidsthed om at være Led af et socialt Hele, en Bevidsthed,
der opstaar, naar alt Menneskelivs Solidaritet erfares og anerkendes.
Uvilkaarligt følger Mennesket Trangen til Samkvem,
og Nytten af dette gaar først bagefter op for det. Spiren til
Solidaritetsfølelse ligger allerede i Forholdet mellem de to
Køn og i Omsorgen for Afkommet; i den menneskelige Verden
foregaar der saa en stadig Opdragelse til Altruisme (et
af Comte dannet Ord til at betegne det Modsatte af Egoisme).
Det enkelte og isolerede Individ er kun en Abstraktion. ---

2. I videre Forstand er ogsaa Hovedretningen i det nittende
% PRINTED AS IS (read at the image, structure pass s1/s2)
Aarhundres engelske Filosofi positivistisk. Ved Midten
af Aarhundredet indtog \textsc{John} \textsc{Stuart} \textsc{Mill} (1806---1873) en
højt anset Plads i den videnskabelige Verden ved sine Arbejder
over Logik \textit{(System of Logic} 1843), sin Nationaløkonomi
og sine etiske og socialreformatoriske Arbejder. Han
var stærkt paavirket af den radikale Etiker og Samfundsreformator
\textsc{Jeremy} \textsc{Bentham}, men tillige af Comte, saavelsom af
den romantiske Retning, der i England repræsenteredes af
Coleridge og Carlyle. Mill beklædte et Embede i det ostindiske
Kompagnis Tjeneste og anvendte sin Fritid til sine Studier.
I sin interessante Selvbiografi har han gjort Rede for sin
Udviklingsgang, i hvilken saa mange forskellige Retninger
greb ind, uden at han ret formaaede at sammenarbejde dem.

I sin Logik stiller han sig i skarpere Modsætning til Kant,
end Comte gjorde. Han forsøger at opfatte al Tænkning som
blot Association, særlig som Berøringsassociation. Grundformen
for al Slutning finder han i den Kendsgerning, at naar to
Fænomener hyppigt ere forekomne sammen, vente vi det andet,
naar det ene er indtraadt. Det er en Slutning fra Enkelthed
til Enkelthed. Ere mange saadanne Erfaringer gjorte, formulere
vi dem til almindelige Sætninger, der kunne danne
Grundlag for Deduktion. Men al Deduktion forudsætter en
Induktion, tilsidst en Overgang fra Enkelthed til Enkelthed.
Ogsaa Aarsagssætningen grunder sig paa Induktion. --- Mill's
Logik er Modstykket til Hegels. Medens Hegel vilde udlede
alt Tankeindhold af Tankens Former, vilde Mill udlede Tankens
Former af det givne Indhold. Her fremtræder Modsæt%
% --- p. 78 ---
\setcounter{footnote}{0}%
\opage{78}ningen mellem Romantik og Positivisme paa den skarpeste
Maade. --- Betydningen af Mill's Logik ligger i hans Karakteristik
af de specielle Metoder og i den redelige Drøftelse af
de principielle Spørgsmaal.

Som Etiker hævdede Mill i Tilslutning til Bentham „Nyttemoralen“
(Utilitarianismen), der vurderer Handlingerne efter
den Grad, i hvilken de virke til den størst mulige Lykke for
saa stort et Antal Mennesker som muligt (et allerede af Bentham
og tidligere Filosofer opstillet Princip). Den etiske Følelse,
der driver til at handle efter dette Princip, udfolder sig
under Indflydelse af forskellige Motiver, blandt hvilke Sympatien
indtager en fremragende Plads.

3. En stor Udvidelse fik Positivismen ved Udviklingshypotesen,
der fremtraadte kort efter Aarhundredets Midte. Allerede
før Darwins Offentliggørelse af sin Teori havde Herbert
Spencer hævdet, at man ved de Erfaringer, de menneskelige
Tanker og Følelser skylde deres Udvikling, ikke blot maatte
tænke paa det enkelte Individs Erfaringer, man maatte tage
Hensyn til Slægten, hvorfra det stammer. Men det var især
ved Darwins Forbindelse af Udviklingslæren med Ideen om
Kampen for Livet, at dette Synspunkt vandt mere almindelig
Opmærksomhed.

\textsc{Charles} \textsc{Darwin} (1809---1881) faar med Rette en Plads
i Filosofiens Historie, dels fordi hans Resultater (ligesom
Kopernikus', Galilei's og Newton's) ere af saa stor Betydning
for Verdensanskuelsen, dels fordi hans Opfattelse af Videnskaben
og dens Grænser har filosofisk Interesse. Darwin selv
levede som stille Forsker i landlig Ensomhed, efterat han i
sin Ungdom paa en treaarig Rejse om Jorden havde indsamlet
et stort Forraad af Iagttagelser. --- Hans Forklaring af
Arternes Oprindelse \textit{(Origin of Species} 1859) er ganske i Positivismens
Aand. At antage, at hver Art er skabt af Gud, er
for Darwin kun at sige, at man ikke véd, hvorledes den er
bleven til. Derimod støtter Darwin sig til, at Organismerne
variere, og at disse Variationer kunne være mere eller mindre
gavnlige for Organismen, naar den skal skaffe sig Føde,
Lys og Varme og sikre sit Afkom og sig selv mod Farer og
% --- p. 79 ---
\setcounter{footnote}{0}%
\opage{79}Fjender, eller som Darwin med et billedligt Udtryk kalder
det, naar den skal bestaa i Kampen for Tilværelsen. De Variationer,
som vare afgjort uheldige i denne Kamp, vilde efterhaanden
uddø; kun de mere eller mindre hensigtsmæssige
kunne bestaa. Denne Lære havde for Darwin den Betydning
at tjene til en Tanketraad (the intelligible thread), ved hvilken
en stor Mængde Kendsgerninger kunde forbindes. En Grænse
% PRINTED AS IS (read at the image, structure pass s1/s2)
for sin Hypetese saa han i Spørgsmaalet om Oprindelsen til
de Variationer, mellem hvilke Kampen for Tilværelsen (eller,
% PRINTED AS IS (read at the image, odd spots o6)
som Darwin ogsaa kalder den, den naturlige Udvælgelse) staar,
En lignende Gaade finder han i Livets Oprindelse overhovedet.
--- Senere Forskere have hævdet, at kun ved betydelige,
„springende“ Variationer (Mutationer, som deres Opdager
Hugo de Vries har kaldt dem) blive nye Arter til, ved hvilke
det saa bliver et Spørgsmaal, om de kunne bestaa i Kampen
for Tilværelsen; denne er altsaa ikke egentlig Aarsag til nye
Arters Opstaaen.

Af Darwins Selvbiografi og Breve \textit{(Life and Letters of
Charles Darwin} 1887) ser man, at hans sidste Stade med
Hensyn til Forholdet mellem Teleologi og Mekanisme var et
lignende som Kants. Ligesaa lidt som han kunde tro, at Verden,
som vi kende den, var Værk af en bevidst Hensigt
(design), ligesaa lidt kunde han blive staaende ved, at den
skulde skyldes Tilfældet eller den blinde Kraft (chance, brute
force). Han erklærede sig selv for en Agnostiker (d. v. s. En,
der erkender sin Uvidenhed angaaende det, der ligger ud
over Erfaringen).

\textsc{Herbert} \textsc{Spencer} (1820---1903) gik fra det praktiske Liv
over til filosofisk Forsken. Han begyndte som Ingeniør, fik
derefter Interesse for sociale Spørgsmaal og blev derved ført
til psykologiske og filosofiske Studier. Hans Filosofi udsprang
af Trangen til at finde en Grundtanke, der kunde sammenfatte
en Række af naturvidenskabelige, psykologiske og sociologiske
Studier. I sine \textit{Essays} har han givet en Række, tildels
glimrende skrevne Afhandlinger om enkelte Spørgsmaal.
Senere udgav han i ti Bind sin „syntetiske Filosofi“, hvis
% --- p. 80 ---
\setcounter{footnote}{0}%
\opage{80}første Bind \textit{(First Principles} 1861) giver Grundlaget for hans
Verdensanskuelse: Udviklingsbegrebet.

I sin Erkendelseslære viser Spencer, under Indflydelse af
den skotske Kantianer William Hamilton, at kun det, der
staar i Forhold til noget Andet, kan erkendes. Et Ubetinget
og Altomfattende kan derfor ikke erkendes. Enhver Erkendelse
giver kun et Udsnit af en uendelig Realitet, som kun
kan udtrykkes ved Symbolers Hjælp, ikke i Begrebets Form.
Spencer er, med denne Motivering, Agnostiker ligesom Darwin.
--- Men al den Erkendelse, vi kunne vinde, samler sig
om Udviklingsbegrebet. Vi have kun forstaaet et Fænomen,
naar vi have fulgt det gennem dets forskellige Trin til dets
nuværende Skikkelse. Det viser sig, at der her er visse Træk,
der stadigt komme igen, ligegyldigt hvilke Fænomener vi
have for os, og som tilsammen tjene til at definere, hvad vi
forstaa ved Ordet Udvikling. Der er tre saadanne Hovedejendommeligheder.

1) Udvikling er Koncentration. Der sker en Overgang fra
løsere til fastere Sammenhæng, naar et Fænomen bliver til.
Saaledes naar en Sanddynge danner sig paa Strandbredden,
ved Stjernesystemers og Kloders Tilbliven, ved Bevidsthedslivets
og Samfundslivets Udvikling. --- 2) Men dernæst er
Udvikling Differentiation, Overgang fra en mere ensartet til
en mere forskelligartet Tilstand. Medens Koncentration gaar
for sig ved, at ensartede Dele paavirkes af samme Kraft
(Sandskornene f. Ex. af Vinden), sker Differentiationen ved,
at de ensartede Dele indenfor den ved Koncentrationen dannede
Helhed paavirkes af forskellige Kræfter (som naar en
Del ligger i Skygge, en anden i Sol). Indenfor Helheden bliver
der da Forskellighed mellem Delene. Indenfor Solsystemet
udforme sig de enkelte Kloder, og paa samme Klode bliver
der Forskel mellem de indre og de ydre Dele, og mellem de
forskellige Dele af Overfladen. I Organismen fremtræder der
større Forskelligheder mellem Organerne, jo mere Væxten
skrider frem, og det organiske Liv paa Jorden fremtræder paa
senere Trin med flere forskellige Arter end paa tidligere Trin.
Ligeledes gaar det med Bevidsthedslivet, hvor Forskellig%
% --- p. 81 ---
\setcounter{footnote}{0}%
\opage{81}hederne paa senere Trin gøre sig stærkest gældende, og i Samfundet
er Arbejdets Deling et Udtryk for den stigende Differentiation.
--- 3) Naar Koncentrationen ikke falder bort ved
Differentiationen, men den faste Sammenhæng bevares trods
Forskellighedernes Opstaaen, gaas der over fra en mere ubestemt
til en bestemt ordnet Tilstand. Her fremtræder Udviklingen
som Determination.

Denne tredie Udviklingsform minder om, hvad Hegel
kaldte den højere Enhed, i hvilken Modsætninger forenedes
eller gik op. Men medens Hegel lærte, at saadanne Enheder
dannes med logisk Nødvendighed, kan Spencer ikke garantere
en saadan Fremskriden. Og ligesaalidt som andre Begreber
har Udviklingsbegrebet Gyldighed ud over Erfaringen.

I sin Etik \textit{(Data of Ethics.} 1879) bygger Spencer (ligesom
Hume og Comte, iøvrigt ogsaa Darwin) paa Sympatien som
Grundlag og søger at vise, hvorledes der ud fra den kan begrundes
en Livsførelse, i hvilken den menneskelige Udvikling
naar sin fuldkomneste Form: som Koncentration viser det
etiske Liv sig i Selvbeherskelsen, som Differentiation i den
rige Udfoldelse af Evnerne, Medfølelsen med Andre giver
Anledning til, og som Determination i Samvittighedsfuldhed
og Retfærdighed.

\closerule

\chapterheadone{VIII. NYERE PROBLEMBEHANDLINGER.}{Nyere Problembehandlinger}

I den romantiske Idealisme og den positivistiske Realisme
staa to Tankeretninger overfor hinanden, der have det fælles,
at de enkelte Problemer overvejende drøftes ud fra almindelige
Synspunkter. I en vis Modsætning til begge Retninger
træde Bestræbelser, der indgaaende dvæle ved enkelte Problemer,
maaske ved et eneste, og lade de almindelige Synspunkter
træde i Baggrunden eller komme som et sidste Ud%
% --- p. 82 ---
\setcounter{footnote}{0}%
\opage{82}blik. Idet vi ordne Fremstillingen efter den Karakteristik af
de fire Hovedproblemer, der er given i Indledningen, give vi
i det Følgende en Oversigt over saadanne Bestræbelser i den
sidste Del af det nittende Aarhundrede og i Begyndelsen af
det tyvende.

\sectionhead{a. Det psykologiske Problem.}

α. Af stor Betydning blev her Grundlæggelsen af en experimental
Psykologi ved \textsc{Gustav Theodor Fechner} (1801---
1887), der fra Fysiken gik over til Filosofien. I sin Verdensanskuelse
staar han Romantiken nær, idet han dristigt anvender
poetiske Analogier; men han forener med denne Idealisme
en exakt naturvidenskabelig Synsmaade. Da Loven om Energiens
Bestaaen efter hans Opfattelse gør det umuligt at antage
et Aarsagsforhold mellem det Aandelige og det Materielle, antager
han, at Forholdet er som mellem den indre og den
ydre Side af en Kugleflade: Idealisten ser kun hin, Materialisten
kun denne, og dog se de egentligt det Samme. Men
Fechner stiller sig nu den Opgave at udfinde, hvorledes disse
to Synspunkter forholde sig til hinanden, og det er ad denne
Vej han er kommen til experimental Psykologi \textit{(Elemente der
Psychophysik.} 1861). Ligesom Kepler, med hvem han er
aandsbeslægtet, arbejdede han sig fra poetisk Spekulation frem
til exakte Undersøgelser. Ved Psykofysik forstaas Undersøgelsen
af Forholdet mellem det Aandelige og det Materielle. Ved
at undersøge Forholdet mellem de fysiske Paavirkninger og
de tilsvarende Sansefornemmelser kom Fechner til det Resultat,
at Paavirkning og Fornemmelse ikke ere ligefrem proportionale,
men at jo stærkere en Paavirkning er, des mere
maa den forøges, hvis en vis Ændring af Fornemmelsen skal opstaa.
Paavirkningerne maa voxe i en Kvotientrække, for at
Fornemmelserne kunne voxe i en Differensrække (Se Psykologi
V A, 3). Ved egne og Andres Forsøg fandt han denne
Lov bekræftet. Dog er den (saavidt den gælder) neppe, som
Fechner mente, en Lov for Forholdet mellem Sjæleliv og
Hjerneliv. Disse ere sandsynligvis (hvad der ogsaa stemmer
% --- p. 83 ---
\setcounter{footnote}{0}%
\opage{83}bedst med Fechners egen Identitetshypotese) ligefrem proportionale,
hvorimod Fechners Lov udtrykker Forholdet mellem
Sjæl eller Hjerne paa den ene Side, de fysiske Paavirkninger
paa den anden Side.

Hvad Fechner havde grundlagt, førte \textsc{Wilhelm} \textsc{Wundt}
(f. 1832), der fra Fysiologien gik over til Filosofien, videre og
organiserede den nye Videnskab ved sit Værk \textit{Physiologische
Psychologie} (1874). En hel Række andre Problemer foruden
det af Fechner behandlede optoges nu, saaledes den Hastighed,
hvormed psykiske Fænomener opstaa, Rumsopfattelsens
Udvikling, Forestillingsassociationens Love o. s. v. --- En
Række af Forskere have fortsat, hvad Fechner havde grundlagt
og Wundt organiseret. I Danmark har \textsc{Alfred} \textsc{Lehmann}
arbejdet paa dette Omraade.

I Frankrig har man især søgt at belyse Psykologien ved
Studiet af det syge Sjæleliv, hvis Forhold og Ytringer ofte
lade Sjælelivets Ejendommelighed fremtræde paa særligt
slaaende Maade. Den berømte Nervepatolog \textsc{Charcot} og
hans Skole fik i denne Henseende stor Indflydelse; i den
nyeste Tid virker \textsc{Pierre} \textsc{Janet}, der fra Filosofien er gaaet
over til Psychiatrien, i denne Retning.

β. Den beskrivende Psykologi fandt betydelige Dyrkere i
\textsc{William} \textsc{James} (f. 1842, d. 1910) \textit{(Principles of Psychology}
1890) og \textsc{Henri} \textsc{Bergson} (f. 1859) (\textit{Les données immédiates de
% PRINTED AS IS (read at the image, sweep u4)
la concience.} 1888), der begge lægge stor Vægt paa det uvilkaarlige
Sjæleliv i dets kontinuerlige Strøm, i Modsætning til
de Forskelle og Modsætninger, praktiske og teoretiske Fornødenheder
foranledige os til at opstille. Allerede i en Afhandling
fra 1884 havde James paavist den oprindelige sjælelige
Livsstrøm (the original flux of life, the stream of thought)
som det Grundlag, hvoraf vor udtrykkelige Bevidsthed stadigt
% PRINTED AS IS (read at the image, structure pass s1/s2)
øser. Bergson fremhæver særligt, at alle kvantitive Forskelle
kun opstilles som Midler til klar og tydelig Opfattelse og
Forstaaelse af det umiddelbart Givne, hvis kontinuerlige Forløb
kun frembyder kvalitative Forskelle. Med stor Fremstillingskunst
forsøge de to Psykologer at aabne vort Blik for
det oprindeligt og umiddelbart Givne, der allerede i den dag%
% --- p. 84 ---
\setcounter{footnote}{0}%
\opage{84}lige Sprogbrug og senere i Videnskaben udstykkes paa mekanisk
Maade.

γ. Med Hensyn til Bevidsthedslivets ejendommelige Natur
gøres mere og mere den Opfattelse gældende, som Leibniz og
Kant havde udtalt, at vi maa søge denne Ejendommelighed
deri, at al sjælelig Livsytring kundgør sig som en Sammenfatten,
en Syntese af et mere eller mindre mangfoldigt Indhold,
hvis Elementer faa deres Karakter bestemt ved selve
denne Sammenfatten. Denne Hypotese om Bevidsthedslivets
Natur træder i Modsætning dels til den skolastiske Lære om
en Sjælesubstans (der hyldedes af Descartes og hans spiritualistiske
Efterfølgere), dels til den atomistiske Opfattelse af
Sjælelivet, der hyldedes af den engelske „Associationspsykologi“
(Hume, Stuart Mill). Til den slutter Wilhelm Wundt sig;
den søges begrundet i den Fremstilling af Psykologien,
som Forfatteren af denne Oversigt har givet, ligesom den
ogsaa gøres gældende af Pierre Janet paa Grundlag af hans
Undersøgelser af det syge Sjæleliv. (Smlgn. Psykologi II, 5;
V B, 5). (I sin Bog \textit{La synthèse mentale,} 1908, har \textsc{Dwelshauer}
fremstillet Syntesebegrebet, saaledes som det fremtræder
hos de tre nævnte Forfattere og hos Leibniz og Kant).

Syntesebegrebet fører naturligt til at betone Sjælelivet som
en ejendommelig Energi og derved ogsaa til at fremhæve
Villieslivet som Hovedsiden af Bevidsthedslivet. Paa dette
Punkt har Schopenhauer øvet en gennemgribende Indflydelse.
Der gør sig en Voluntarisme gældende i den psykologiske
Teori, i Modsætning til Intellektualisme og Emotionalisme. I
\textsc{Alfred} \textsc{Fouillée}'s \textit{Psychologie des idée-forces} (1893) har
denne Synsmaade faaet sit mest interessante Udtryk. Alle
sjælelige Fænomener ere Symptomer paa en Trang eller
Stræben (appétition), der fremtræder som Lyst eller Ulyst,
alt efter som den begunstiges eller hæmmes. Skelnen og
Foretrækken (discernement et préférence) ere oprindeligt Et,
og enhver Sansning, enhver Tankeakt betegner en mere eller
mindre bevidst og bestemt Retning af vor Livstrang.

% --- p. 85 ---
\setcounter{footnote}{0}%

\opage{85}\sectionhead{b. Tilværelsesproblemet.}

Romantikens Filosofi havde trot at kunne reformere Naturvidenskaben.
Og dog havde netop samtidigt med denne
Filosofis Tilbliven, i Slutningen af det attende Aarhundrede,
Naturvidenskaben gjort store Fremskridt. Den gamle Overbevisning
om Materiens Bestaaen blev bekræftet af \textsc{Lavoisier}
ved Vejning, altsaa formedelst den kvantitative Metode, og en
Række af Forskere (Priestley, Saussure o. fl.) paaviste Hovedlovene
for Stofskiftet i Plante- og Dyreverdenen og indordnede
derved det organiske Liv i Naturens store Kredsløb.

Et nyt Opsving fik Naturvidenskaben omtrent ved Aar
1840, navnligt ved \textsc{Robert} \textsc{Mayers} Opdagelse af Energiens
Bestaaen (1842). Tanker, som allerede Descartes, Huyghens
og Leibniz havde udtalt, fik nu erfaringsmæssig Bekræftelse,
idet det viste sig, at det ikke blot er paa det rent mekaniske
Omraade, at ingen Kraft forgaar eller opstaar, men at dette
ogsaa gælder under Vexelvirkningen mellem de specielle Naturkræfter
indbyrdes, idet der kan godtgøres et bestemt kvantitativt
Forhold (Ækvivalens) mellem den Kraftytring (f. Ex.
Bevægelse), der forsvinder, og den (f. Ex. Varme), der opstaar.

Dertil kom endnu (1859) Darwin's Hypotese om Arternes
Opstaaen. Naturvidenskaben havde saaledes paa storartet
Maade paavist indre Sammenhæng, hvor man forhen kun
havde set Brudstykker og Spring. Spørgsmaalet var nu, hvilken
Indflydelse disse Resultater vilde faa paa den almindelige
Verdensanskuelse. Hos Herbert Spencer have vi allerede
truffet et Forsøg paa at besvare dette Spørgsmaal. Ligeledes
hos Fechner, der fornyede Spinozas Identitetshypotese. I en
glimrende Fremstilling søgte \textsc{Albert Lange} (1828---1875)
\textit{(Geschichte des Materialismus} 1865) at godtgøre Umuligheden
og Uberettigelsen af at indskyde psykiske Led i de materielle
Aarsagers Rækker, hvorfor Identitetshypotesen bliver den
eneste Mulighed. Materialismen er uholdbar, fordi de materielle
Processer, af hvilke den vil udlede de sjælelige Fænomener,
selv kun kendes af os gennem Fornemmelser og Fore%
% --- p. 86 ---
\setcounter{footnote}{0}%
\opage{86}stillinger, altsaa gennem sjælelige Processer. Materien er et
Begreb, der er dannet ifølge de Love, der gælde for vor
Sansning og Tænkning. Saaledes kombinerer Lange Tilværelsesproblemet
og Erkendelsesproblemet.

Men foruden den spinozistiske Hypotese (som i det Væsentlige
hyldes af Wundt og ogsaa søges begrundet i den
Fremstilling af Psykologien, Forfatteren af denne Oversigt har
givet) fremtræde ogsaa Forsøg i materialistisk og i spiritualistisk
Retning.

I den tyske Literatur gav Naturvidenskabens Opsving Anledning
til en Række materialistiske Skrifter, der havde den
Fortjeneste at udbrede naturvidenskabelige Kundskaber og
Ideer i større Kredse. Med Oppositionen mod Dogmatik og
spiritualistisk Spekulation forbandt sig hos de tyske Materialister
i det nittende Aarhundrede, ligesom hos de aandsbeslægtede
Franskmænd ved Midten af det attende Aarhundrede,
en idealistisk Interesse for Humanitet og Fremskridt.
Idealistisk Sindelag er meget vel foreneligt med teoretisk Materialisme:
Materialisten kan anerkende de aandelige Fænomeners
og Bestræbelsers høje Værd, selvom han betragter
dem som blotte Atomsvingninger. --- Den mest fremragende
af disse Forfattere var Fysiologen Jacob Moleschott (1822
---1893). Født i Holland, virkede han i sin Ungdom som
Docent i Heidelberg, og da han blev afsat her paa Grund af
sine Anskuelser, drog han til Zürich og senere til Italien,
hvor han udøvede en betydningsfuld fysiologisk Lærervirksomhed.
I sin Bog \textit{Der Kreislauf des Lebens} (1852) betragter
han Kemien som den højeste Videnskab, fordi den viser,
hvorledes Stoffet, og med Stoffet Livet, med Livet igen Tanken,
gennemgaar sit store Kredsløb. I sin Selvbiografi (\textit{An
meine Freunde. Lebenserinnerungen.} 1895) fremstiller han
sine Tankers Historie og erklærer, at han egentligt kun har
kæmpet mod Dualismen, og at hans Lære paa Grund af den
indre Sammenhæng mellem Aand, Kraft og Stof, ligesaa vel
vilde kunne kaldes Idealisme som Materialisme. Han slutter sig
altsaa egentligt til Identitetshypotesen. Undertiden udtaler sig
ogsaa Lægen \textsc{Louis} \textsc{Büchner} (1824---1899) i denne Retning,
% --- p. 87 ---
\setcounter{footnote}{0}%
\opage{87}skønt han mener at hylde den rene Materialisme. Hans Bog
\textit{Kraft und Stoff} (1855) er meget uklar, hvad de principielle
Spørgsmaal angaar, men er skrevet i et populært Sprog og
vandt en meget stor Læsekreds.

I Modsætning til disse Forfattere stillede \textsc{Hermann} \textsc{Lotze}
(1817---1881) sig den Opgave at fastholde den romantiske
Periodes Idealisme i Sammenhæng med og paa Grundlag af
den moderne Naturvidenskabs Metoder og Resultater. Han
begyndte sin videnskabelige Virksomhed paa en Gang som
Naturforsker og Filosof, men gik senere helt over til Filosofien,
for hvilken han virkede en lang Række af Aar ved
Universitetet i Göttingen. Som Naturforsker var det hans
Bestræbelse at behandle Medicinen og Fysiologien som mekaniske
Naturvidenskaber uden den da endnu hyppige Beraabelse
paa en særlig „Livskraft“. Sine filosofiske Ideer har
han udviklet i \textit{Mikrokosmus} (1864---68) og i \textit{Drei Bücher
der Metaphysik} (1879).

Medens Romantikens Filosofi gik Deduktionens og Konstruktionens
Vej, er Lotze klar over, at de Tankemidler, der
staa til Raadighed ved Begrundelsen af vor Verdensanskuelse,
ere Analysen og Analogien. --- Den mekaniske Naturopfattelse
betragter alle Fænomener som bestemte ved Vexelvirkning af
Elementer, som ere udstrakte, selvom deres Udstrækning er
overordentligt lille. Men hvad der er udstrakt, maa bestaa af
Dele, kan ikke være absolut enkelt. Og Udstrækning er desuden
en Egenskab, der maa forklares efter samme Princip
som alle Sansekvaliteter, nemlig ved en Vexelvirkning af Elementer;
men de Elementer, hvis Vexelvirkning skal forklare
Udstrækningskvaliteten, kunne ikke selv være udstrakte, maa
altsaa være simplere end de Elementer, ved hvilke Naturvidenskaben
stanser; de kunne da kun være Kraftpunkter.
Og skulle de staa i Vexelvirkning med hverandre, kunne
de ikke være absolut selvstændige og forskellige; de maa
være Led i en Ursubstans, som omfatter dem alle og gør
det muligt, at det enes Tilstand kan bestemme det andets. ---
Længere end til de ubestemte Begreber om uudstrakte Elementer
og om en Ursubstans kunne vi ikke komme ad Ana%
% --- p. 88 ---
\setcounter{footnote}{0}%
\opage{88}lysens Vej. Kun ad Analogiens Vej kunne vi faa en positiv
Bestemmelse af disse Begreber. Lotze er enig med Leibniz
i, at al Metafysik beror paa Analogi. Idet han gaar ud fra, at
Tilværelsen i sig selv enten maa opfattes i Analogi med det
Materielle eller med det Aandelige, hævder han, at denne
sidste Analogi er den eneste mulige, fordi det Aandelige er
os umiddelbart givet i vor egen Natur, medens vi kun kende
det Materielle, vort eget Legeme og vor egen Hjerne indbefattede,
som Objekt. Vort Legeme er kun den Del af Omverdenen,
der er os nærmest. Tilværelsens Elementer og den
Ursubstans, der lægger sig for Dagen i denne Vexelvirkning,
maa da tænkes som besjælede i forskellige Grader. Kun paa
denne Maade faa vi en forstandig Verdensanskuelse.

Forholdet mellem Sjæl og Legeme frembyder for Lotze
ingen Vanskelighed, da de, ifølge hans metafysiske Idealisme,
bestaa af ensartede Elementer. (Se nærmere om hans Standpunkt
i denne Henseende Psykologi II 8 c).

\sectionhead{c. Erkendelsesproblemet.}

Henimod Slutningen af det nittende Aarhundrede banede
den Indsigt sig Vej, at baade den romantiske og den positivistiske
Reaktion mod Kant's kritiske Filosofi var gaaet for vidt.
Den kritiske Retning var i Aarhundredets Løb ikke forsvunden,
men den blev nu af større Betydning, og man fordybede
sig igen med Iver i Kant's Skrifter. I Tyskland har især
\textsc{Albert} \textsc{Lange} banet Vej for den saakaldte Nykriticisme eller
Nykantianisme. I Frankrig staar \textsc{Charles} \textsc{Renouvier}
(1815---1903) som den energiske Fornyer af den kritiske
Filosofi.

Kant's Opgave havde været at finde de Grundbegreber og
de Forudsætninger, paa hvilke en videnskabelig Forstaaelse
af Tilværelsen beror. I en vis Forstand kan man sige, at
disse Grundbegreber og Forudsætninger ere de Midler, ved
hvis Hjælp vi kunne naa det Maal, vor intellektuelle Interesse
lader os stræbe efter. Men Grundbegreberne (Rum, Tid, Størrelse,
Kausalitet) vare for Kant tillige Udtryk for den menne%
% --- p. 89 ---
\setcounter{footnote}{0}%
\opage{89}skelige Aands Former, for vor intellektuelle Natur, og det
var for ham intet Tilfælde, at vi i Videnskaben netop arbejde
med disse Grundbegreber. Derved adskiller \textit{Kriticismen} sig
fra en erkendelsesteoretisk Opfattelse, der er fremtraadt i de
senere Aar, og som udelukkende begrunder Videnskabens
Grundbegreber og Forudsætninger ved den Nytte, den teoretiske
og praktiske Betydning, de have. Denne Opfattelse, den
saakaldte \textit{Pragmatisme,} er hævdet af William James i nøje
Sammenhæng med hans ovenfor omtalte psykologiske Opfattelse,
ifølge hvilken vi kun, naar Nødvendigheden kræver
det, forlade den uvilkaarlige og umiddelbare Henstrømmen,
hvori vort Sjæleliv kundgør sig. De Tanker, vi udvikle, og de
Sætninger, vi opstille, fremkaldes stedse ved bestemte Interesser.
Stode vi aldeles ligegyldige, vilde vi ikke danne bestemte
Begreber eller drage bestemte Slutninger. (W. \textsc{James}:
\textit{Pragmatism.} 1907).

En pragmatisk Opfattelse blev ogsaa gjort gældende fra
naturvidenskabelig Side. Fremragende Naturforskere, der tænkte
over deres Videnskabs Principer, bleve opmærksomme paa,
at Videnskabens Begrebsbestemmelser og Forudsætninger maa
søge deres Retfærdiggørelse deri, at de muliggøre en intellektuel
Bearbejdelse og Fremstilling af Kendsgerningerne. Kun
herpaa beror deres Nødvendighed, som derfor ikke er ubetinget,
saalænge den Mulighed ikke er udelukket, at andre
Begreber og Forudsætninger end de, der nu anvendes, ligesaa
godt kunde føre til samme Maal. Den berømte engelske Fysiker
Maxwell udtalte sig allerede 1855 i denne Retning, og
senere har \textsc{Ernst} \textsc{Mach} paa Grundlag af Studier over Fysikens
Historie gjort det samme Synspunkt gældende. Begge
Forskere lægge stor Vægt paa de Analogier, ved hvis Hjælp
Videnskaben overfører Forestillinger fra det ene Erfaringsomraade
til det andet; saaledes beror den mekaniske Naturopfattelse
paa en stor, gennemført Analogi mellem Massers
Bevægelse i Rummet og Tingenes kvalitative Forandringer
(Temperatur, Elektricitet o. s. v.) (Ernst Mach: \textit{Die Mechanik
in ihrer} Entwickelung. 1883. --- \textit{Erkenntnis und} Irrtum. 1905).
Mach bruger istedetfor Ordet Pragmatisme hellere Udtrykket
% --- p. 90 ---
\setcounter{footnote}{0}%
\opage{90}økonomisk Erkendelsesteori, idet han lægger Vægt paa, at vi
indskrænke vore Definitioner og Hypoteser til det strengt
Nødvendige.

\textsc{Richard} \textsc{Avenarius} (1843---1896) har fremstillet den
økonomiske eller pragmatiske Erkendelseslære fra et psykologisk
Synspunkt, idet han har søgt at give, hvad man kunde
kalde Problemernes Naturhistorie. \textit{(Philosophie als} Denken
\textit{der Welt nach dem Prinzip des kleinsten Kraftmaasses.} 1876.
\textit{Kritik der reinen} Erfahrung. 1888---1890). Et Problem forudsætter
en Spænding mellem Individet og den omgivende Verden.
En saadan Spænding opstaar, naar de Paavirkninger, der
komme fra Omverdenen, kræve enten en større eller en
mindre Udløsning af Energi, end Individet formaar at yde.
Naar derimod Paavirkning (P) og Energi (E) svare til hinanden
(saa at P = E), er der intet Problem og ingen Uro;
Individet genkender sin Omverden, føler sig hjemme og har
Tillid til sine Forestillinger og Iagttagelser.

Kræves der et større Arbejde, end Individet kan yde
(altsaa P > E), føler Individet sig fremmed i sin Verden.
Genkendelse er ikke mulig, og det Givne staar som fuldt af
Modsigelser, Afvigelser og Undtagelser. Ved enhver Udvidelse
af Erfaringen og af den aandelige Horizont opstaar derfor
nye Problemer. Stigende Kultur medfører Forøgelse af Problemerne.
--- Er der omvendt mere Energi tilrede, end Omverdenen
kræver (saa at P < E), vil der opstaa en Trang,
der higer ud over det Givne, en romantisk Gæring eller en
praktisk Idealisme vil optræde. Uroen vil holde sig, indtil
enten Energien er udtømt eller Omverdenen er omdannet
efter Individets Trang.

Udligningen (Deproblematiseringen) er ikke altid en virkelig
Løsning. Der naas maaske kun et foreløbigt eller rent
individuelt Synspunkt, uden at Muligheden af en ny Spænding
(Problematisering) er udelukket. Definitiv og universel
bliver Deproblematiseringen kun, naar en fuldstændig Tillæmpelse
til det Givne er opnaat, og alle foreløbige og individuelle
Elementer i Opfattelsen ere udskilte. Dette opnaas,
naar Erkendelsen væsentligt bestaar i kvantitativ Beskrivelse,
% --- p. 91 ---
\setcounter{footnote}{0}%
\opage{91}der stedse paaviser det Følgende som Ækvivalent til det
Foregaaende. Da er den rene Erfarings Standpunkt naat. ---

I den danske Literatur peger K. \textsc{Kroman}'s Bog \textit{Vor Naturerkendelse}
(1883) i pragmatisk Retning, forsaavidt den i
Selvopholdelsestrangen finder det psykologiske Grundlag for
de Postulater, der opstille den videnskabelige Erkendelses
Principer (Identitetsprincipet, Aarsagsprincipet). ---

Medens Kriticismen lærer, at vi opstille de Grundbegreber
og Forudsætninger, \textit{vi efter vor Aands Natur maa opstille,}
lærer Pragmatismen, at vi opstille de Grundbegreber og Forudsætninger,
\textit{vi have Brug for.} Det er klart, at de Midler, vi
raade over til at løse vore Opgaver, maa være bestemte ved
vor Natur og afhænge af denne; der maa være en indre
Sammenhæng mellem vor Erkendelses Midler og vor Aands
Natur. Syntesen er paa en Gang Midlet til Erkendelse og
Form for vor Aands uvilkaarlige Virken. Pragmatismen holder
sig til de Anledninger, der føre til Udvikling af vore Betingelser
(vore „Former“); Kriticismen er tilbøjelig til at anse
„Formerne“ som givne en Gang for alle.

\medskip\noindent Tillæg. \textbf{Aarsagssætningens Historie i den nyere Tid.}

1. Sætningen opstilles som selvfølgelig, endog i den bestemte
Form, at der ikke kan være Mere i Virkningen end i
Aarsagen. Der skelnes ikke mellem Grund (ratio) og Aarsag
(causa). (\textsc{Descartes} og \textsc{Spinoza}).

2. \textsc{Hume}'s Kritik af Aarsagssætningen: denne kan hverken
bevises ved Slutning eller ved Iagttagelse, men grunder sig
kun paa Vane eller erhvervet Instinkt.

3. \textsc{Kant}: Aarsagsbegrebet er en Kategori, en af vor Tænknings
Grundformer, en speciel Art af Syntese. Aarsagssætningen
bevises ved, at vi kun formedelst den kunne skelne
mellem blot Forestilling og virkelig Erfaring; skal der være
Erfaring i streng Betydning (som forskellig fra blot Forestilling),
maa Aarsagsbegrebet gælde. Der er kun Analogi mellem
Grund og Aarsag.

4. Aarsagsbegrebet er et Middel til Erkendelse og Orien%
% --- p. 92 ---
\setcounter{footnote}{0}%
\opage{92}tering i Verden; Aarsagssætningen er blot en Arbejdshypotese.
\textit{(Pragmatismen).}

5. \textit{Den i Psykologi V D, 3 fremstillede Opfattelse} er enig
med Kant i, at Aarsagsbegrebet er en af de Former, i hvilke
Syntesen, Bevidsthedens formale Enhed, lægger sig for Dagen,
--- og enig med Hume i, at der ikke kan føres noget afsluttende
Bevis for Aarsagssætningens Gyldighed. Overfor Pragmatismen
fremhæves, at de Midler og Maader, vor Erkendelse
anvender, meget vel kunne være Udtryk for vor Aands Natur,
ja, maa være det, selvom vor Naturs Former ikke en
Gang for alle kunne findes og fastsættes. I Tankevirksomheden
ytrer sig en væsentlig Side ved vor Natur, og en saadan
Ytring er baade Maal og Middel. Den intellektuelle Interesse
har sin selvstændige Betydning i det menneskelige
Aandsliv.

\sectionhead{d. Vurderingsproblemet.}

α. Paa Etikens Omraade stod Brydningen henimod Slutningen
af det nittende Aarhundrede mellem Kants Fornuftmoral
og den af Bentham og Stuart Mill hævdede Nyttemoral;
hin holdt sig til den indre, frie Underkastelse under
en almengyldig Lov, denne til Handlingernes paaviselige
Virkninger. I en skarpsindig Bog \textit{(Methods of Ethics} 1877)
paaviste \textsc{Henry} \textsc{Sidgwick} (1838---1900), at der indenfor
Nyttemoralen maatte gøres en skarp Forskel mellem to Standpunkter:
det ene tager kun Hensyn til det enkelte Individs
egen Nytte, det andet forstaar ved Nytte den universelle
Nytte, hele Samfundets Nytte. Hvad der er nyttigt fra det
ene Standpunkt, er ikke altid nyttigt fra det andet Standpunkt,
og der kan her opstaa et vanskeligt, maaske uløseligt Problem.
I en mesterlig Analyse af den almindeligt anerkendte
praktiske Moral (the morality of common sense) viser Sidgwick,
at den har en universal Nyttemoral til sin Forudsætning.

I den Fremstilling af Etiken, denne Oversigts Forfatter
udgav 1887, paavistes, at Kants Etik og Benthams og Mill's
Nyttemoral trods al Modsætning hvilede paa en fælles For%
% --- p. 93 ---
\setcounter{footnote}{0}%
\opage{93}udsætning, nemlig at menneskelige Handlinger skulle vurderes
efter en almengyldig Maalestok. Selve denne Maalestok (den
Grundværdi, der bestemmer de enkelte Værdier) kan naturligvis
ikke igen vurderes. Det er den historiske og psykologiske
Undersøgelses Sag at paavise, hvorledes et Standpunkt,
der anerkender og anvender en saadan Maalestok, bliver til;
det er Etikens Sag at anvende den forudsatte Maalestok paa
de enkelte Livsforhold. I selve Forudsætningen om en almengyldig
Maalestok ligger, at den individuelle Etik maa være en
Del af den sociale. Det store Problem er, hvorledes et enkelt
Individ kan føres over fra en Grundværdi til en anden.

I den nyeste Tid er den sociale Karakter af den etiske
Vurdering især bleven stærkt indskærpet af franske Forskere.
Allerede hos Comte var Etiken en Del af Sociologien; først
i hans senere Periode udskilte han den. Den franske sociologiske
Skole med \textsc{Émile} \textsc{Durkheim} i Spidsen lægger den
sociale Solidaritet, grundet dels paa Blodets Baand, dels paa
Arbejdsfællesskab, til Grund for Moralen. Den Enkeltes Samvittighed
bestemmes, bevidst og ubevidst, ved Samfundets Indflydelse
gennem Overleveringer, Sæder, Indretninger og Tugt.
(\textsc{Durkheim}: \textit{La Division du Travail Social.} 1893). --- Medens
Durkheim overvejende undersøger den praktiske Morals historiske
og psykologiske Udvikling, gaa \textsc{Léon} \textsc{Bourgeois} og
hans Efterfølgere ud fra, at ethvert Individ formedelst Samfundets
Solidaritet staar i en Gæld til de tidligere Slægter,
som han maa afbetale ved Virken paa sin Samtids Opgaver.
Af Solidariteten udspringer Samfundets Ret og Individets Pligt.
Léon Bourgeois' Bog \textit{Solidarité} er udkommen i dansk Oversættelse
(1910). --- I dansk Literatur gaar N. \textsc{Bang}'s Tankegang
i \textit{Begrebet Moral} (1897) i lignende Retning.

β. Som Avenarius paaviste, opstaar der praktiske Problemer,
naar der hos Individet rører sig en større Fylde af
Energi, end Omgivelserne kunne lægge Beslag paa. En lignende
Tanke rører sig hos forskellige Forfattere i nyeste Tid,
der stærkt betone Spontaneiteten, den uvilkaarlige Stræben og
Trang til Expansion. I James' og Bergson's beskrivende Psykologi
kommer denne Side af Sjælelivet stærkt frem, og den
% --- p. 94 ---
\setcounter{footnote}{0}%
\opage{94}er ogsaa bleven anvendt paa Vurderingsproblemet, mest lidenskabeligt
og hensynsløst af \textsc{Friedrich} \textsc{Nietzsche} (1844---
% PRINTED AS IS (read at the image, structure pass s1/s2)
1900), der krævede en hel ny Livsvurdening, med Omvurdering
af alle hidtil anerkendte Værdier. I sit Ungdomsskrift \textit{Die
Geburt der Tragœdie} (1872) stiller han den tragisk-poetiske
Livsopfattelse, repræsenteret ved Dionysos og Apollon, i Modsætning
til Sokrates' intellektuelle Optimisme. Videnskaben
(Sokrates) staar under Kunsten (Apollon), der igen staar under
det overstrømmende Liv (Dionysos), der endog tager det
op med Lidelse og Død. Men den kraftige Livsudfoldelse,
uden hvilken Tilværelsen ikke faar positiv Værdi, fandt han
saa kun mulig, naar en herskende Stand repræsenterer Kraften
og Selvhævdelsen, medens en Kreds af Arbejdere og Spidsborgere
(die Heerde) besørger det underordnede Arbejde.
Aristokratiet er til for sin egen Skyld, da det udtrykker
Livets højeste Kraft og umiddelbare Værdi; i det naas Historiens
Maal, medens „Hjorden“ kun er Kopi, Middel eller Hindring
\textit{(Jenseits von Gut und} Böse. 1886. --- \textit{Die Genealogie
der Moral} 1887). I det digteriske Værk \textit{Also sprach Zarathustra,}
især i Udkastet til dets Afslutning, antydes dog Muligheden
af, at denne sociale Dualisme kunde overvindes, naar
Omvurderingen fuldendtes.

Ogsaa ud fra andre Standpunkter ere beslægtede Tanker
udtalte. --- \textsc{Jean} \textsc{Marie} Guyau's (1854---1888) Hovedtanke
var Livets Trang til Udfoldelse, Meddelelse og Hengivelse,
den spontane Virken af Sjælens Evner, der stedse paany fører
ud over opnaaede Resultater og hæmmende Skranker. I denne
Trang finder Guyau Muligheden for Sympati (Altruisme) og
Højsind, som den engelske Skole forklarede ved Forestillingsassociation
eller ved Kampen for Livet. \textit{(Esquisse d'une Morale
sans obligation ni sanction.} 1885). Og ud af denne Trang
vil der ogsaa kunne udvikle sig en ophøjet og modig Livsstemning,
der tør tage Kampen op med Tilværelsens Disharmonier,
selvom den ikke kan støtte sig til overnaturlige Garantier.
(\textit{L'irréligion de l'avenir.} 1887). --- Forfatteren af denne
historiske Oversigt har i sin Religionsfilosofi (1901) søgt at
vise, at der bag al Magi og Religion ligger en Trang til at fast%
% --- p. 95 ---
\setcounter{footnote}{0}%
\opage{95}holde den stadige Bestaaen af de Goder (Værdier), der for
Mennesket paa givet Trin ere de højeste. Denne Trang behøver
ikke at falde bort, selvom den ikke længere kan udpræges
i Myte, Dogme og Kultus, men maa søge Udtryk i
poetiske Symboler, og selvom Eftertanken fører til den Indsigt,
som allerede Værdiernes Historie under Menneskeslægtens
Udvikling fører til, at vore Værdier ikke kunne have vedvarende
Bestaaen i den Form, vi kende. Der kan være en
kontinuerlig Udvikling paa Værdiernes Omraade, som paa andre
Omraader, og det er den, der kan være Genstand for vor
Tro. --- \textsc{William James} viser i sit Skrift \textit{Varieties of religious
experience} (1902), støttet til et Studium af religiøse Fænomener,
saavel som til hvad han havde paavist i sin Psykologi,
hvor ringe en Del af vort aandelige Væsen, der kan forklares
tydeligt og bestemt. Gennem mange Grader gaar det Bevidste
tilbage i det Ubevidste, eller det Underbevidste, og vore første
Forudsætninger stamme ofte fra denne „subliminale“ Region,
fra det, der er under Bevidsthedens Tærskel. En uvilkaarlig
Overbevisning er det Dybeste i os (the deep thing in us).
James er tilbøjelig til at betragte Indflydelserne fra hin dunkle
Region som de Veje, ad hvilke en højere Tingenes Orden
virker i os. Enhver nærmere Bestemmelse af disse Indflydelsers
Oprindelse er en Fortolkning, der vil falde forskelligt ud
hos forskellige Individer. De fleste Mennesker er for tilbøjelige
til en dogmatisk Formulering af deres uvilkaarlige Trang.
Religion er ifølge James Menneskets Totalstemning overfor
Livet, en Stemning, der i sig indeholder Summen af Menneskets
Erfaringer om Livet. Derved opnaas en Enhed og Harmoni
i det aandelige Liv, som ellers vil mangle. Og den Livsfylde,
der kan strømme fra de ubevidste Kilder og stadigt
nærer vort bevidste Liv, kan være stærk nok til, at Tragedier
og Skibbrud aldrig for os betegne Erfaringens sidste Ord.

\closerule

\end{document}
